\documentclass[10pt,journal,onecolumn]{IEEEtran}

% ---------- Packages ----------
\usepackage[T1]{fontenc}
\usepackage{newtxtext}
\usepackage{amsmath}
\usepackage{amssymb}
\usepackage{newtxmath}
\usepackage{booktabs}
\usepackage{array}
\usepackage{tabularx}
\usepackage{multirow}
\usepackage{listings}
\usepackage{xcolor}
\usepackage[hidelinks]{hyperref}

% ---------- Paragraph breaking ----------
\emergencystretch=3em
\tolerance=2000

% ---------- Code style ----------
\definecolor{codegray}{gray}{0.94}
\definecolor{codebar}{gray}{0.65}
\lstdefinestyle{srs}{
  basicstyle=\ttfamily\footnotesize,
  backgroundcolor=\color{codegray},
  breaklines=true,
  showstringspaces=false,
  frame=single,
  rulecolor=\color{codebar},
  framerule=0.3pt,
  columns=fullflexible,
  keepspaces=true,
  aboveskip=4pt,
  belowskip=4pt
}
\lstset{style=srs}

% ---------- Commands ----------
\newcommand{\reqid}[1]{\textbf{#1}}

% ---------- Component divider ----------
\newcommand{\srsdivider}[1]{%
  \clearpage
  \vspace*{1.4em}
  \begin{center}
    {\LARGE\bfseries #1}\\[0.6em]
    \rule{\textwidth}{0.8pt}
  \end{center}
  \vspace{1.2em}
  \setcounter{section}{0}%
  \phantomsection
  \addcontentsline{toc}{section}{#1}%
}

% ---------- Title / Authors ----------
\title{NeuroShell Core Framework\\[0.4em] Combined Software Requirements Specification\\[0.15em] Components 01--04 (IRE, Planner, AEERE, AVAE)}
\author{\IEEEauthorblockN{NeuroShell Core Team}
\IEEEauthorblockA{Cyber Operations Research Group\\
MSc Cybersecurity Research Project R26-CS-018}}

\begin{document}

\maketitle

\begin{abstract}
This Software Requirements Specification (SRS) combines the four component specifications of the NeuroShell Core Framework into a single system-level reference: Component 01 (Intent Recognition Engine), Component 02 (Dynamic Planner and RAG Engine), Component 03 (Adaptive Execution and Error Recovery Engine), and Component 04 (AI Vulnerability Analysis Engine). The document records the complete functional and non-functional requirements, external interfaces, data requirements, and verification and validation characteristics of each component, together forming the autonomous NLU-driven offensive security microservices pipeline.
\end{abstract}

\begin{IEEEkeywords}
Software Requirements Specification, NeuroShell Core, Intent Recognition, Natural Language Understanding, RAG Planning, Error Recovery, Vulnerability Analysis, FastAPI, Cybersecurity Automation, Microservices.
\end{IEEEkeywords}

% =================== Document Control ===================
\section*{Document Control}
\begin{center}
\begin{tabularx}{0.9\columnwidth}{>{\bfseries}l >{\raggedright\arraybackslash}X}
Document ID & NS-SRS-001\\
Version & 1.0\\
Date & 2026-09-05\\
Project & NeuroShell Core Framework --- R26-CS-018\\
Reference & MSc Cybersecurity Research Project (R26-CS-018)\\
Classification & Internal / Academic Research\\
Status & Current\\
Components & C1 IRE / C2 Planner / C3 AEERE / C4 AVAE\\
\end{tabularx}
\end{center}

\section*{Revision History}
\begin{center}
\begin{tabularx}{0.9\columnwidth}{l l l X}
\toprule
\textbf{Version} & \textbf{Date} & \textbf{Author} & \textbf{Changes}\\
\midrule
1.0 & 2026-09-05 & NeuroShell Core Team & Initial combined system SRS (C1--C4)\\
\bottomrule
\end{tabularx}
\end{center}

\tableofcontents

\clearpage
\srsdivider{Component 01 --- Intent Recognition Engine (IRE)}
\section{Introduction}

\subsection{Purpose}
This Software Requirements Specification (SRS) defines the complete functional and non-functional requirements for the \textbf{NeuroShell Intent Recognition Engine (IRE)}, designated Component 01 of the NeuroShell Framework. The IRE is a domain-specific Natural Language Understanding (NLU) gateway that converts natural-language offensive-security operator commands into validated, schema-constrained JSON intent contracts.

This document is intended for:
\begin{itemize}
\item \textbf{Research evaluators} assessing the IRE's design and empirical performance;
\item \textbf{Developers} maintaining or extending the pipeline;
\item \textbf{Downstream consumers} (e.g., Component 02 Dynamic Planner) integrating with the IRE API;
\item \textbf{Examiners} reviewing the MSc Cybersecurity research project (ref.~R26-CS-018).
\end{itemize}

\subsection{Scope}

\subsubsection{Product Scope}
The IRE is a \textbf{self-contained FastAPI microservice} that:
\begin{enumerate}
\item Accepts natural-language offensive-security commands via a REST API;
\item Applies multi-layer validation (adversarial detection, schema, regex, network architecture, scope);
\item Produces a validated \texttt{IntentSchema} or \texttt{PlannerContract} JSON output;
\item Enforces role-based access control, session context, and hallucination detection;
\item Maintains a durable, privacy-preserving audit log.
\end{enumerate}

\subsubsection{In Scope}
\begin{itemize}
\item All source code under \texttt{src/} (pipeline, middleware, validation, inference, schemas, API, audit, utilities);
\item Configuration under \texttt{config/} (settings, features, RBAC policy, enforcement policy, LoRA config);
\item Grounding data under \texttt{data/} (CVE grounding, alias map);
\item Evaluation harness under \texttt{eval/} (diagnose, metrics, experiment scripts);
\item Test suite under \texttt{tests/} (unit and integration);
\item Deployment artifacts (Dockerfile, docker-compose.yml, scripts).
\end{itemize}

\subsubsection{Out of Scope}
\begin{itemize}
\item Component 02 (Dynamic Planner), Component 03 (AEERE Executor), and Component 04 (AVAE Vulnerability Analyzer) --- these are downstream consumers referenced only for interface-design justification;
\item Frontend UI (separate project);
\item Network infrastructure and physical security;
\item LoRA fine-tuning execution (Phase 2, planned but not yet implemented).
\end{itemize}

\subsection{Definitions, Acronyms, Abbreviations}
\begin{table}[h]
\caption{Definitions, Acronyms, and Abbreviations.}
\label{tab:definitions}
\begin{tabularx}{\columnwidth}{>{\ttfamily}l X}
\toprule
\textbf{Term} & \textbf{Definition}\\
\midrule
IRE & Intent Recognition Engine\\
NLU & Natural Language Understanding\\
LLM & Large Language Model\\
CVE & Common Vulnerabilities and Exposures\\
RBAC & Role-Based Access Control\\
CIDR & Classless Inter-Domain Routing\\
RFC-1918 & Address Allocation for Private Internets\\
SRS & Software Requirements Specification\\
Tier-1 & Single-sample LLM inference path\\
Tier-2 & Conditional N-sample self-consistency path\\
DAcc & Decision Accuracy\\
CFAcc & Contract-Field Accuracy\\
FP & False Positive (incorrectly blocked well-formed input)\\
S5/S6/S6B/S7 & Pipeline validation stages (Schema/Regex/Network/Scope)\\
M1--M11 & Pipeline middleware components\\
PlannerContract & Machine-readable JSON hand-off object for Component 02\\
Grounding Hash & SHA-256 digest of CVE and alias datasets for cache invalidation\\
\bottomrule
\end{tabularx}
\end{table}

\subsection{References}
\begin{thebibliography}{10}
\bibitem{R1} NeuroShell IRE README. \emph{README.md}, NeuroShell IRE Repository.
\bibitem{R2} IRE Research Extraction. \emph{IRE-Doc/research-extraction.md}.
\bibitem{R3} Experiment Results (900-record). \emph{IRE-Doc/NeuroShell\_IRE\_Test\_New\_Results.md}.
\bibitem{R4} Data Collection Plan. \emph{IRE-Doc/Data\_Collection\_Plan.md}.
\bibitem{R5} Results Section Draft. \emph{IRE-Doc/Results\_Section\_Draft.md}.
\bibitem{R6} Gaps and Future Works. \emph{IRE-Doc/Gaps \& Future Works.md}.
\bibitem{R7} IEEE Standard 830-1998, \emph{Recommended Practice for Software Requirements Specifications}, IEEE Computer Society.
\bibitem{R8} OWASP Top 10 (2021), owasp.org.
\bibitem{R9} RFC 1918 --- \emph{Address Allocation for Private Internets}, IETF.
\bibitem{R10} \emph{CVE Format Standard}, mitre.org.
\end{thebibliography}

\subsection{Overview}
The remainder of this document is organized as follows:
\begin{itemize}
\item \textbf{Section 2} describes the overall product perspective, user classes, and operating environment;
\item \textbf{Section 3} specifies all system features and functional requirements in detail;
\item \textbf{Section 4} defines external interface requirements (API, hardware, software, communication);
\item \textbf{Section 5} specifies non-functional requirements (performance, security, reliability);
\item \textbf{Section 6} covers data requirements (models, retention, grounding data);
\item \textbf{Section 7} describes verification, validation, and empirical results;
\item \textbf{Section 8} covers deployment procedures;
\item \textbf{Section 9} provides appendices with enumerations, policy tables, and schema definitions.
\end{itemize}
% =================== 2. OVERALL DESCRIPTION ===================
\section{Overall Description}

\subsection{Product Perspective}
The NeuroShell IRE is \textbf{Component 01} of the NeuroShell Framework, an autonomous penetration testing system. The IRE operates as an \textbf{independently deployable microservice} that bridges the gap between human operator intent (natural language) and machine-executable actions (validated JSON contracts).

\begin{lstlisting}
                +------------------------------------+
                |     Operator / Frontend UI         |
                |   (natural-language command)       |
                +----------------+------------------+
                                 | POST /parse
                                 v
                +------------------------------------+
                |     NeuroShell IRE (C01)           |
                |  NLU Gateway + Validation Stack    |
                |  Port 8001 (FastAPI)               |
                +--------+----------------+----------+
                         |                |
                  POST /parse/plan       POST /execute
                         |                |
                +--------v-------+  +----v--------------+
                | Component 02   |  | Component 03      |
                |   Dynamic      |  |  AEERE Executor   |
                |   Planner      |  |  (Port 8003)      |
                | (Port 8002)    |  +-------------------+
                +--------+-------+
                         |
                +--------v-------+
                | Component 04   |
                |    AVAE        |
                |  Vuln Analy.   |
                | (Port 8004)    |
                +----------------+
\end{lstlisting}

The IRE does \textbf{not} depend on downstream components for core operation. Component 02 (Planner), Component 03 (Executor), and Component 04 (Vulnerability Analyzer) are optional integrations controlled by feature flags.

\subsection{Product Functions}
The IRE provides the following core functions:

\begin{table*}[t]
\caption{IRE core product functions.}
\label{tab:functions}
\begin{tabularx}{\textwidth}{l X}
\toprule
\textbf{Function} & \textbf{Description}\\
\midrule
\textbf{Intent Classification} & Classify natural-language commands into one of 9 intent types.\\
\textbf{Target Extraction} & Extract and validate target specification (IP, subnet, domain, URL, hostname).\\
\textbf{Entity Extraction} & Extract ports, modifiers, CVE IDs, tool hints, and scheduling.\\
\textbf{Hallucination Detection} & Detect 6 classes of LLM hallucination via structured validation.\\
\textbf{Adversarial Detection} & Detect prompt injection, jailbreak, and encoded payload attacks.\\
\textbf{RBAC Enforcement} & Enforce role-based access control (4 roles, 9 intents).\\
\textbf{Session Management} & Multi-turn context injection with configurable memory window.\\
\textbf{Semantic Caching} & Exact-hash plus cosine-similarity response cache with grounding invalidation.\\
\textbf{Self-Consistency Voting} & Tier-2 N-sample consensus with semantic entropy uncertainty quantification.\\
\textbf{Contract Generation} & Produce \texttt{PlannerContract} JSON with tool synthesis and command generation.\\
\textbf{Audit Logging} & Privacy-preserving (SHA-256 hashed) append-only audit trail.\\
\textbf{Escalation Queue} & Human-in-the-loop review for high-severity findings with HTTP 202 pattern.\\
\bottomrule
\end{tabularx}
\end{table*}

\subsection{User Classes and Characteristics}
\begin{table*}[t]
\caption{User classes, privileges, and permitted intents.}
\label{tab:roles}
\begin{tabularx}{\textwidth}{l l X l}
\toprule
\textbf{Role} & \textbf{Privilege Level} & \textbf{Permitted Intents} & \textbf{Description}\\
\midrule
\textbf{viewer} & Read-only & None (pre-inference fast-fail) & Observation-only; cannot execute any commands.\\
\textbf{analyst} & Standard & NETWORK\_SCAN, VULNERABILITY\_AUDIT, SERVICE\_ENUMERATION, PASSIVE\_RECON, AMBIGUOUS, REJECTED & Passive and observational operations only.\\
\textbf{operator} & Elevated & All 9 intents & Active operations including exploitation and password attacks.\\
\textbf{admin} & Full & All 9 intents & System administration including model switching and policy management.\\
\bottomrule
\end{tabularx}
\end{table*}

\textbf{Default role:} \texttt{analyst} (when no role is specified in the request).

\textbf{Unknown roles:} Treated as \texttt{viewer} (most restrictive --- implicitly denied everywhere).

\subsection{Operating Environment}

\subsubsection{Hardware}
\begin{table}[h]
\caption{Hardware requirements.}
\label{tab:hardware}
\begin{tabularx}{\columnwidth}{l X X}
\toprule
\textbf{Component} & \textbf{Minimum} & \textbf{Recommended}\\
\midrule
CPU & 4 cores (2.0 GHz+) & 8 cores (3.0 GHz+)\\
RAM & 8 GB & 16 GB\\
GPU & None (CPU-only supported) & NVIDIA CUDA-capable GPU\\
Storage & 10 GB free & 50 GB SSD\\
Network & Loopback (localhost) & LAN (100 Mbps+)\\
\bottomrule
\end{tabularx}
\end{table}

\subsubsection{Software}
\begin{table}[h]
\caption{Software environment.}
\label{tab:software}
\begin{tabularx}{\columnwidth}{l l X}
\toprule
\textbf{Component} & \textbf{Version} & \textbf{Purpose}\\
\midrule
Python & 3.11+ & Runtime\\
Ollama & 0.32.x & LLM inference daemon\\
Docker & 24.x (optional) & Containerized deployment\\
Docker Compose & 2.x (optional) & Multi-service orchestration\\
\bottomrule
\end{tabularx}
\end{table}

\subsubsection{LLM Models}
\begin{table}[h]
\caption{LLM models.}
\label{tab:models}
\begin{tabularx}{\columnwidth}{l l X}
\toprule
\textbf{Model} & \textbf{Size} & \textbf{Purpose}\\
\midrule
\texttt{gemma4:latest} & 9.6 GB & Primary inference (default)\\
\texttt{qwen2.5-coder:7b} & 4.7 GB & Alternative inference backbone\\
\bottomrule
\end{tabularx}
\end{table}

\subsection{Design and Implementation Constraints}
\begin{enumerate}
\item \textbf{Language:} Python 3.11+ (async-first with FastAPI/uvicorn).
\item \textbf{Schema Validation:} Pydantic v2 with \texttt{extra="forbid"} (strict schema enforcement).
\item \textbf{Inference Runtime:} Ollama (local daemon, no cloud API calls).
\item \textbf{Async Concurrency:} Single-worker uvicorn (pipeline is CPU-bound during LLM inference).
\item \textbf{Rate Limiting:} slowapi, 30 requests/minute per IP.
\item \textbf{Input Limit:} 2000 characters per command.
\item \textbf{Session Limit:} 1024 concurrent sessions, 8 turns per session, 30-minute TTL.
\item \textbf{Cache Limits:} 256 entries (inference LRU), 512 entries (semantic cache).
\item \textbf{Privacy:} Raw input never stored in plaintext; SHA-256 hashing only.
\end{enumerate}

\subsection{Assumptions and Dependencies}

\subsubsection{Assumptions}
\begin{enumerate}
\item Ollama daemon is running and accessible at the configured URL.
\item Required LLM models are pre-pulled and available.
\item The operator has valid credentials for the assigned RBAC role.
\item Network connectivity exists between IRE and the Ollama daemon.
\item The engagement scope CIDR is correctly configured for the target environment.
\end{enumerate}

\subsubsection{Dependencies}
\begin{table}[h]
\caption{External dependencies.}
\label{tab:deps}
\begin{tabularx}{\columnwidth}{l l X}
\toprule
\textbf{Dependency} & \textbf{Type} & \textbf{Purpose}\\
\midrule
Ollama & Runtime & LLM inference engine\\
\texttt{data/cve\_grounding.json} & Data & CVE reference dataset (20 entries)\\
\texttt{data/alias\_map.json} & Data & Exploit-name to CVE alias mappings (17 entries)\\
\texttt{all-MiniLM-L6-v2} & ML Model & Sentence embeddings for semantic cache\\
\texttt{en\_core\_web\_sm} & ML Model & spaCy English NLP model\\
\bottomrule
\end{tabularx}
\end{table}
% =================== 3. SYSTEM FEATURES ===================
\section{System Features and Functional Requirements}

\subsection{Pipeline Architecture}

\subsubsection{Overview}
The IRE implements a \textbf{7-stage pipeline with 11 middleware components}, organized as follows:
\begin{lstlisting}
PRE-INFERENCE MIDDLEWARE (M1-M4)
    v
7-STAGE PIPELINE (S1-S7)
    v
POST-INFERENCE MIDDLEWARE (M5-M11)
    v
RESPONSE ASSEMBLY + AUDIT
\end{lstlisting}

\reqid{Requirement FR-PIPE-001:} The pipeline shall process each request through the full stage sequence (S1$\rightarrow$S7) unless an earlier stage raises a blocking exception.

\reqid{Requirement FR-PIPE-002:} Each middleware and stage shall be wrapped in individual try/except blocks so that a failure in non-critical middleware degrades gracefully rather than aborting the request.

\reqid{Requirement FR-PIPE-003:} The pipeline shall accumulate \texttt{ValidationFindings} across all stages (S5/S6/S6B/S7) and include them in the response payload regardless of pass/fail status.

\reqid{Requirement FR-PIPE-004:} Every non-passed finding shall be written to the SQLite audit log via \texttt{\_log\_audit\_findings}.

\subsubsection{Stage Sequence}
\begin{table*}[t]
\caption{Pipeline stage sequence (S1--S7).}
\label{tab:stages}
\begin{tabularx}{\textwidth}{l l X X X}
\toprule
\textbf{Stage} & \textbf{Module} & \textbf{Function} & \textbf{Input} & \textbf{Output}\\
\midrule
S1 & \texttt{InputNormalizer} & NFKC Unicode normalization, control-char strip, jargon standardization & Raw text & Normalized text\\
S2 & \texttt{AliasResolver} & Alias to CVE canonical mapping with provenance annotation & Normalized text & Enriched text\\
S3 & \texttt{OllamaInferenceEngine} & Tier-1/Tier-2 LLM inference with self-consistency & Enriched text & \texttt{IntentSchema}\\
S4 & \texttt{JSONParser} & Strip thought tags, fences; extract JSON object & LLM output text & Parsed dict\\
S5 & \texttt{SchemaValidator} & Pydantic v2 \texttt{IntentSchema} validation & Dict & \texttt{IntentSchema}\\
S6 & \texttt{RegexValidator} & IP/CIDR/CVE/port/metachar regex checks & \texttt{IntentSchema} & Validated \texttt{IntentSchema}\\
S6B & \texttt{NetworkArchitectureValidator} & Type consistency, \texttt{ipaddress} parse, scope containment, CIDR sanity & \texttt{IntentSchema} & Validated \texttt{IntentSchema}\\
S7 & \texttt{ScopeGuard} & 7 injection regexes on raw input plus RFC-1918 private-range check & Raw input + \texttt{IntentSchema} & (\texttt{IntentSchema}, warnings)\\
\bottomrule
\end{tabularx}
\end{table*}

\subsubsection{Middleware Sequence}
\begin{table*}[t]
\caption{Middleware sequence (M1--M11).}
\label{tab:middleware}
\begin{tabularx}{\textwidth}{l l l X}
\toprule
\textbf{Middleware} & \textbf{Module} & \textbf{Phase} & \textbf{Function}\\
\midrule
M1 & \texttt{AdversarialDetector} & Pre-inference & Weighted threat scoring $\rightarrow$ block/allow\\
M2 & \texttt{RBACGuard.pre\_inference\_check} & Pre-inference & Restricted-role fast-fail\\
M3 & \texttt{SessionContextStore.inject\_context} & Pre-inference & Prepend session turn history\\
M4 & \texttt{SemanticCache.lookup} & Pre-inference & Exact-hash then cosine cache hit\\
M5 & \texttt{RBACGuard.post\_inference\_check} & Post-inference & Role $\times$ intent authorization\\
M7 & \texttt{SemanticCache.store} & Post-inference & Store successful response\\
M8 & \texttt{IREResponseV2.from\_v1} & Post-inference & Assemble V2 response with metadata\\
M9 & \texttt{SubIntentClassifier.enrich} & Post-inference & Hierarchical sub-intent classification\\
M10 & \texttt{SessionContextStore.update} & Post-inference & Record completed turn\\
M11 & Session info enrich & Post-inference & Attach turn count to response\\
\bottomrule
\end{tabularx}
\end{table*}

\subsection{Pre-Inference Middleware (M1--M4)}

\subsubsection{M1 --- AdversarialDetector}
\reqid{Requirement FR-ADV-001:} The adversarial detector shall scan raw input against 15 lexical injection patterns (weights 0.8--1.0) and 7 structural anomaly patterns (weights 0.5--0.9) plus 3 heuristic checks.

\reqid{Requirement FR-ADV-002:} Any single signal with weight $\geq$ 0.95 shall trigger immediate blocking (auto-block), regardless of cumulative score.

\reqid{Requirement FR-ADV-003:} The cumulative threat score shall be capped at 1.0, and blocking shall occur when the score meets or exceeds the configured threshold (default 0.70).

\reqid{Requirement FR-ADV-004:} Scanning shall short-circuit as soon as the threshold is reached (efficiency optimization).

\reqid{Requirement FR-ADV-005:} The adversarial detector may be disabled via the \texttt{adversarial\_detection} feature flag.

Detection categories:
\begin{table}[h]
\caption{Adversarial detection layers.}
\label{tab:advlayers}
\begin{tabularx}{\columnwidth}{l >{\raggedright\arraybackslash}p{0.18\columnwidth} X >{\raggedright\arraybackslash}p{0.3\columnwidth}}
\toprule
\textbf{Layer} & \textbf{Patterns} & \textbf{Weight Range} & \textbf{Examples}\\
\midrule
A: Lexical & 15 regex patterns & 0.8--1.0 & \texttt{ignore previous instructions}, \texttt{jailbreak}, \texttt{$<$|im\_start|$>$}\\
B: Structural & 7 regex patterns & 0.5--0.9 & Hex encoding, base64 payloads, \texttt{$<$script$>$}, \texttt{eval()}\\
B: Heuristic & 3 checks & 0.3--0.4 & Special-char density $>$30\%, token $>$100 chars, repeated n-grams\\
\bottomrule
\end{tabularx}
\end{table}

Decision rule:
\begin{equation}
\mathrm{BLOCK} \iff \bigl(\exists\, i \in M(x):\ w_{i} \geq 0.95\bigr) \lor \bigl(S(x) \geq \tau\bigr),
\end{equation}
\begin{equation}
S(x) = \min\Bigl(1,\ \sum_{i \in M(x)} w_{i}\Bigr), \qquad \tau = 0.70.
\end{equation}

\subsubsection{M2 --- RBACGuard (Pre-Inference)}
\reqid{Requirement FR-RBAC-001:} The pre-inference check shall fast-fail for any role in the \texttt{RESTRICTED\_ROLES} set (\texttt{viewer}).

\reqid{Requirement FR-RBAC-002:} An \texttt{RBACError} shall be raised with \texttt{intent="ANY"} for restricted roles.

\reqid{Requirement FR-RBAC-003:} The pre-check may be disabled via the \texttt{rbac} feature flag.

\subsubsection{M3 --- SessionContextStore}
\reqid{Requirement FR-SES-001:} The session context store shall maintain per-session turn logs with configurable limits: max 8 turns, max 1024 sessions, 30-minute TTL, LRU eviction.

\reqid{Requirement FR-SES-002:} Context shall be injected in the following format:
\begin{lstlisting}
Session context:
[Turn 1] NETWORK_SCAN on SUBNET "192.168.1.0/24" ports=[22,80]
[Turn 2] VULNERABILITY_AUDIT on IP "10.0.0.5" cves=[CVE-2021-44228]
Current command: <user command>
\end{lstlisting}

\reqid{Requirement FR-SES-003:} Context injection shall occur BEFORE input normalization and BEFORE cache lookup, so cache keys include session state.

\reqid{Requirement FR-SES-004:} Only successful responses shall be recorded as turns (error responses are not stored).

\subsubsection{M4 --- SemanticCache}
\reqid{Requirement FR-CACHE-001:} The semantic cache shall support two-tier lookup: exact SHA-256 hash match first, then cosine similarity over all entries.

\reqid{Requirement FR-CACHE-002:} The cosine similarity threshold shall be configurable (default 0.82).

\reqid{Requirement FR-CACHE-003:} Each cache entry shall carry a \texttt{grounding\_hash} (SHA-256 of \texttt{cve\_grounding.json} + \texttt{alias\_map.json}). Entries with stale grounding hashes shall be purged on contact.

\reqid{Requirement FR-CACHE-004:} The cache shall use LRU eviction at capacity (default 512 entries).

\reqid{Requirement FR-CACHE-005:} Only successful responses shall be cached.

\reqid{Requirement FR-CACHE-006:} The embedding model (\texttt{all-MiniLM-L6-v2}) shall be lazy-loaded on first use.

\subsection{Core Pipeline Stages (S1--S7)}

\subsubsection{S1 --- InputNormalizer}
\reqid{Requirement FR-NORM-001:} Input normalization shall apply: NFKC Unicode normalization, control-char stripping (ord $<$ 32), think-token removal (\texttt{ thinking}, \texttt{$<$|think|$>$}), whitespace collapse.

\reqid{Requirement FR-NORM-002:} A 20-term jargon standardization table shall map offensive-security slang to canonical terms (e.g., \texttt{syn scan} $\rightarrow$ \texttt{SYN stealth scan}, \texttt{pwn} $\rightarrow$ \texttt{exploit}, \texttt{privesc} $\rightarrow$ \texttt{privilege escalation}).

\reqid{Requirement FR-NORM-003:} Empty or whitespace-only input shall raise an \texttt{InputValidationError} with an HTTP 200 error envelope.

\subsubsection{S2 --- AliasResolver}
\reqid{Requirement FR-ALIAS-001:} The alias resolver shall load \texttt{data/alias\_map.json} (17 active entries) and compile word-boundary patterns sorted longest-first.

\reqid{Requirement FR-ALIAS-002:} Aliases shall be rewritten to \texttt{CANONICAL (alias: ORIGINAL)} format so the LLM sees provenance.

\reqid{Requirement FR-ALIAS-003:} Each alias shall be rewritten only at its first occurrence per input.

\subsubsection{S3 --- OllamaInferenceEngine (Tier-1/Tier-2)}
\reqid{Requirement FR-INF-001:} The inference engine shall support two tiers of inference with conditional gating.

\reqid{Requirement FR-INF-002:} Tier-1 shall execute a single Ollama API call at the configured temperature (default 0.1).

\reqid{Requirement FR-INF-003:} The Tier-1 confidence proxy shall be computed as
\begin{equation}
c_{\mathrm{base}} = \begin{cases} o.\mathrm{confidence}, & \text{if parsed to } [0,1],\\ 0.5, & \text{otherwise.} \end{cases}
\end{equation}
\begin{equation}
c_{1} = \begin{cases} \min(c_{\mathrm{base}},\,0.4), & i \in \mathcal{H},\\ \min(c_{\mathrm{base}},\,0.5), & i \notin \mathcal{H},\ t_{v} \text{ missing},\\ c_{\mathrm{base}}, & \text{otherwise,} \end{cases}
\end{equation}
where $i = \mathrm{intent}$, $t_{v} = \mathrm{target.value}$, and $\mathcal{H} = \{\mathrm{AMBIGUOUS},\ \mathrm{REJECTED},\ \emptyset\}$.

\reqid{Requirement FR-INF-004:} If $c_{1} \geq \gamma$ ($\gamma = 0.7$), the Tier-1 output shall be returned immediately (entropy logged as 0, band ``low'', tier ``1'').

\reqid{Requirement FR-INF-005:} Tier-2 shall execute when $c_{1} < \gamma$, generating $N=5$ total samples (Tier-1 reused $+$ $N-1$ resamples) at a higher temperature (default 0.7).

\reqid{Requirement FR-INF-006:} Tier-2 outputs shall be clustered by semantic-equivalence signature
\begin{equation}
\begin{aligned}
\sigma(o) ={}& \bigl(\mathrm{intent},\ \mathrm{target\_type},\ \mathrm{target\_value.lower()},\\
&\quad \mathrm{tuple}(\mathrm{sorted}(cve\_ids))\bigr),
\end{aligned}
\end{equation}
which deliberately ignores ports and modifiers.

\reqid{Requirement FR-INF-007:} Semantic entropy shall be computed as
\begin{equation}
H(x) = -\sum_{j} p_{j}\, \ln p_{j}, \qquad p_{j} = |C_{j}|/N.
\end{equation}

\reqid{Requirement FR-INF-008:} Entropy bands shall be mapped via configurable thresholds:
\begin{itemize}
\item $H < 0.3$ $\rightarrow$ ``low'';
\item $0.3 \leq H \leq 0.8$ $\rightarrow$ ``medium'';
\item $H > 0.8$ $\rightarrow$ ``high''.
\end{itemize}

\reqid{Requirement FR-INF-009:} The resolution method for Tier-2 shall be plurality vote (largest cluster wins; ties resolve by insertion order).

\reqid{Requirement FR-INF-010:} The inference engine shall maintain a two-layer cache: in-memory LRU (256 entries) plus JSON disk persistence, keyed by \texttt{SHA-256(model\_name|enriched\_input)}.

\subsubsection{S4 --- JSONParser}
\reqid{Requirement FR-JSON-001:} The JSON parser shall strip \texttt{ thinking...response}, \texttt{$<$|think|$>$}, and markdown code fences.

\reqid{Requirement FR-JSON-002:} It shall extract the substring from the first \texttt{\{} to the last \texttt{\}} (inclusive).

\reqid{Requirement FR-JSON-003:} It shall validate that the result is a Python dict; failure raises \texttt{JSONParseError}.

\subsubsection{S5 --- SchemaValidator}
\reqid{Requirement FR-SCH-001:} The schema validator shall validate the parsed dict against the \texttt{IntentSchema} Pydantic v2 model via \texttt{model\_validate()}.

\reqid{Requirement FR-SCH-002:} On validation failure, a \texttt{ValidationFinding} shall be created with \texttt{severity="block"} and the appropriate \texttt{HallucinationClass} mapped from the failing field:
\begin{table*}[t]
\caption{Schema error to hallucination-class mapping.}
\label{tab:schmap}
\begin{tabularx}{\textwidth}{>{\ttfamily}l >{\ttfamily}X}
\toprule
\textbf{Error Field/Condition} & \textbf{Hallucination Class}\\
\midrule
\texttt{ports} in error path & FABRICATED\_PARAMETER\\
\texttt{cve\_ids} in error path & FABRICATED\_CVE\\
\texttt{target} or \texttt{intent} in error path & TARGET\_TYPE\_MISMATCH\\
\texttt{rejection\_reason}, \texttt{confidence}, AMBIGUOUS, REJECTED in message & UNGROUNDED\_CONFIDENCE\\
Default fallback & UNGROUNDED\_CONFIDENCE\\
\bottomrule
\end{tabularx}
\end{table*}

\reqid{Requirement FR-SCH-003:} Schema constraints include:
\begin{itemize}
\item \texttt{intent} $\in$ 9-class \texttt{IntentType} enum;
\item \texttt{target.type} $\in$ 6-class \texttt{TargetType} enum;
\item \texttt{confidence} $\in$ [0.0,~1.0];
\item \texttt{ports[$i$]} $\in$ [1,~65535];
\item \texttt{cve\_ids[$i$]} matches \texttt{\^{}CVE-\textbackslash d\{4\}-\textbackslash d\{4,7\}\$};
\item \texttt{intent=REJECTED} $\Rightarrow$ \texttt{rejection\_reason} required;
\item \texttt{intent=AMBIGUOUS} $\Rightarrow$ \texttt{confidence} $<$ 0.5;
\item No extra fields allowed (\texttt{extra="forbid"}).
\end{itemize}

\subsubsection{S6 --- RegexValidator}
\reqid{Requirement FR-REG-001:} Four sequential checks shall be performed:
\begin{table*}[t]
\caption{RegexValidator checks.}
\label{tab:regex}
\begin{tabularx}{\textwidth}{l X l l}
\toprule
\textbf{Check} & \textbf{Description} & \textbf{Failure Class} & \textbf{Severity}\\
\midrule
1. Target Format & Type-specific regex: IPv4, IPv6, CIDR, Domain, Hostname, URL & TARGET\_TYPE\_MISMATCH & block\\
2. CVE Syntax + Grounding & \texttt{\^{}CVE-\textbackslash d\{4\}-\textbackslash d\{4,7\}\$} syntax; then lookup in \texttt{cve\_grounding.json} & FABRICATED\_CVE / UNGROUNDED\_CONFIDENCE & block / from policy\\
3. Port Bounds & $\forall p$: $1 \leq p \leq 65535$ & FABRICATED\_PARAMETER & block\\
4. Shell Metacharacters & Scan for \texttt{[;\&|`\$!$>$]} in \texttt{target.value} and modifiers & FABRICATED\_PARAMETER & block\\
\bottomrule
\end{tabularx}
\end{table*}

\reqid{Requirement FR-REG-002:} CVE syntax failures shall produce severity from \texttt{EnforcementPolicy} (default: block).

\reqid{Requirement FR-REG-003:} Syntactically valid but ungrounded CVEs (not in \texttt{cve\_grounding.json}) shall produce \texttt{UNGROUNDED\_CONFIDENCE} with severity from policy (default: warn --- advisory only).

\subsubsection{S6B --- NetworkArchitectureValidator}
\reqid{Requirement FR-NET-001:} Four-step validation for IP/SUBNET targets:
\begin{table*}[t]
\caption{Network architecture validation steps.}
\label{tab:net}
\begin{tabularx}{\textwidth}{l X >{\raggedright\arraybackslash}p{0.2\textwidth}}
\toprule
\textbf{Step} & \textbf{Check} & \textbf{Failure Class}\\
\midrule
1. Type Consistency & IP must not contain ``/''; SUBNET must contain ``/''; UNKNOWN must be empty & TARGET\_TYPE\_MISMATCH\\
2. Parseability & \texttt{ipaddress.ip\_network()} or \texttt{ip\_address()} parse & format error\\
3. Scope Containment & IP: $\mathrm{addr} \in \Omega$; SUBNET: $\mathrm{net} \subseteq \Omega$ (\texttt{subnet\_of}) & OUT\_OF\_SCOPE\_TARGET\\
4. CIDR Sanity & Prefix $\geq$ \texttt{max\_cidr\_prefix} (30); reject /32 subnet; IPv6 $<$ /64 rejected & OUT\_OF\_SCOPE\_TARGET\\
\bottomrule
\end{tabularx}
\end{table*}

\reqid{Requirement FR-NET-002:} Domain, URL, Hostname, and UNKNOWN types shall auto-pass the scope check (scope validation not applicable).

\reqid{Requirement FR-NET-003:} Engagement mode \texttt{"disabled"} shall short-circuit all validation; \texttt{"warn"} shall downgrade all severities to advisory warnings.

\subsubsection{S7 --- ScopeGuard}
\reqid{Requirement FR-SCOPE-001:} The scope guard shall scan the \textbf{raw input} (not the enriched/normalized form) against 7 injection regex patterns. Any match shall produce an unconditional hard block.

\reqid{Requirement FR-SCOPE-002:} Private-range membership shall be checked against 7 network ranges:
\begin{itemize}
\item \texttt{10.0.0.0/8}, \texttt{172.16.0.0/12}, \texttt{192.168.0.0/16};
\item \texttt{127.0.0.0/8} (loopback), \texttt{169.254.0.0/16} (link-local);
\item \texttt{::1/128} (IPv6 loopback), \texttt{fc00::/7} (IPv6 ULA).
\end{itemize}

\reqid{Requirement FR-SCOPE-003:} Public (non-private) targets shall trigger \texttt{OUT\_OF\_SCOPE\_TARGET} with severity determined by \texttt{EnforcementPolicy}:
\begin{itemize}
\item \texttt{"strict"} mode $\rightarrow$ block;
\item \texttt{"research"} mode $\rightarrow$ warn (advisory).
\end{itemize}

\subsection{Post-Inference Middleware (M5--M11)}

\subsubsection{M5 --- RBACGuard (Post-Inference)}
\reqid{Requirement FR-RBAC-010:} The post-inference check shall verify that the detected intent is permitted for the caller's role using \texttt{INTENT\_ROLE\_POLICY}.

\reqid{Requirement FR-RBAC-011:} Denied intents shall raise \texttt{RBACError} listing the detected intent and the role's permitted intents.

\subsubsection{M7 --- SemanticCache (Store)}
\reqid{Requirement FR-CACHE-010:} Successful responses shall be stored with their grounding hash for future invalidation.

\subsubsection{M8 --- IREResponseV2 Assembly}
\reqid{Requirement FR-RESP-001:} The V2 response shall include all V1 fields plus: \texttt{sub\_intent}, \texttt{secondary\_intents}, \texttt{uncertainty\_band}, \texttt{raw\_entropy}, \texttt{tier\_triggered}, \texttt{resolution\_method}, \texttt{cache\_hit}, \texttt{rbac\_role}, \texttt{session\_turns}, \texttt{validation\_findings}, \texttt{schema\_version}.

\reqid{Requirement FR-RESP-002:} Error responses shall follow the 200-envelope pattern: HTTP 200 with \texttt{status:"error"}, carrying \texttt{error}, \texttt{stage}, \texttt{detail}, \texttt{field}, \texttt{validation\_findings}, and \texttt{latency\_ms}.

\subsubsection{M9 --- SubIntentClassifier}
\reqid{Requirement FR-SUB-001:} The sub-intent classifier shall be a first-match-wins rule table over 27 sub-intent types across 7 parent intents.

\reqid{Requirement FR-SUB-002:} Classification shall be gated: disabled for AMBIGUOUS/REJECTED intents and for confidence $<$ 0.6.

\subsubsection{M10/M11 --- Session Update}
\reqid{Requirement FR-SES-010:} After successful processing, the session store shall record the turn with intent, target, ports, CVE IDs, and confidence.
\subsection{Schema and Contract System}

\subsubsection{IntentSchema (IRE Internal)}
The \texttt{IntentSchema} is the IRE's internal validated representation:
\begin{table*}[t]
\caption{IntentSchema fields.}
\label{tab:intentschema}
\begin{tabularx}{\textwidth}{>{\ttfamily}l l l X}
\toprule
\textbf{Field} & \textbf{Type} & \textbf{Required} & \textbf{Constraints}\\
\midrule
\texttt{intent} & IntentType & Yes & 9-class enum\\
\texttt{target} & Target & Yes & \texttt{\{type: TargetType, value: str\}}\\
\texttt{ports} & List[int] & No & Each $\in$ [1, 65535]\\
\texttt{modifiers} & List[str] & No & Free text\\
\texttt{cve\_ids} & List[str] & No & Each matches \texttt{\^{}CVE-\textbackslash d\{4\}-\textbackslash d\{4,7\}\$}\\
\texttt{tool\_hint} & Optional[str] & No & Suggested tool\\
\texttt{schedule} & Optional[str] & No & Cron expression\\
\texttt{confidence} & float & Yes & $\in$ [0.0, 1.0]\\
\texttt{rejection\_reason} & Optional[str] & Yes if REJECTED & Required when \texttt{intent=REJECTED}\\
\bottomrule
\end{tabularx}
\end{table*}

\subsubsection{PlannerContract (Component 02 Hand-off)}
The \texttt{PlannerContract} extends the intent schema with downstream-consumption fields:
\begin{table}[h]
\caption{PlannerContract fields.}
\label{tab:plannercontract}
\begin{tabularx}{\columnwidth}{>{\ttfamily}l l X}
\toprule
\textbf{Field} & \textbf{Type} & \textbf{Purpose}\\
\midrule
\texttt{sub\_intent} & Optional[SubIntentType] & Hierarchical refinement for C02 command construction\\
\texttt{tool\_hint} & Optional[str] & Tool suggestion (nmap, gobuster, nikto, manual, none)\\
\texttt{schedule} & Optional[str] & Cron expression; None = execute immediately\\
\texttt{scope\_warnings} & List[str] & SCOPE\_WARNING and ARCH\_WARNING entries\\
\bottomrule
\end{tabularx}
\end{table}

Executability predicates:
\begin{itemize}
\item \texttt{is\_executable} $\rightarrow$ False if REJECTED/AMBIGUOUS, or if any \texttt{scope\_warning} contains ``SCOPE\_VIOLATION'' or ``ARCH\_WARNING'';
\item \texttt{is\_immediate} $\rightarrow$ True if \texttt{schedule} is None;
\item \texttt{has\_scope\_issues} $\rightarrow$ True if any \texttt{scope\_warnings} present.
\end{itemize}

\subsection{Hallucination Taxonomy and Enforcement}

\subsubsection{HallucinationClass Enum (6 Classes)}
\begin{table*}[t]
\caption{Hallucination classes.}
\label{tab:hallucination}
\begin{tabularx}{\textwidth}{>{\ttfamily}l X X}
\toprule
\textbf{Class} & \textbf{Description} & \textbf{Detection Stage}\\
\midrule
FABRICATED\_CVE & Malformed or non-existent CVE ID & S5 (schema), S6 (regex+grounding)\\
TARGET\_TYPE\_MISMATCH & Structural mismatch between target type and value & S5 (schema), S6 (format), S6B (consistency)\\
CONTRADICTORY\_ACTION\_TARGET & Incompatible action-target pairing & S5 (schema)\\
FABRICATED\_PARAMETER & Out-of-bounds or injected parameters & S5 (schema), S6 (port/metachar)\\
UNGROUNDED\_CONFIDENCE & Valid syntax but unindexed CVE; low-confidence output & S6 (grounding), S5 (confidence)\\
OUT\_OF\_SCOPE\_TARGET & Target outside authorized engagement scope & S6B (scope), S7 (RFC-1918)\\
\bottomrule
\end{tabularx}
\end{table*}

\subsubsection{Enforcement Policy}
\begin{table}[h]
\caption{Default enforcement by class.}
\label{tab:enforcement}
\begin{tabularx}{\columnwidth}{>{\ttfamily}l l X}
\toprule
\textbf{HallucinationClass} & \textbf{Default Action} & \textbf{Reasoning}\\
\midrule
FABRICATED\_PARAMETER & BLOCK & Execution risk from out-of-bounds parameters or injection\\
CONTRADICTORY\_ACTION\_TARGET & BLOCK & Safety invariant violation\\
TARGET\_TYPE\_MISMATCH & BLOCK & Structural mismatch prevents downstream execution\\
FABRICATED\_CVE & BLOCK & Complete LLM hallucination\\
UNGROUNDED\_CONFIDENCE & WARN & Novel but plausible CVEs; advisory only\\
OUT\_OF\_SCOPE\_TARGET & Dynamic & BLOCK in strict mode; WARN in research mode\\
\bottomrule
\end{tabularx}
\end{table}

\subsubsection{ValidationFinding Schema}
Each validation check produces a \texttt{ValidationFinding}:
\begin{table*}[t]
\caption{ValidationFinding fields.}
\label{tab:vf}
\begin{tabularx}{\textwidth}{>{\ttfamily}l l X}
\toprule
\textbf{Field} & \textbf{Type} & \textbf{Values}\\
\midrule
\texttt{validator} & str & Which component performed the check\\
\texttt{passed} & bool & True if check passed\\
\texttt{hallucination\_class} & Optional[HallucinationClass] & Assigned class on failure\\
\texttt{detail} & str & Human-readable description\\
\texttt{severity} & Literal["block","warn","escalate"] & Enforcement action\\
\bottomrule
\end{tabularx}
\end{table*}

\subsection{Contract Builder and Tool Synthesis}

\subsubsection{Tool Synthesis Table}
\begin{table*}[t]
\caption{Tool synthesis per intent.}
\label{tab:toolsynth}
\begin{tabularx}{\textwidth}{>{\ttfamily}l l X}
\toprule
\textbf{Intent} & \textbf{Tool} & \textbf{Command Template}\\
\midrule
NETWORK\_SCAN & nmap & \texttt{nmap [flags] $<$target$>$}\\
VULNERABILITY\_AUDIT & nikto (URL) / nmap & \texttt{nikto -h $<$target$>$ [-p port] [-ssl]}\\
SERVICE\_ENUMERATION & nmap & \texttt{nmap -sV [flags] $<$target$>$}\\
EXPLOITATION & manual & Requires human authorization\\
PASSWORD\_ATTACK & manual & Lockout risk warning\\
DIRECTORY\_BRUTEFORCE & gobuster & \texttt{gobuster dir -u $<$url$>$ -w $<$wordlist$>$ -q}\\
PASSIVE\_RECON & none & Information gathering only\\
AMBIGUOUS & none & Requires clarification\\
REJECTED & none & Command rejected\\
\bottomrule
\end{tabularx}
\end{table*}

\subsubsection{Priority Ladder}
\begin{table}[h]
\caption{Contract priority ladder.}
\label{tab:priority}
\begin{tabularx}{\columnwidth}{l X}
\toprule
\textbf{Priority} & \textbf{Condition}\\
\midrule
SCHEDULED & Cron expression is set\\
CRITICAL & Intent is EXPLOITATION\\
HIGH & Intent is VULNERABILITY\_AUDIT with CVE IDs\\
MEDIUM & NETWORK\_SCAN, SERVICE\_ENUMERATION, or VULNERABILITY\_AUDIT (no CVEs)\\
LOW & All other intents\\
\bottomrule
\end{tabularx}
\end{table}

\subsubsection{Nmap Flag Mapping}
\begin{table}[h]
\caption{Modifier to nmap-flag mapping.}
\label{tab:nmapflags}
\begin{tabularx}{\columnwidth}{>{\ttfamily}l >{\ttfamily}l}
\toprule
\textbf{Modifier} & \textbf{Flag}\\
\midrule
\texttt{stealth} & \texttt{-sS} (SYN scan)\\
\texttt{aggressive} & \texttt{-A}\\
\texttt{version\_detect} & \texttt{-sV}\\
\texttt{os\_detect} & \texttt{-O}\\
\texttt{udp} & \texttt{-sU}\\
\texttt{full\_port} & \texttt{-p-}\\
\texttt{verbose} & \texttt{-v}\\
\texttt{stealth} timing & \texttt{-T1}\\
Default timing & \texttt{-T3}\\
\bottomrule
\end{tabularx}
\end{table}

\subsection{Audit and Escalation System}

\subsubsection{Audit Logger}
\reqid{Requirement FR-AUDIT-001:} The audit logger shall maintain an append-only SQLite \texttt{audit\_log} table.

\reqid{Requirement FR-AUDIT-002:} Raw user input shall NEVER be stored in plaintext; only SHA-256 hashes shall be recorded.

\reqid{Requirement FR-AUDIT-003:} Each finding shall be recorded with: timestamp, session\_id, raw\_input\_hash, intent\_summary, hallucination\_class, severity, enforcement\_action, and status.

\subsubsection{Escalation Queue}
\reqid{Requirement FR-ESC-001:} Findings with \texttt{severity="escalate"} shall be recorded with \texttt{status="pending"} in the audit log.

\reqid{Requirement FR-ESC-002:} The \texttt{POST /parse/plan} endpoint shall return HTTP 202 with \texttt{escalation\_id} when an ESCALATE-severity finding is encountered, withholding the contract.

\reqid{Requirement FR-ESC-003:} Admins may approve or reject escalations via \texttt{POST /admin/escalations/\{id\}/resolve}, which releases or discards the cached contract.
% =================== 4. EXTERNAL INTERFACES ===================
\section{External Interface Requirements}

\subsection{User Interfaces}
The IRE does not have a direct user interface. It provides a REST API consumed by:
\begin{itemize}
\item Frontend UI (separate project, \texttt{apps/web/});
\item Downstream components (C02, C03, C04);
\item CLI tools and scripts.
\end{itemize}

\subsection{Hardware Interfaces}
\begin{table}[h]
\caption{Hardware interfaces.}
\label{tab:hwinterface}
\begin{tabularx}{\columnwidth}{l X}
\toprule
\textbf{Interface} & \textbf{Specification}\\
\midrule
Network & TCP/IP over Ethernet or loopback\\
Ollama Communication & HTTP/1.1 to Ollama daemon (default port 11434)\\
Storage & Filesystem for SQLite audit log, JSON cache files, grounding data\\
\bottomrule
\end{tabularx}
\end{table}

\subsection{Software Interfaces}
\begin{table*}[t]
\caption{Software interface map.}
\label{tab:swinterfaces}
\begin{tabularx}{\textwidth}{l l X}
\toprule
\textbf{Interface} & \textbf{Protocol} & \textbf{Target}\\
\midrule
Ollama API & HTTP POST \texttt{/api/generate} & Local Ollama daemon\\
Component 02 & HTTP POST \texttt{/plan} & Dynamic Planner (port 8002, optional)\\
Component 03 & HTTP POST & AEERE Executor (port 8003, optional)\\
Component 04 & HTTP POST & AVAE Vulnerability Analyzer (port 8004, optional)\\
SQLite & File I/O & \texttt{data/audit\_log.db}\\
Filesystem & Read & \texttt{data/cve\_grounding.json}, \texttt{data/alias\_map.json}\\
\bottomrule
\end{tabularx}
\end{table*}

\subsection{Communication Interfaces}

\subsubsection{REST API Base}
\begin{itemize}
\item \textbf{Protocol:} HTTP/1.1 (HTTPS recommended for production);
\item \textbf{Content-Type:} \texttt{application/json};
\item \textbf{Authentication:} \texttt{X-API-Key} header (Bearer-style);
\item \textbf{Rate Limiting:} 30 requests/minute per IP (slowapi);
\item \textbf{CORS:} Allow all origins (development mode).
\end{itemize}

\subsubsection{Response Headers}
\begin{table}[h]
\caption{IRE response headers.}
\label{tab:headers}
\begin{tabularx}{\columnwidth}{>{\ttfamily}l X}
\toprule
\textbf{Header} & \textbf{Description}\\
\midrule
\texttt{X-IRE-Schema-Version} & Schema version(s) supported: \texttt{"1,2"}\\
\texttt{X-IRE-Latency-Ms} & Pipeline processing latency in milliseconds\\
\texttt{X-IRE-Intent} & Detected intent type\\
\texttt{X-IRE-Cache-Hit} & Cache hit type: \texttt{"exact"}, \texttt{"semantic"}, or \texttt{"none"}\\
\texttt{X-Chain-Cache-Hit} & Chain cache hit (on \texttt{/execute} only)\\
\bottomrule
\end{tabularx}
\end{table}

\subsection{API Endpoint Specifications}

\subsubsection{GET /}
\textbf{Purpose:} Service information (no auth required).

\textbf{Response 200:}
\begin{lstlisting}
{
  "service": "NeuroShell IRE",
  "version": "2.0.0",
  "status": "running",
  "docs": "/docs",
  "endpoints": { ... }
}
\end{lstlisting}

\subsubsection{GET /health}
\textbf{Purpose:} System health check (no auth required).

\textbf{Response 200:}
\begin{lstlisting}
{
  "status": "healthy",
  "ollama": { "connected": true, "model": "gemma4:latest" },
  "pipeline": { "stages": 7, "features": { ... } },
  "sessions": { "active": 0, "max": 1024 },
  "semantic_cache": { "size": 0, "max_size": 512 },
  "inference_cache": { "size": 0 }
}
\end{lstlisting}

\textbf{Response 503:} Service unavailable.

\subsubsection{POST /parse}
\textbf{Purpose:} Core intent recognition endpoint.

\textbf{Request:}
\begin{lstlisting}
{
  "command": "scan 192.168.1.0/24 for open ports",
  "session_id": "optional-session-id",
  "role": "analyst"
}
\end{lstlisting}

\textbf{Constraints:}
\begin{itemize}
\item \texttt{command}: required, non-empty, $\leq$ 2000 characters;
\item \texttt{role}: optional, default \texttt{"analyst"}.
\end{itemize}

\textbf{Response 200 (success):}
\begin{lstlisting}
{
  "status": "success",
  "intent": "NETWORK_SCAN",
  "target": { "type": "SUBNET", "value": "192.168.1.0/24" },
  "ports": [],
  "modifiers": [],
  "cve_ids": [],
  "tool_hint": null,
  "schedule": null,
  "confidence": 0.85,
  "rejection_reason": null,
  "scope_warnings": [],
  "latency_ms": 27650,
  "sub_intent": "SYN_STEALTH",
  "uncertainty_band": "low",
  "raw_entropy": 0.0,
  "tier_triggered": 1,
  "resolution_method": "single_sample",
  "cache_hit": "none",
  "rbac_role": "analyst",
  "validation_findings": [],
  "schema_version": 2
}
\end{lstlisting}

\textbf{Response 200 (error envelope):}
\begin{lstlisting}
{
  "status": "error",
  "error": "SCOPE_VIOLATION",
  "stage": "scope_guard",
  "detail": "Target outside authorized scope",
  "latency_ms": 150,
  "validation_findings": [...]
}
\end{lstlisting}

\textbf{Other Responses:} 400 (bad request), 401 (unauthorized), 403 (RBAC denied), 413 (payload too large), 429 (rate limited), 503 (service unavailable).

\subsubsection{POST /parse/plan}
\textbf{Purpose:} Intent recognition plus \texttt{PlannerContract} generation for Component 02.

\textbf{Request:} Same as \texttt{/parse}.

\textbf{Response 200:} \texttt{PlannerContract} JSON (see Section 3.5.2).

\textbf{Response 202:} Escalation pending --- returns \texttt{escalation\_id} and \texttt{status: "pending\_review"}, withholds contract.

\textbf{Response 422:} Contract validation failure.

\subsubsection{POST /execute}
\textbf{Purpose:} Full synchronous C1$\rightarrow$C2$\rightarrow$C3$\rightarrow$C4 chain execution.

\textbf{Request:} Same as \texttt{/parse}.

\textbf{Response 200:} Combined multi-component result (IRE parse + planner output + executor output + vulnerability analysis).

\subsubsection{GET /metrics}
\textbf{Purpose:} Pipeline performance metrics (requires auth).

\textbf{Response 200:}
\begin{lstlisting}
{
  "intent_distribution": { ... },
  "error_counts": { ... },
  "latency_percentiles": { "p50": 16498, "p95": 70893, "p99": ... },
  "entropy_records": { ... },
  "cache_hits": { "exact": 0, "semantic": 0 }
}
\end{lstlisting}

\subsubsection{GET/DELETE /session/\{session\_id\}}
\textbf{Purpose:} Session management (requires auth).

\textbf{GET Response 200:} Session info with turns list.

\textbf{DELETE Response 200:}
\begin{lstlisting}
{ "cleared": true }
\end{lstlisting}

\subsubsection{Admin Endpoints}
\begin{table*}[t]
\caption{Administrative endpoints.}
\label{tab:admin}
\begin{tabularx}{\textwidth}{l >{\ttfamily}l X}
\toprule
\textbf{Method} & \textbf{Path} & \textbf{Purpose}\\
\midrule
GET & \texttt{/admin/rbac/\{role\}} & Role permissions summary\\
GET & \texttt{/admin/features} & Feature flags + tuning params\\
POST & \texttt{/admin/features/reload} & Reload features.yaml\\
GET & \texttt{/admin/models} & List available LLM models + active\\
POST & \texttt{/admin/models/switch} & Switch active model (persists to .env)\\
GET & \texttt{/admin/enforcement-policy} & Full enforcement policy\\
GET & \texttt{/admin/escalations} & Escalation queue\\
POST & \texttt{/admin/escalations/\{id\}/resolve} & Approve/reject escalation\\
GET & \texttt{/admin/audit/summary} & Audit log summary\\
GET & \texttt{/admin/chain-cache} & Chain cache stats\\
DELETE & \texttt{/admin/chain-cache} & Clear chain cache\\
GET & \texttt{/engagement/scope} & Current engagement scope config\\
\bottomrule
\end{tabularx}
\end{table*}
% =================== 5. NON-FUNCTIONAL REQUIREMENTS ===================
\section{Non-Functional Requirements}

\subsection{Performance Requirements}

\subsubsection{Latency (Measured on CPU-only Hardware)}
\begin{table*}[t]
\caption{Latency metrics.}
\label{tab:latency}
\begin{tabularx}{\textwidth}{l X X X}
\toprule
\textbf{Metric} & \textbf{Raw LLM} & \textbf{Full Pipeline} & \textbf{Notes}\\
\midrule
p50 & 11.5--13.5 s & 16.5 s & Pipeline adds $\sim$3--5 s validation overhead\\
p95 & 16.5--63.3 s & 70.9 s & Dominated by Tier-2 spikes on uncertain inputs\\
p99 & 17.4--93.7 s & $\sim$150 s & Rare Tier-2 chains\\
Mean & 13.4--25.3 s & 28.8 s & Weighted by Tier-1 majority\\
\bottomrule
\end{tabularx}
\end{table*}

\reqid{Requirement NFR-PERF-001:} Pipeline validation overhead (S5--S7) shall not exceed 5 seconds for single-sample requests.

\reqid{Requirement NFR-PERF-002:} Adversarial detection (M1) shall complete in $<$ 5 ms.

\reqid{Requirement NFR-PERF-003:} The semantic cache lookup shall complete in $<$ 50 ms for up to 512 entries.

\textbf{Note:} All latency figures reflect \textbf{CPU-only inference} on an ASUS VivoBook X513EA (Intel i5-1135G7, 4C/8T, 24 GB DDR4). GPU deployment would significantly reduce Tier-1 and Tier-2 latencies.

\subsubsection{Throughput}
\begin{table}[h]
\caption{Throughput limits.}
\label{tab:throughput}
\begin{tabularx}{\columnwidth}{l X}
\toprule
\textbf{Metric} & \textbf{Value}\\
\midrule
Rate limit & 30 requests/minute per IP\\
Concurrent sessions & 1024\\
Cache capacity & 256 (inference) + 512 (semantic)\\
\bottomrule
\end{tabularx}
\end{table}

\subsection{Safety Requirements}
\reqid{Requirement NFR-SAF-001:} The pipeline shall be fail-closed: unknown roles are denied by absence from permit lists; schema violations hard-fail.

\reqid{Requirement NFR-SAF-002:} The pipeline bias shall be toward the safe cell: every observed false positive is an over-block (false rejection), never a dangerous misclassification.

\reqid{Requirement NFR-SAF-003:} Exploitation and password attack intents shall require at least the \texttt{operator} role.

\reqid{Requirement NFR-SAF-004:} Shell metacharacters in target/modifiers shall be blocked regardless of RBAC role.

\reqid{Requirement NFR-SAF-005:} The pipeline shall produce ESCALATE findings (not ALLOW) when enforcement severity is \texttt{"escalate"}, deferring to human review.

\subsection{Security Requirements}

\subsubsection{Input Security}
\reqid{Requirement NFR-SEC-001:} The pipeline shall defend against 4 adversary capabilities:
\begin{enumerate}
\item \textbf{T1 --- Prompt injection:} Lexical injection patterns, structural anomaly detection, LLM self-rejection.
\item \textbf{T2 --- Privilege abuse:} Pre-inference role fast-fail, post-inference role$\times$intent policy.
\item \textbf{T3 --- Model hallucination:} Schema/regex/network/scope validation stack.
\item \textbf{T4 --- Out-of-scope targets:} CIDR containment, RFC-1918 check, scope mode enforcement.
\end{enumerate}

\subsubsection{Data Security}
\reqid{Requirement NFR-SEC-002:} Raw user input shall NEVER be stored in plaintext. Only SHA-256 hashes shall be recorded in the audit log.

\reqid{Requirement NFR-SEC-003:} API authentication shall use the \texttt{X-API-Key} header. The default development key (\texttt{dev\_insecure\_key}) shall NOT be used in production.

\reqid{Requirement NFR-SEC-004:} CORS shall be restricted to specific origins in production (current default: allow all, development mode only).

\subsubsection{Injection Defense Layers}
\begin{table}[h]
\caption{Injection defense layers.}
\label{tab:defense}
\begin{tabularx}{\columnwidth}{l >{\raggedright\arraybackslash}p{0.28\columnwidth} X}
\toprule
\textbf{Layer} & \textbf{Target} & \textbf{Mechanism}\\
\midrule
M1 AdversarialDetector & Raw input & 15 lexical + 7 structural + 3 heuristic patterns\\
S7 ScopeGuard & Raw input & 7 injection regexes (independent of M1)\\
S4 JSONParser & LLM output & Strip thought tags, fences\\
S5 SchemaValidator & Parsed JSON & Pydantic \texttt{extra="forbid"} (no extra fields)\\
S6 RegexValidator & Validated schema & Shell metacharacter scan\\
\bottomrule
\end{tabularx}
\end{table}

\subsection{Reliability and Availability}
\reqid{Requirement NFR-REL-001:} Non-critical middleware failures (session context, adversarial detection, cache store) shall degrade gracefully without aborting the pipeline.

\reqid{Requirement NFR-REL-002:} The audit logger shall be append-only with no delete operations.

\reqid{Requirement NFR-REL-003:} The semantic cache shall self-heal via grounding-hash invalidation when reference datasets change.

\reqid{Requirement NFR-REL-004:} Docker health checks shall verify IRE (30s interval) and Ollama (15s interval) connectivity.

\subsection{Maintainability}
\reqid{Requirement NFR-MNT-001:} All pipeline stages and middleware shall be individually configurable via feature flags in \texttt{config/features.yaml}.

\reqid{Requirement NFR-MNT-002:} Enforcement policy shall be environment-configurable (research/strict/production) without code changes.

\reqid{Requirement NFR-MNT-003:} The test suite shall achieve $\geq$ 80\% line coverage (current: 434 unit + 48 integration = 482 tests).

\subsection{Portability}
\reqid{Requirement NFR-PORT-001:} The IRE shall run on Linux, macOS, and Windows with Python 3.11+.

\reqid{Requirement NFR-PORT-002:} Docker deployment shall use the \texttt{python:3.11-slim} base image.

\reqid{Requirement NFR-PORT-003:} All configuration shall be environment-variable-driven via the \texttt{.env} file or Docker environment.
% =================== 6. DATA REQUIREMENTS ===================
\section{Data Requirements}

\subsection{Logical Data Model}

\subsubsection{Core Entities}
\begin{lstlisting}
IntentSchema
+-- intent: IntentType (9 values)
+-- target: Target {type: TargetType, value: str}
+-- ports: List[int]
+-- modifiers: List[str]
+-- cve_ids: List[str]
+-- tool_hint: Optional[str]
+-- schedule: Optional[str]
+-- confidence: float [0.0, 1.0]
+-- rejection_reason: Optional[str]

PlannerContract (extends IntentSchema)
+-- sub_intent: Optional[SubIntentType]
+-- scope_warnings: List[str]
+-- executability predicates

IREResponseV2 (extends IntentSchema)
+-- validation_findings: List[ValidationFinding]
+-- uncertainty metadata
+-- session context
+-- XAI metadata (optional)
\end{lstlisting}

\subsubsection{ValidationFinding}
\begin{lstlisting}
ValidationFinding
+-- validator: str
+-- passed: bool
+-- hallucination_class: Optional[HallucinationClass]
+-- detail: str
+-- severity: "block" | "warn" | "escalate"
\end{lstlisting}

\subsubsection{CacheEntry (Semantic Cache)}
\begin{lstlisting}
CacheEntry
+-- key_hash: str (SHA-256)
+-- embedding: List[float] (384-dim)
+-- response: IREResponse
+-- grounding_hash: str (SHA-256)
+-- hit_count: int
+-- created_at: float
+-- last_hit_at: float
\end{lstlisting}

\subsubsection{AuditRecord (SQLite)}
\begin{lstlisting}
audit_log
+-- id: INTEGER PRIMARY KEY
+-- timestamp: TEXT (ISO 8601)
+-- session_id: TEXT (nullable)
+-- raw_input_hash: TEXT (SHA-256)
+-- intent_summary: TEXT
+-- hallucination_class: TEXT
+-- severity: TEXT
+-- enforcement_action: TEXT
+-- status: TEXT ("recorded" | "pending" | "approved" | "rejected")
+-- resolved_by: TEXT
+-- resolved_at: TEXT
+-- resolution_note: TEXT
+-- cached_contract_json: TEXT
\end{lstlisting}

\subsection{Data Retention and Privacy}
\begin{table*}[t]
\caption{Data types, retention, and privacy.}
\label{tab:retention}
\begin{tabularx}{\textwidth}{l X X}
\toprule
\textbf{Data Type} & \textbf{Retention} & \textbf{Privacy}\\
\midrule
Raw user input & Transient (in-memory only) & Never stored\\
Audit log & Append-only, persistent & SHA-256 hashed input only\\
Semantic cache & LRU eviction (512 max) & In-memory + optional disk\\
Inference cache & LRU eviction (256 max) & Disk: \texttt{eval/.ollama\_inference\_cache.json}\\
Session context & 30-minute TTL, LRU eviction & In-memory only, never persisted\\
\bottomrule
\end{tabularx}
\end{table*}

\subsection{Grounding Data}

\subsubsection{CVE Grounding Dataset (\texttt{data/cve\_grounding.json})}
\begin{itemize}
\item \textbf{Purpose:} Reference dataset for CVE validation (Check 2 of RegexValidator);
\item \textbf{Format:} JSON with 20 curated entries plus a \texttt{\_meta} block;
\item \textbf{Schema per entry:} \texttt{\{id, description, affected\_service, severity\}};
\item \textbf{Refresh:} Manual via \texttt{scripts/update\_cve\allowbreak{}grounding.py} (CISA KEV catalog);
\item \textbf{Grounding Hash:} SHA-256 digest computed on each cache lookup; stale entries purged.
\end{itemize}

\subsubsection{Alias Map (\texttt{data/alias\_map.json})}
\begin{itemize}
\item \textbf{Purpose:} Map exploit names and slang to canonical CVE identifiers;
\item \textbf{Entries:} 10 exploit-name to CVE mappings + 8 slang/term normalizations;
\item \textbf{Format:} JSON object \texttt{\{alias: canonical, ...\}};
\item \textbf{Processing:} Word-boundary regex compilation, longest-first matching.
\end{itemize}

\subsection{Configuration Data}
\begin{table*}[t]
\caption{Configuration files.}
\label{tab:config}
\begin{tabularx}{\textwidth}{>{\ttfamily}l l X}
\toprule
\textbf{File} & \textbf{Format} & \textbf{Purpose}\\
\midrule
\texttt{config/settings.py} & pydantic-settings & Environment-driven configuration (30+ fields)\\
\texttt{config/features.yaml} & YAML & Feature flags (10) + tuning parameters (11)\\
\texttt{config/rbac\_policy.py} & Python & Role to intent permission table\\
\texttt{config/enforcement\_policy.py} & Python & HallucinationClass to action mapping\\
\texttt{config/lora\_config.yaml} & YAML & Phase-2 fine-tuning plan\\
\bottomrule
\end{tabularx}
\end{table*}

% =================== 7. VERIFICATION AND VALIDATION ===================
\section{Verification and Validation}

\subsection{Test Suite}
\begin{table*}[t]
\caption{Test suite summary.}
\label{tab:testsuite}
\begin{tabularx}{\textwidth}{l l p{3cm} X}
\toprule
\textbf{Category} & \textbf{Count} & \textbf{Files} & \textbf{Notable Coverage}\\
\midrule
Unit tests & 434 & 22 files under \texttt{tests/unit/} & adversarial (42), sub\_intent (46), regex (40), semantic\_cache (37), session\_context (35)\\
Integration tests & 48 & 5 files under \texttt{tests/integration/} & API v1 (12), API v2 (33), audit escalation (3)\\
\textbf{Total} & \textbf{482} & \textbf{27 files} & \\
\bottomrule
\end{tabularx}
\end{table*}

Key behavioral assertions:
\begin{itemize}
\item Tier-2 gating calls the LLM exactly once when confident (\texttt{test\_self\_consistency\_entropy.py});
\item Near-duplicate cache isolation and grounding invalidation (\texttt{test\_semantic\_cache\_adversarial.py});
\item IPv6 and production-mode edge coverage (\texttt{test\_uncovered\_source\_behaviors.py}).
\end{itemize}

\subsection{Evaluation Harness}
\begin{table*}[t]
\caption{Evaluation harness scripts.}
\label{tab:evalharness}
\begin{tabularx}{\textwidth}{>{\ttfamily}l X X}
\toprule
\textbf{Script} & \textbf{Purpose} & \textbf{Path}\\
\midrule
\texttt{diagnose.py} & Full-pipeline runner $\rightarrow$ JSONL debug rows + per-category report & \texttt{eval/diagnose.py}\\
\texttt{metrics.py} & Decision/accuracy/catch/cache metrics & \texttt{eval/metrics.py}\\
\texttt{category\_metrics.py} & NormalizedRow-based metrics engine (canonical) & \texttt{eval/category\_metrics.py}\\
\texttt{run\_baseline.py} & Raw unconstrained LLM baseline scorer & \texttt{eval/run\_baseline.py}\\
\texttt{exp1\_raw\_vs\_pipeline.py} & Exp1: raw vs. pipeline, 4 configs & \texttt{eval/Data Collection/}\\
\texttt{exp2\_model\_head\_to\_head.py} & Exp2: Gemma vs. Qwen & \texttt{eval/Data Collection/}\\
\texttt{exp3\_ablation\_study.py} & Exp3: 8-config $\times$ 900-record ablation & \texttt{eval/Data Collection/}\\
\texttt{exp4\_hallucination\_analysis.py} & Exp4: per-class deep-dive (250 records) & \texttt{eval/Data Collection/}\\
\bottomrule
\end{tabularx}
\end{table*}

\subsection{Empirical Results Summary (900-Record Canonical Runs)}

\subsubsection{Experiment 1 --- Raw LLM vs. Full Pipeline}
\begin{table}[h]
\caption{Experiment 1 results.}
\label{tab:exp1}
\begin{tabularx}{\columnwidth}{l X X X X}
\toprule
\textbf{Metric} & \textbf{Raw Gemma} & \textbf{Raw Qwen} & \textbf{Pipeline Gemma} & \textbf{Pipeline Qwen}\\
\midrule
Intent Accuracy & 78.0\% & 69.0\% & \textbf{86.0\%} & \textbf{82.0\%}\\
Hallucination Catch & 0.0\% & 0.0\% & \textbf{72.0\%} & \textbf{64.0\%}\\
False Positive Rate & 32.0\% & 31.0\% & \textbf{14.0\%} & \textbf{15.0\%}\\
\bottomrule
\end{tabularx}
\end{table}

\noindent\textbf{Pipeline lift:} +8pp/+13pp accuracy, +72pp/+64pp catch, $-18$pp/$-16$pp FP.

\subsubsection{Experiment 2 --- Model Head-to-Head}
\begin{table}[h]
\caption{Experiment 2 results.}
\label{tab:exp2}
\begin{tabularx}{\columnwidth}{l X X X}
\toprule
\textbf{Metric} & \textbf{Gemma 4} & \textbf{Qwen 2.5 Coder} & \textbf{Delta}\\
\midrule
Intent Accuracy & 86.0\% & 82.0\% & $-4.0$pp\\
Hallucination Catch & 72.0\% & 64.0\% & $-8.0$pp\\
FP Rate & 14.0\% & 15.0\% & $+1.0$pp\\
p50 Latency & 27,650 ms & 32,755 ms & $+18\%$\\
\bottomrule
\end{tabularx}
\end{table}

\subsubsection{Experiment 3 --- Component Ablation}
\begin{table*}[t]
\caption{Experiment 3 ablation results.}
\label{tab:exp3}
\begin{tabularx}{\textwidth}{l X X X X X}
\toprule
\textbf{Config} & \textbf{DAcc} & \textbf{Catch} & \textbf{FP} & \textbf{CFAcc} & \textbf{Mean Latency}\\
\midrule
\textbf{Full Pipeline} & \textbf{73.5\%} & \textbf{84\%} & 34\% & 87.5\% & 28,805 ms\\
No Scope Guard & 64.0\% & 73\% & 44\% & 83.9\% & 31,404 ms\\
No RBAC Guard & 74.5\% & 64\% & 22\% & 88.2\% & 28,260 ms\\
No Network Validator & 68.0\% & 82\% & 44\% & 83.9\% & 27,450 ms\\
No Hallucination Validator & 65.4\% & 62\% & 44\% & 82.1\% & 26,722 ms\\
Validation Stack Disabled & 64.0\% & 41\% & 13\% & 87.4\% & 25,198 ms\\
Self-Consistency Disabled & 67.5\% & 80\% & 45\% & 85.5\% & 11,809 ms\\
Semantic Cache Disabled & 66.0\% & 79\% & 47\% & 86.8\% & 22,949 ms\\
\bottomrule
\end{tabularx}
\end{table*}

\textbf{Component contribution ranking (by accuracy impact):}
\begin{enumerate}
\item ScopeGuard ($+9.5$pp);
\item Entire Validation Stack ($+9.5$pp);
\item SchemaValidator + RegexValidator ($+8.1$pp);
\item Semantic Cache ($+7.5$pp);
\item Tier-2 Self-Consistency ($+6.0$pp);
\item NetworkArchitectureValidator ($+5.5$pp);
\item RBACGuard ($-1.0$pp).
\end{enumerate}

\subsubsection{Experiment 4 --- Hallucination Deep-Dive (250 Records)}
\begin{table}[h]
\caption{Experiment 4 per-category catch rates.}
\label{tab:exp4}
\begin{tabularx}{\columnwidth}{l X X >{\raggedright\arraybackslash}p{0.15\columnwidth}}
\toprule
\textbf{Category} & \textbf{Total} & \textbf{Caught} & \textbf{Catch \%}\\
\midrule
rbac\_violation & 12 & 12 & \textbf{100.0\%}\\
adversarial\_injection & 45 & 43 & \textbf{95.6\%}\\
fabricated\_parameter & 32 & 28 & 87.5\%\\
target\_type\_mismatch & 32 & 26 & 81.2\%\\
fabricated\_cve & 32 & 25 & 78.1\%\\
contradictory\_action\_target & 32 & 19 & 59.4\%\\
out\_of\_scope & 45 & 21 & \textbf{46.7\%}\\
ambiguous & 20 & 8 & 40.0\%\\
\textbf{Total} & \textbf{250} & \textbf{182} & \textbf{72.8\%}\\
\bottomrule
\end{tabularx}
\end{table}

\textbf{Largest structural gap:} \texttt{llm\_no\_scope\_detection} (35\% of all misses) --- public/out-of-scope targets reach the LLM before any scope gate.

\subsection{Metric Definitions}
\begin{table}[h]
\caption{Metric definitions.}
\label{tab:metrics}
\begin{tabularx}{\columnwidth}{>{\bfseries}l X}
\toprule
\textbf{Metric} & \textbf{Formula}\\
\midrule
DAcc (Decision Accuracy) & \texttt{sum}/\texttt{N\_total} of correct decisions: well\_formed$\rightarrow$ALLOWED+correct fields, flawed$\rightarrow$BLOCKED/ESCALATED\\
Catch (Flawed-Input Catch) & \texttt{|flawed records with BLOCKED/ESCALATED|}/\texttt{N\_flawed}\\
FP (False Positive Rate) & \texttt{|well\_formed records with non-ALLOWED decision|}/\texttt{N\_well\_formed}\\
CFAcc (Contract-Field Accuracy) & \texttt{|ALLOWED records with correct intent+target|}/\texttt{|ALLOWED well\_formed|}\\
\bottomrule
\end{tabularx}
\end{table}
% =================== 8. DEPLOYMENT ===================
\section{Deployment}

\subsection{Docker Deployment (Recommended)}
\begin{lstlisting}
# 1. Clone and configure
cd apps/services/neuroshell-ire
cp .env.example .env
# Edit .env: set IRE_API_KEY, OLLAMA_MODEL

# 2. Build and start
docker-compose up -d --build

# 3. Verify
curl http://localhost:8001/health
\end{lstlisting}

Services:
\begin{table*}[t]
\caption{Docker services.}
\label{tab:docker}
\begin{tabularx}{\textwidth}{>{\ttfamily}l >{\ttfamily}l >{\ttfamily}l >{\raggedright\arraybackslash}X}
\toprule
\textbf{Service} & \textbf{Container} & \textbf{Port} & \textbf{Health Check}\\
\midrule
ire-service & neuroshell-ire & 8001 & \texttt{curl -f http://localhost:8001/health}\\
ollama & neuroshell-ollama & 11434 & \texttt{curl -f http://localhost:11434/api/tags}\\
\bottomrule
\end{tabularx}
\end{table*}

\noindent\textbf{Volumes:} \texttt{ollama\_models} mounted at \texttt{/root/.ollama} (persists downloaded models).

\subsection{Manual Deployment}
\begin{lstlisting}
# 1. Prerequisites
python --version  # 3.11+
ollama --version  # 0.32.x

# 2. Virtual environment
python -m venv .venv
.venv\Scripts\activate  # Windows
# source .venv/bin/activate  # Linux/macOS

# 3. Dependencies
pip install -r requirements.txt
python -m spacy download en_core_web_sm

# 4. Configuration
cp .env.example .env
# Edit .env: set IRE_API_KEY, OLLAMA_MODEL, OLLAMA_BASE_URL

# 5. Pull LLM model
ollama pull gemma4:e4b

# 6. Start server
python scripts/run_dev.py
\end{lstlisting}

\subsection{Configuration Management}

\subsubsection{Environment Variables (.env)}
\begin{table*}[t]
\caption{Environment variables.}
\label{tab:envvars}
\begin{tabularx}{\textwidth}{>{\ttfamily}l >{\ttfamily}l X}
\toprule
\textbf{Variable} & \textbf{Default} & \textbf{Description}\\
\midrule
\texttt{IRE\_API\_KEY} & \texttt{dev\_insecure\_key} & API authentication key\\
\texttt{OLLAMA\_BASE\_URL} & \texttt{http://localhost:11434} & Ollama daemon URL\\
\texttt{OLLAMA\_MODEL} & \texttt{gemma4:e4b} & Active LLM model\\
\texttt{SCOPE\_MODE} & \texttt{research} & Scope enforcement mode\\
\texttt{ENGAGEMENT\_SCOPE} & \texttt{192.168.0.0/16} & Authorised target CIDR\\
\texttt{LOG\_LEVEL} & \texttt{INFO} & Logging level\\
\texttt{PORT} & \texttt{8001} & Server port\\
\bottomrule
\end{tabularx}
\end{table*}

\subsubsection{Feature Flags (config/features.yaml)}
\begin{table}[h]
\caption{Feature flags.}
\label{tab:flags}
\begin{tabularx}{\columnwidth}{>{\ttfamily}l >{\ttfamily}l X}
\toprule
\textbf{Flag} & \textbf{Default} & \textbf{Description}\\
\midrule
\texttt{session\_context} & true & Multi-turn session memory\\
\texttt{rbac} & true & Role-based access control\\
\texttt{adversarial\_detection} & true & Adversarial input detection\\
\texttt{semantic\_cache} & true & Embedding-based response cache\\
\texttt{sub\_intent} & true & Hierarchical sub-intent classification\\
\texttt{schema\_v2} & true & V2 schema routing\\
\texttt{xai} & false & XAI token attribution\\
\texttt{uncertainty\_estimation} & false & Monte Carlo dropout\\
\texttt{continual\_learning} & false & EWC active learning\\
\texttt{speculative\_decoding} & false & Draft model acceleration\\
\bottomrule
\end{tabularx}
\end{table}

% =================== 9. APPENDICES ===================
\appendices

\section{IntentType Enumeration (9 Values)}
\begin{table*}[t]
\caption{Intent types.}
\label{tab:appA}
\begin{tabularx}{\textwidth}{>{\ttfamily}l X >{\raggedright\arraybackslash}p{0.4\textwidth}}
\toprule
\textbf{Value} & \textbf{Category} & \textbf{Description}\\
\midrule
NETWORK\_SCAN & Active Recon & Port scanning, host discovery, OS detection\\
VULNERABILITY\_AUDIT & Assessment & CVE-specific or service-specific vulnerability checks\\
DIRECTORY\_BRUTEFORCE & Enumeration & Web directory/file brute-forcing\\
SERVICE\_ENUMERATION & Enumeration & Service version detection, banner grabbing\\
EXPLOITATION & Attack & Remote code execution, privilege escalation\\
PASSWORD\_ATTACK & Attack & Brute force, credential stuffing, password spray\\
PASSIVE\_RECON & Recon & DNS lookup, WHOIS, OSINT gathering\\
AMBIGUOUS & Special & Confident $<$ 0.5; requires clarification\\
REJECTED & Special & Command rejected (injection, policy, etc.)\\
\bottomrule
\end{tabularx}
\end{table*}

\section{TargetType Enumeration (6 Values)}
\begin{table*}[t]
\caption{Target types.}
\label{tab:appB}
\begin{tabularx}{\textwidth}{>{\ttfamily}l X X}
\toprule
\textbf{Value} & \textbf{Format} & \textbf{Example}\\
\midrule
IP & IPv4 address & \texttt{192.168.1.1}\\
SUBNET & CIDR notation & \texttt{192.168.1.0/24}\\
DOMAIN & Domain name & \texttt{example.com}\\
URL & HTTP/HTTPS URL & \texttt{https://example.com/api}\\
HOSTNAME & Single-label hostname & \texttt{webserver01}\\
UNKNOWN & Empty/unspecified & (empty string)\\
\bottomrule
\end{tabularx}
\end{table*}

\section{SubIntentType Enumeration (27 Values)}
\begin{table}[h]
\caption{Sub-intent types by parent intent.}
\label{tab:appC}
\begin{tabularx}{\columnwidth}{>{\ttfamily}l >{\ttfamily}X}
\toprule
\textbf{Parent Intent} & \textbf{Sub-Intents}\\
\midrule
NETWORK\_SCAN & ICMP\_DISCOVERY, SYN\_STEALTH, UDP\_SWEEP, VERSION\_DETECT, OS\_DETECT, FULL\_PORT\\
EXPLOITATION & RCE, PRIV\_ESC, LATERAL, PERSISTENCE\\
VULNERABILITY\_AUDIT & CVE\_SPECIFIC, SERVICE\_SPECIFIC, FULL\_SCAN\\
SERVICE\_ENUMERATION & SMB, SNMP, VERSION\_DETECT, BANNER\_GRAB\\
PASSWORD\_ATTACK & SPRAY, CREDENTIAL\_STUFFING, BRUTE\_FORCE\\
PASSIVE\_RECON & DNS, WHOIS, OSINT\\
DIRECTORY\_BRUTEFORCE & API, FILES, WEB\\
\bottomrule
\end{tabularx}
\end{table}

\section{HallucinationClass Enumeration (6 Values)}
\begin{table}[h]
\caption{Hallucination classes and enforcement.}
\label{tab:appD}
\begin{tabularx}{\columnwidth}{>{\ttfamily}l X X}
\toprule
\textbf{Class} & \textbf{Detection} & \textbf{Default Enforcement}\\
\midrule
FABRICATED\_CVE & S5 + S6 (syntax) & BLOCK\\
TARGET\_TYPE\_MISMATCH & S5 + S6 (format) + S6B & BLOCK\\
CONTRADICTORY\_ACTION\_TARGET & S5 (schema) & BLOCK\\
FABRICATED\_PARAMETER & S5 + S6 (port/metachar) & BLOCK\\
UNGROUNDED\_CONFIDENCE & S5 (confidence) + S6 (grounding) & WARN\\
OUT\_OF\_SCOPE\_TARGET & S6B (scope) + S7 (RFC-1918) & Dynamic (BLOCK/WARN)\\
\bottomrule
\end{tabularx}
\end{table}

\section{RBAC Policy Table}
\begin{table*}[t]
\caption{RBAC policy matrix.}
\label{tab:appE}
\begin{tabularx}{\textwidth}{>{\ttfamily}l c c c c}
\toprule
\textbf{Intent} & \textbf{viewer} & \textbf{analyst} & \textbf{operator} & \textbf{admin}\\
\midrule
NETWORK\_SCAN & --- & \checkmark & \checkmark & \checkmark\\
VULNERABILITY\_AUDIT & --- & \checkmark & \checkmark & \checkmark\\
DIRECTORY\_BRUTEFORCE & --- & --- & \checkmark & \checkmark\\
SERVICE\_ENUMERATION & --- & \checkmark & \checkmark & \checkmark\\
EXPLOITATION & --- & --- & \checkmark & \checkmark\\
PASSWORD\_ATTACK & --- & --- & \checkmark & \checkmark\\
PASSIVE\_RECON & --- & \checkmark & \checkmark & \checkmark\\
AMBIGUOUS & --- & \checkmark & \checkmark & \checkmark\\
REJECTED & --- & \checkmark & \checkmark & \checkmark\\
\bottomrule
\end{tabularx}
\end{table*}

\noindent\textbf{Sensitive intents (require operator+):} EXPLOITATION, PASSWORD\_ATTACK, DIRECTORY\_BRUTEFORCE.
\section{Enforcement Policy Table}
\begin{table*}[t]
\caption{Enforcement policy by environment mode.}
\label{tab:appF}
\begin{tabularx}{\textwidth}{>{\ttfamily}l l l l}
\toprule
\textbf{HallucinationClass} & \textbf{research} & \textbf{strict} & \textbf{production}\\
\midrule
FABRICATED\_PARAMETER & BLOCK & BLOCK & BLOCK\\
CONTRADICTORY\_ACTION\_TARGET & BLOCK & BLOCK & BLOCK\\
TARGET\_TYPE\_MISMATCH & BLOCK & BLOCK & BLOCK\\
FABRICATED\_CVE & BLOCK & BLOCK & BLOCK\\
UNGROUNDED\_CONFIDENCE & WARN & WARN & WARN\\
OUT\_OF\_SCOPE\_TARGET & WARN & BLOCK & BLOCK\\
\bottomrule
\end{tabularx}
\end{table*}

\section{Adversarial Detection Patterns}

\subsection{Layer A: Lexical (15 patterns)}
\begin{table*}[t]
\caption{Layer A lexical patterns.}
\label{tab:appG1}
\begin{tabularx}{\textwidth}{>{\ttfamily}X >{\raggedright\arraybackslash}p{0.08\textwidth}}
\toprule
\textbf{Pattern} & \textbf{Weight}\\
\midrule
\texttt{ignore (all~)?(previous|prior|above) instructions?} & 0.9\\
\texttt{forget (everything|all|your instructions?)} & 0.9\\
\texttt{you are now (?!scanning|checking|running)} & 0.8\\
\texttt{act as (a~)?(different|new|another|evil|unrestricted)} & 0.8\\
\texttt{jailbreak} & 1.0\\
\texttt{dan mode} & 1.0\\
\texttt{prompt injection} & 0.9\\
\texttt{(system|assistant|user):\textbackslash{}s*you (must|shall|will)} & 0.85\\
\texttt{$<$|im\_(start|end)|$>$} & 1.0\\
\texttt{$<$|think|$>$} & 0.95\\
\texttt{[INST]} or \texttt{[/INST]} & 0.9\\
\texttt{\#\#\# (instruction|system|human|assistant)} & 0.85\\
\texttt{disregard (all~)?(prior|previous|your)} & 0.85\\
\texttt{new instructions:} & 0.8\\
\texttt{override (previous~)?(instructions?|rules?|guidelines?)} & 0.85\\
\bottomrule
\end{tabularx}
\end{table*}

\subsection{Layer B: Structural (7 patterns + 3 heuristics)}
\begin{table*}[t]
\caption{Layer B structural patterns.}
\label{tab:appG2}
\begin{tabularx}{\textwidth}{>{\ttfamily}X >{\raggedright\arraybackslash}p{0.08\textwidth}}
\toprule
\textbf{Pattern} & \textbf{Weight}\\
\midrule
\texttt{(\textbackslash{}x[0-9a-f]\{2\})\{4,\}} & 0.6\\
\texttt{(\%[0-9a-f]\{2\})\{4,\}} & 0.5\\
\texttt{base64:\textbackslash{}s*[A-Za-z0-9+/]\{20,\}} & 0.7\\
\texttt{$<$script[\textbackslash{}s$>$]} & 0.8\\
\texttt{\textbackslash{}\$\{[\^{}]\}\{0,50\}\textbackslash{}\}} & 0.6\\
\texttt{eval( |exec(|\_\_import\_\_()} & 0.9\\
\texttt{rm -rf|del /[sq]|format c:} & 0.8\\
\bottomrule
\end{tabularx}
\end{table*}

\begin{table}[h]
\caption{Layer B heuristics.}
\label{tab:appG3}
\begin{tabularx}{\columnwidth}{l X X}
\toprule
\textbf{Heuristic} & \textbf{Threshold} & \textbf{Weight}\\
\midrule
Special-char density & $>$ 30\% (for inputs $>$ 20 chars) & 0.4\\
Single token length & $>$ 100 chars & 0.3\\
Repeated n-gram & $>$ 3 occurrences of 3/4/5-grams ($\geq$ 10 chars) & 0.35\\
\bottomrule
\end{tabularx}
\end{table}

\section{Configuration Parameters (Complete)}
\begin{table*}[t]
\caption{All configuration parameters.}
\label{tab:appH}
\begin{tabularx}{\textwidth}{>{\ttfamily}l X l l}
\toprule
\textbf{Parameter} & \textbf{Default} & \textbf{Type} & \textbf{Source}\\
\midrule
\texttt{ire\_api\_key} & \texttt{dev\_insecure\_key} & str & .env\\
\texttt{port} & 8001 & int & .env\\
\texttt{ollama\_base\_url} & \texttt{http://localhost:11434} & str & .env\\
\texttt{ollama\_model} & \texttt{gemma4:e4b} & str & .env\\
\texttt{ollama\_timeout\_seconds} & 60 & int & .env\\
\texttt{inference\_temperature} & 0.1 & float & .env\\
\texttt{inference\_max\_tokens} & 2048 & int & .env\\
\texttt{self\_consistency\_trigger\_threshold} & 0.7 & float & .env\\
\texttt{self\_consistency\_samples} & 5 & int & .env\\
\texttt{self\_consistency\_temperature} & 0.7 & float & .env\\
\texttt{entropy\_low\_threshold} & 0.3 & float & .env\\
\texttt{entropy\_high\_threshold} & 0.8 & float & .env\\
\texttt{scope\_mode} & research & literal & .env\\
\texttt{engagement\_scope} & \texttt{192.168.0.0/16} & str & .env\\
\texttt{engagement\_name} & Lab Environment & str & .env\\
\texttt{engagement\_mode} & warn & literal & .env\\
\texttt{max\_cidr\_prefix} & 30 & int & .env\\
\texttt{min\_cidr\_prefix} & 32 & int & .env\\
\texttt{lru\_cache\_max\_size} & 256 & int & .env\\
\texttt{planner\_enabled} & false & bool & .env\\
\texttt{planner\_url} & \texttt{http://127.0.0.1:8002} & str & .env\\
\texttt{planner\_api\_key} & \texttt{neuroshell-c2-secret-key} & str & .env\\
\texttt{planner\_timeout\_seconds} & 300 & int & .env\\
\texttt{component3\_enabled} & false & bool & .env\\
\texttt{component3\_url} & \texttt{http://127.0.0.1:8003} & str & .env\\
\texttt{component4\_enabled} & false & bool & .env\\
\texttt{component4\_url} & \texttt{http://127.0.0.1:8004} & str & .env\\
\texttt{session\_ttl\_minutes} & 30 & int & features.yaml\\
\texttt{session\_max\_turns} & 8 & int & features.yaml\\
\texttt{session\_max\_sessions} & 1024 & int & features.yaml\\
\texttt{semantic\_cache\_threshold} & 0.82 & float & features.yaml\\
\texttt{semantic\_cache\_max\_size} & 512 & int & features.yaml\\
\texttt{adversarial\_classifier\_threshold} & 0.7 & float & features.yaml\\
\texttt{sub\_intent\_confidence\_threshold} & 0.6 & float & features.yaml\\
\bottomrule
\end{tabularx}
\end{table*}

\section{Exception Hierarchy}
\begin{lstlisting}
IREError (base)
+-- JSONParseError          (stage: "json_parser")
+-- SchemaValidationError   (stage: "schema_validator")
+-- RegexValidationError    (stage: "regex_validator")
+-- NetworkValidationError  (stage: "network_validator")
|   +-- extra: validation_type
+-- ScopeError              (stage: "scope_guard")
+-- InferenceError          (stage: "inference")
+-- RBACError               (stage: "rbac_guard")
    +-- extra: role, denied_intent
\end{lstlisting}

\section{Dependency Versions}
\begin{table*}[t]
\caption{Package dependency versions.}
\label{tab:appJ}
\begin{tabularx}{\textwidth}{>{\ttfamily}l >{\ttfamily}l X}
\toprule
\textbf{Package} & \textbf{Version} & \textbf{Purpose}\\
\midrule
\texttt{fastapi} & 0.115.0 & Async REST framework\\
\texttt{uvicorn[standard]} & 0.30.6 & ASGI server\\
\texttt{pydantic} & $\geq$2.9, $<$3.0 & Schema validation\\
\texttt{pydantic-settings} & 2.4.0 & Environment config\\
\texttt{jsonschema} & 4.23.0 & JSON Schema support\\
\texttt{spacy} & 3.7.6 & NLP preprocessing\\
\texttt{ollama} & 0.3.3 & LLM client\\
\texttt{httpx} & 0.27.2 & HTTP client\\
\texttt{structlog} & 24.4.0 & Structured logging\\
\texttt{sentence-transformers} & 3.0.1 & Semantic cache embeddings\\
\texttt{slowapi} & 0.1.9 & Rate limiting\\
\texttt{pytest} & 8.3.2 & Test framework\\
\texttt{pytest-asyncio} & 0.23.8 & Async test support\\
\bottomrule
\end{tabularx}
\end{table*}

\section{Hashing Algorithms}
\begin{table*}[t]
\caption{Hash algorithm uses.}
\label{tab:appK}
\begin{tabularx}{\textwidth}{l l X}
\toprule
\textbf{Use} & \textbf{Algorithm} & \textbf{Input}\\
\midrule
Audit input privacy & SHA-256 & Raw user input\\
Semantic cache exact key & SHA-256 & Lowercased, stripped input text\\
Inference cache key & SHA-256 & \texttt{model\_name|enriched\_input}\\
Grounding invalidation & SHA-256 & Concatenation of \texttt{cve\_grounding.json} + \texttt{alias\_map.json}\\
\bottomrule
\end{tabularx}
\end{table*}

\vspace{1em}
\noindent\textit{--- End of Document ---}

\srsdivider{Component 02 --- Dynamic Planner \& RAG Engine}
\section{Introduction}

\subsection{Purpose}

This Software Requirements Specification (SRS) defines the complete functional and non-functional requirements for the \textbf{NeuroShell Dynamic Planner \& RAG Engine}, designated Component 02 of the NeuroShell Framework. The Dynamic Planner (also referred to as "C2") is a retrieval-augmented generation (RAG) service that converts validated intent contracts from the Component 01 Intent Recognition Engine (IRE) into \textbf{tool-grounded, safety-checked Kali Linux command plans} that are forwarded to Component 03 (AEERE) for execution.

This document is intended for:

\begin{itemize}
\item \textbf{Research evaluators} assessing the C2 design and empirical performance
\item \textbf{Developers} maintaining or extending the RAG pipeline and endpoint surface
\item \textbf{Downstream consumers} (Component 03 AEERE; Component 01 IRE chain orchestration) integrating with the C2 API
\item \textbf{Examiners} reviewing the MSc Cybersecurity research project (ref R26-CS-018)
\end{itemize}

\subsection{Scope}

\subsubsection{Product Scope}

The Dynamic Planner is a \textbf{self-contained FastAPI microservice} that:

\begin{enumerate}
\item Accepts a validated \texttt{PlanRequest} (enriched Flat planner contract or nested \texttt{IREIntentContract}) over REST
\item Interprets the intent contract into flat planning parameters (intent, target, ports, modifiers, CVE IDs, tool hint)
\item Retrieves tool documentation from a local ChromaDB vector knowledge base via semantic embedding
\item Synthesizes an executable Kali Linux command (or up to 5-step command sequence) through a grounded LLM call (Ollama)
\item Validates every synthesized command through a safety filter and a structural command validator
\item Maintains per-session context in Redis (with in-memory fallback) and records operational metrics
\item \textbf{Never executes} any command --- execution is delegated exclusively to Component 03 (AEERE)
\end{enumerate}

\subsubsection{In Scope}

\begin{itemize}
\item All source code under \texttt{app/} and \texttt{src/} (wrapper app, intent interpretation, RAG retrieval, context assembly, command synthesis, tool selection, safety/validation, session, metrics)
\item Configuration (\texttt{app/config.py}, \texttt{.env}, \texttt{.env.example})
\item Knowledge base (\texttt{knowledge\_base/} raw docs + ChromaDB vector store) and ingestion script (\texttt{scripts/ingest\_docs.py})
\item Tool mapping (\texttt{data/tool\_map.json})
\item Test suite under \texttt{tests/} (unit and integration)
\item Companion documentation under \texttt{docs/}
\item Deployment artifacts (Dockerfile, Modelfile, docker-compose, runtime scripts)
\end{itemize}

\subsubsection{Out of Scope}

\begin{itemize}
\item Component 01 (IRE), Component 03 (AEERE Executor), Component 04 (AVAE Vulnerability Analyzer) --- referenced only for interface design
\item Actual command execution / shell interaction --- \textbf{explicitly out of scope by design}
\item Frontend UI (separate project)
\item LoRA fine-tuning pipeline execution (adapter artifact \texttt{models/\allowbreak gemma2-\allowbreak c2-\allowbreak lora.\allowbreak gguf} is an external deployment dependency, not shipped in-repo)
\item Network infrastructure and physical security
\end{itemize}

\subsection{Definitions, Acronyms, Abbreviations}

\begin{table}[h]
\caption{Table 1.}
\label{tab:t1}
\begin{tabularx}{\columnwidth}{l >{\raggedright\arraybackslash}X}
\toprule
\textbf{Term} & \textbf{Definition}\\
\midrule
C1/C2/C3/C4 & Component 01 (IRE) / 02 (Planner) / 03 (AEERE) / 04 (AVAE)\\
RAG & Retrieval-Augmented Generation\\
LLM & Large Language Model\\
RAG KB & Vector knowledge base (ChromaDB, collection \texttt{tool\_documentation})\\
ChromaDB & Embedded vector database used for semantic retrieval\\
Ollama & Local LLM inference server used for command synthesis\\
KB & Knowledge base\\
IREIntentContract & C1$\rightarrow$C2 nested intent contract schema\\
PlannerContract & C1 output intent contract (flat mapping into the C2 \texttt{PlanRequest})\\
PlanRequest & C2 inbound request schema (format A: flat; format B: nested)\\
PlannerOutput & C2 outbound plan schema consumed by C3\\
Tool hint & Preferred tool supplied by the caller or resolved by RAG\\
Multi-step & A plan mode that synthesizes a sequence of up to 5 commands\\
RFC-1918 & Address Allocation for Private Internets\\
E2E latency & End-to-end latency recorded by \texttt{MetricsCollector} per plan request\\
Vector retriever & Semantically retrieves top-k documentation chunks (default top\_k=3)\\
Score threshold & Cosine-distance derived similarity threshold (default 0.30)\\
\bottomrule
\end{tabularx}
\end{table}

\subsection{References}

\begin{table}[h]
\caption{Table 2.}
\label{tab:t2}
\begin{tabularx}{\columnwidth}{l >{\raggedright\arraybackslash}X >{\raggedright\arraybackslash}X}
\toprule
\textbf{Ref ID} & \textbf{Document} & \textbf{Location}\\
\midrule
R1 & Component 02 README / system docs & \texttt{README.md}, \texttt{docs/01-15}\\
R2 & IRE ↔ C2 Integration Analysis & \texttt{neuroshell-\allowbreak ire/\allowbreak IRE-\allowbreak Doc/\allowbreak Component 02 Integration Analysis.\allowbreak md}\\
R3 & C2 session / metrics / admin API reference & \texttt{docs/\allowbreak 02\_\allowbreak api\_\allowbreak reference.\allowbreak md}\\
R4 & RAG subsystem design & \texttt{docs/\allowbreak 06\_\allowbreak rag\_\allowbreak subsystem.\allowbreak md}\\
R5 & Safety \& validation design & \texttt{docs/\allowbreak 07\_\allowbreak safety\_\allowbreak validation.\allowbreak md}\\
R6 & Configuration \& deployment & \texttt{docs/\allowbreak 10\_\allowbreak configuration\_\allowbreak deployment.\allowbreak md}\\
R7 & Evaluation report (benchmarks) & \texttt{docs/\allowbreak 15\_\allowbreak evaluation\_\allowbreak report.\allowbreak md}\\
R8 & C3 AEERE SRS (downstream consumer) & \texttt{adaptive-\allowbreak execution-\allowbreak error-\allowbreak recovery/\allowbreak SRS-\allowbreak Doc/\allowbreak }\\
R9 & IEEE 830-1998 SRS Standard & IEEE Computer Society\\
R10 & RFC 1918 --- Private Address Space & IETF\\
\bottomrule
\end{tabularx}
\end{table}

\subsection{Overview}

The remainder of this document is organized as follows:

\begin{itemize}
\item \textbf{Section 2} describes the overall product perspective, user classes, and operating environment
\item \textbf{Section 3} specifies all system features and functional requirements in detail
\item \textbf{Section 4} defines external interface requirements (API, hardware, software, communication)
\item \textbf{Section 5} specifies non-functional requirements (performance, security, reliability)
\item \textbf{Section 6} covers data requirements (models, RAG KB, sessions, metrics)
\item \textbf{Section 7} describes verification, validation, and empirical results
\item \textbf{Section 8} covers deployment procedures
\item \textbf{Section 9} provides appendices with enumerations, templates, and schemas
\end{itemize}

\section{Overall Description}

\subsection{Product Perspective}

The NeuroShell Dynamic Planner is \textbf{Component 02} of the NeuroShell Framework, an autonomous natural-language-driven offensive security orchestration framework. C2 is an \textbf{independently deployable microservice} that acts as the command-synthesis layer between intent recognition (C1) and command execution (C3).

\begin{lstlisting}
User natural language request
        │
        ▼
┌──────────────────────────────────────────────┐
│ C1 IRE (Intent Recognition Engine)            │
│  POST /parse/plan → PlannerContract           │
└───────────────────┬──────────────────────────┘
                    │ PlanRequest (flat or nested)
                    ▼
┌──────────────────────────────────────────────┐
│ C2 Dynamic Planner & RAG Engine  (this doc)  │
│  POST /plan → PlannerOutput                  │
└───────────────────┬──────────────────────────┘
                    │ PlannerOutput (command, command_sequence, tool, ...)
                    ▼
┌──────────────────────────────────────────────┐
│ C3 AEERE (Execution & Error Recovery)         │
│  POST /execute → C3ExecutionResult            │
└───────────────────┬──────────────────────────┘
                    ▼
┌──────────────────────────────────────────────┐
│ C4 AVAE (Vulnerability Analysis Engine)       │
│  POST /analyze                                │
└──────────────────────────────────────────────┘
\end{lstlisting}

C2 \textbf{performs no execution}: it produces validated, tool-grounded command plans. This separation enforces a strict safety boundary: an LLM may propose commands, but the safety filter and command validator gate their release.

\subsection{Product Functions Summary}

\begin{itemize}
\item \textbf{Intent interpretation}: flattens \texttt{IREIntentContract} / \texttt{PlannerContract} into planning parameters; resolves default ports; computes multi-step mode
\item \textbf{Semantic RAG retrieval}: sentence-transformers embeddings (\texttt{all-MiniLM-L6-v2}, 384-dim) over ChromaDB collection \texttt{tool\_documentation}
\item \textbf{Tool selection}: priority resolution from C1 tool hint $\rightarrow$ top RAG retrieval source $\rightarrow$ \texttt{data/tool\_map.json} fallback
\item \textbf{Grounded command synthesis}: Ollama LLM synthesis constrained by system prompt, retrieved documentation, and target reinforcement
\item \textbf{Command sequence support}: multi-step synthesis (up to 5 commands) for \texttt{EXPLOITATION} / \texttt{VULNERABILITY\_AUDIT} when confidence > 0.85
\item \textbf{Dual-layer validation}: \texttt{SafetyFilter} (danger patterns + target scope) and \texttt{CommandValidator} (tool regex + required flags)
\item \textbf{Session context}: Redis-backed per-session targets/commands/tools with 3600 s TTL and in-memory fallback
\item \textbf{Metrics \& administration}: latency metrics, RBAC-separated admin endpoints (\texttt{/kb/update}, \texttt{/admin/*}), health/status endpoints
\end{itemize}

\subsection{User Classes and Characteristics}

\begin{table}[h]
\caption{Table 3.}
\label{tab:t3}
\begin{tabularx}{\columnwidth}{l >{\raggedright\arraybackslash}X >{\raggedright\arraybackslash}X}
\toprule
\textbf{User class} & \textbf{Description} & \textbf{Interactions}\\
\midrule
\textbf{End users / operator} & Human operator via frontend or direct REST & \texttt{POST /plan} (via gateway)\\
\textbf{Upstream C1 (IRE)} & Chain orchestration calling C2 within \texttt{/execute} & \texttt{POST /plan}\\
\textbf{Downstream C3 (AEERE)} & Consumes \texttt{PlannerOutput} command plans & (outbound hand-off, C3 side)\\
\textbf{Administrator} & Manages the knowledge base and runtime state & \texttt{/kb/update}, \texttt{/admin/*}, \texttt{/metrics}\\
\textbf{Maintainer / developer} & Extends the pipeline, RAG KB, or validation rules & source, tests, \texttt{scripts/ingest\_docs.py}\\
\bottomrule
\end{tabularx}
\end{table}

\subsection{Operating Environment}

\begin{itemize}
\item \textbf{Runtime}: Python 3.9+ (deployed image \texttt{python:3.11-slim}), FastAPI + Uvicorn ASGI
\item \textbf{Deployment}: Docker container (\texttt{neuroshell-\allowbreak dynamic-\allowbreak planner-\allowbreak rag}), orchestrated behind Kong API gateway (route \texttt{/\allowbreak api/\allowbreak dynamic-\allowbreak planner-\allowbreak rag/\allowbreak *})
\item \textbf{Dependencies}: ChromaDB (embedded persistent), Ollama (local LLM server), Redis (sessions, optional with fallback), SentenceTransformer (\texttt{all-MiniLM-L6-v2})
\item \textbf{Ports}: container 8000; host mapping \texttt{8003:8000} (docker-compose); dev standalone on 8002 (\texttt{src/api/main.py})
\item \textbf{OS}: Linux (production); Windows/macOS supported for development (\texttt{neuroshell\_run.ps1})
\end{itemize}

\subsection{Design and Implementation Constraints}

\begin{itemize}
\item \textbf{No execution}: C2 must never feed commands to a shell; execution is C3's responsibility
\item \textbf{LLM grounding}: every synthesized command must be traceable to retrieved documentation chunks (\texttt{retrieval\_sources} in \texttt{PlannerOutput})
\item \textbf{Safety-first gating}: any synthesized command containing a blocked pattern or targeting a public address is rejected before release
\item \textbf{Deterministic-first}: intent interpretation, tool selection, and validation follow explicit rule logic; the LLM is used only for synthesis
\item \textbf{Version pinning}: package versions pinned via Poetry \texttt{pyproject.toml} (\texttt{fastapi \textasciicircum\{\}0.\allowbreak 115.\allowbreak 0}, \texttt{chromadb 0.5.0}, \texttt{ollama 0.2.0}, etc.)
\end{itemize}

\subsection{Assumptions and Dependencies}

\begin{itemize}
\item Ollama is reachable and hosts the configured model (\texttt{OLLAMA\_MODEL}, default \texttt{gemma4:e4b}; tuned \texttt{neuroshell-\allowbreak c2-\allowbreak gemma2:\allowbreak latest})
\item ChromaDB knowledge base has been seeded via \texttt{scripts/ingest\_docs.py}
\item Redis (\texttt{REDIS\_URL}) is optional: on failure C2 degrades to in-memory session storage
\item LoRA adapter (\texttt{models/\allowbreak gemma2-\allowbreak c2-\allowbreak lora.\allowbreak gguf}) is provided externally at deployment (not committed to VCS)
\end{itemize}

\section{System Features and Functional Requirements}

\subsection{Feature: Plan Orchestration (POST /plan)}

\subsubsection{Description}

The core endpoint accepts an intent contract, runs the 9-step pipeline (authentication $\rightarrow$ contract ingest $\rightarrow$ intent interpretation $\rightarrow$ query construction $\rightarrow$ RAG retrieval $\rightarrow$ tool selection $\rightarrow$ context assembly $\rightarrow$ LLM synthesis $\rightarrow$ per-command safety + structural validation $\rightarrow$ session update $\rightarrow$ metrics), and returns a \texttt{PlannerOutput}.

\reqid{Requirement FR-PLAN-001:} The service shall expose a \texttt{POST /plan} endpoint that authenticates the caller via the \texttt{X-API-Key} header and returns HTTP 401 with \texttt{\{"detail":\allowbreak  "Invalid API key"\}} on mismatch.

\reqid{Requirement FR-PLAN-002:} The service shall accept both enriched flat requests (extra fields allowed, with convenience property shims) and nested \texttt{intent\_contract} requests (\texttt{IREIntentContract}).

\reqid{Requirement FR-PLAN-003:} A rejected intent (\texttt{status == REJECTED}) shall short-circuit to HTTP 400 \texttt{\{"stage":\allowbreak "intent\_\allowbreak filter","error":\allowbreak "INTENT\_\allowbreak REJECTED"\}} \textbf{before} any RAG retrieval or LLM call, and the attempt shall be recorded with \texttt{status="REJECTED"} in metrics.

\reqid{Requirement FR-PLAN-004:} The pipeline shall process requests through the sequence: interpret $\rightarrow$ construct query $\rightarrow$ retrieve (top\_k, threshold) $\rightarrow$ select tool $\rightarrow$ assemble prompt $\rightarrow$ synthesize $\rightarrow$ validate $\rightarrow$ session update $\rightarrow$ metrics.

\reqid{Requirement FR-PLAN-005:} On success the endpoint shall return HTTP 200 with \texttt{PlannerOutput} populated with \texttt{status="success"}, \texttt{command}, \texttt{command\_sequence}, \texttt{tool}, \texttt{session\_id}, \texttt{intent\_ref}, \texttt{estimated\_\allowbreak duration="medium"}, \texttt{retrieval\_sources}, \texttt{validation\_passed=true}, and \texttt{latency\_ms}.

\reqid{Requirement FR-PLAN-006:} Commands failing the safety filter shall return HTTP 400; commands failing structural validation shall return HTTP 422 with stage/error/detail/latency; unexpected exceptions shall return HTTP 500 with a recorded failed metric.

\reqid{Requirement FR-PLAN-007:} The inbound request \texttt{session\_id} shall be mandatorily present; empty target values shall be recovered from the session store's last target when marked \texttt{UNKNOWN}.

\subsection{Feature: Intent Interpretation}

\subsubsection{Description}

\texttt{IntentInterpreter.\allowbreak interpret(contract)} converts a contract into a flat parameter dictionary used by the query constructor and prompt builder.

\reqid{Requirement FR-INTENT-001:} The interpreter shall extract \texttt{intent, target\_\allowbreak value, target\_\allowbreak type, ports, modifiers, cve\_\allowbreak ids, tool\_\allowbreak hint, confidence, wordlist, suggested\_\allowbreak command, multi\_\allowbreak step}.

\reqid{Requirement FR-INTENT-002:} \texttt{multi\_step} shall be \texttt{True} only when the intent is in \texttt{MULTI\_\allowbreak STEP\_\allowbreak INTENTS = \{"EXPLOITATION", "VULNERABILITY\_\allowbreak AUDIT"\}} and the contract confidence exceeds 0.85.

\reqid{Requirement FR-INTENT-003:} Default ports shall be applied when ports are empty: \texttt{[80]} for \texttt{VULNERABILITY\_AUDIT} (nikto) and \texttt{DIRECTORY\_BRUTEFORCE}; empty list otherwise.

\reqid{Requirement FR-INTENT-004:} Tool hint priority shall be \texttt{primary\_tool} (contract) then \texttt{tool\_hint} (contract).

\reqid{Requirement FR-INTENT-005:} The interpreter shall accept both flat (enriched) and nested (\texttt{intent\_contract}) forms without behavioral difference.

\subsection{Feature: RAG Query Construction and Retrieval}

\subsubsection{Description}

\texttt{QueryConstructor} builds a retrieval query string from parameters using intent-specific templates and modifier text. \texttt{VectorRetriever} embeds the query and retrieves top-k documentation chunks from the local ChromaDB store.

\reqid{Requirement FR-RAG-001:} \texttt{QueryConstructor.\allowbreak build\_\allowbreak query} shall select a template from \texttt{QUERY\_TEMPLATES} for the detected intent (7 intent templates plus a default).

\reqid{Requirement FR-RAG-002:} The query shall expand any recognized modifiers against \texttt{MODIFIER\_TEXT\_MAP} (16 modifiers including \texttt{stealth}, \texttt{aggressive}, \texttt{all-ports}, \texttt{wordlist:common}, etc.).

\reqid{Requirement FR-RAG-003:} \texttt{VectorRetriever.\allowbreak retrieve(query, tool\_\allowbreak hint=None, top\_\allowbreak k=3)} shall embed the query with \texttt{all-MiniLM-L6-v2} and query the \texttt{tool\_documentation} collection using cosine space.

\reqid{Requirement FR-RAG-004:} Retrieval results shall be filtered by a configurable score threshold (\texttt{RAG\_SCORE\_THRESHOLD}, default 0.30); scores are computed as \texttt{1.0 - distance}.

\reqid{Requirement FR-RAG-005:} When a tool hint is supplied, retrieval shall filter candidates to the hinted tool's chunks.

\reqid{Requirement FR-RAG-006:} Each returned chunk shall expose \texttt{content}, \texttt{tool}, and \texttt{score} for ground-truth traceability, and the originating document names shall be surfaced in \texttt{retrieval\_sources} of the \texttt{PlannerOutput}.

\reqid{Requirement FR-RAG-007:} The knowledge base shall be updatable at runtime via \texttt{POST /kb/update}, which chunks content (approx 80/20 word split) and upserts documents with ids \texttt{\{tool\}\_\allowbreak \{uuid\}\_\allowbreak chunk\_\allowbreak \{idx\}}.

\subsection{Feature: Tool Selection and Context Assembly}

\subsubsection{Description}

\texttt{ToolSelector} resolves the execution tool using a strict priority chain. \texttt{ContextAssembler} builds the grounded prompt from retrieved docs and intent parameters.

\reqid{Requirement FR-TOOL-001:} \texttt{ToolSelector.select} shall resolve the tool in the order: caller tool hint $\rightarrow$ top RAG retrieval source $\rightarrow$ \texttt{data/tool\_map.json} mapping $\rightarrow$ \texttt{"unknown"}.

\reqid{Requirement FR-TOOL-002:} The tool map shall map the intents \texttt{NETWORK\_SCAN} $\rightarrow$ \texttt{nmap}, \texttt{SERVICE\_ENUMERATION} $\rightarrow$ \texttt{nmap}, \texttt{VULNERABILITY\_AUDIT} $\rightarrow$ \texttt{nikto}, \texttt{DIRECTORY\_BRUTEFORCE} $\rightarrow$ \texttt{gobuster}; unhandled intents (\texttt{EXPLOITATION}, \texttt{PASSWORD\_ATTACK}, \texttt{PASSIVE\_RECON}) shall return no fallback tool and rely on RAG/hint resolution.

\reqid{Requirement FR-TOOL-003:} \texttt{SUPPORTED\_TOOLS} shall be \texttt{\{"nmap", "nikto", "gobuster"\}}.

\reqid{Requirement FR-ASSM-001:} \texttt{ContextAssembler.\allowbreak assemble} shall render \texttt{=== TOOL DOCUMENTATION ===} sections from retrieved docs, a \texttt{SUGGESTED COMMAND REFERENCE} block when a suggested command exists, and explicit \texttt{COMMAND:} versus \texttt{COMMAND SEQUENCE:} prompts for multi-step mode.

\reqid{Requirement FR-ASSM-002:} The assembly shall include a 4-item RULES block constraining the LLM output format (single command, no commentary, use only retrieved flag names, target must match the supplied value).

\reqid{Requirement FR-ASSM-003:} The assembled prompt shall reinforce the exact target value with text such as "IMPORTANT: Use this exact target... Output ONLY one command."

\subsection{Feature: Command Synthesis}

\subsubsection{Description}

\texttt{CommandSynthesizer} calls the local Ollama chat endpoint with the grounded prompt and parses the response into one or more commands.

\reqid{Requirement FR-SYN-001:} Synthesis shall invoke \texttt{ollama.chat()} via the configured \texttt{OLLAMA\_MODEL} (default \texttt{gemma4:e4b}) with a frozen system prompt describing it as a Kali Linux command generator for the NeuroShell C2 framework that outputs only valid, executable commands.

\reqid{Requirement FR-SYN-002:} Decoding parameters shall be deterministic-leaning: temperature 0.0, top\_p 0.9, top\_k 40, num\_predict 128 (as defined in \texttt{Modelfile}).

\reqid{Requirement FR-SYN-003:} \texttt{parse\_command\_output} shall strip markdown fences and number prefixes, filter out commentary lines, and extract candidate commands.

\reqid{Requirement FR-SYN-004:} Multi-step sequences shall be capped at 5 commands; each command in the sequence shall be individually validated.

\reqid{Requirement FR-SYN-005:} The output dict shall expose \texttt{command}, \texttt{command\_sequence}, \texttt{retrieval\_sources}, and \texttt{query\_used}.

\subsection{Feature: Safety Filter and Command Validation}

\subsubsection{Description}

Two independent validation layers gate every synthesized command before it may be released to the response.

\reqid{Requirement FR-SAFETY-001:} \texttt{SafetyFilter.\allowbreak validate(command, target)} shall reject commands containing any of 10 \texttt{BLOCKED\_PATTERNS} (e.g., \texttt{rm -rf}, \texttt{mkfs}, \texttt{dd if=...of=/dev/}, fork-bomb \texttt{:()\{ :|:\& \};:}, \texttt{wget|curl ...|bash}, \texttt{nc -e /bin*}, \texttt{bash -i >\&/dev/tcp}, \texttt{> /dev/sd*}, \texttt{chmod 777 /}).

\reqid{Requirement FR-SAFETY-002:} The safety filter shall reject targets outside the allowed scope: \texttt{192.168.*}, \texttt{10.*}, \texttt{172.(16--31).*}, \texttt{127.*}, and \texttt{localhost}.

\reqid{Requirement FR-SAFETY-003:} The safety filter shall allow-list 9 tools (\texttt{ALLOWED\_TOOLS}: nmap, nikto, gobuster, hydra, metasploit, whois, dig, curl, wget); use of other tools shall produce a non-fatal warning.

\reqid{Requirement FR-SAFETY-004:} Target normalization for scope checks shall use \texttt{urlparse}-based extraction (\texttt{\_normalize\_target}) with IP/hostname validation.

\reqid{Requirement FR-VALID-001:} \texttt{CommandValidator.\allowbreak validate(command, expected\_\allowbreak tool)} shall reject fatal conditions: empty command, multi-line commands, and a binary that does not match the expected tool's pattern.

\reqid{Requirement FR-VALID-002:} The validator shall emit warnings for missing tool-required flags (nmap: \texttt{-\allowbreak s, -\allowbreak p, -\allowbreak -\allowbreak top-\allowbreak ports, -\allowbreak A, -\allowbreak sn}; nikto: \texttt{-h}; gobuster: \texttt{-u, -w}) and for lack of an RFC-1918 IP in the command.

\reqid{Requirement FR-VALID-003:} Validated commands shall be returned only when both the safety filter and the validator pass; otherwise the endpoint returns the corresponding 400/422 error.

\subsection{Feature: Session Context Management}

\subsubsection{Description}

\texttt{SessionContextManager} persists per-session planning context in Redis with a 3600-second TTL, degrading to an in-memory fallback on Redis errors.

\reqid{Requirement FR-SES-001:} The session store shall persist \texttt{session\_\allowbreak id, targets\_\allowbreak seen[], commands\_\allowbreak run[], tools\_\allowbreak used[], created\_\allowbreak at, last\_\allowbreak updated} under key \texttt{session:\{session\_id\}}.

\reqid{Requirement FR-SES-002:} Session entries shall expire after 3600 seconds of inactivity (sliding TTL).

\reqid{Requirement FR-SES-003:} Latest target recovery shall be available via \texttt{get\_\allowbreak last\_\allowbreak target(session\_\allowbreak id)} and used when a request carries an empty/\texttt{UNKNOWN} target.

\reqid{Requirement FR-SES-004:} On \texttt{redis.RedisError}, the manager shall transparently fall back to an in-memory dictionary and log a warning; session updates shall still succeed.

\reqid{Requirement FR-SES-005:} An administrator shall be able to clear a session via \texttt{POST /\allowbreak admin/\allowbreak session/\allowbreak clear?session\_\allowbreak id=.\allowbreak .\allowbreak .\allowbreak } (admin auth required).

\subsection{Feature: Metrics and Administration}

\subsubsection{Description}

\texttt{MetricsCollector} records in-memory statistics; admin endpoints manage the knowledge base, sessions, metrics, and status.

\reqid{Requirement FR-MET-001:} Each plan request shall record \texttt{\{timestamp, session\_\allowbreak id, intent, tool, latency\_\allowbreak ms, success, safety\_\allowbreak flags\}} via \texttt{MetricsCollector.record}.

\reqid{Requirement FR-MET-002:} \texttt{GET /metrics} (API-key protected) shall return total, successful, failed, and average-latency statistics from \texttt{get\_summary()}.

\reqid{Requirement FR-MET-003:} \texttt{POST /\allowbreak admin/\allowbreak metrics/\allowbreak reset} shall clear the in-memory metrics store (admin auth required).

\reqid{Requirement FR-ADM-001:} \texttt{POST /kb/update} shall ingest content (max 51,200 characters) and return \texttt{\{status, chunks\_\allowbreak added, kb\_\allowbreak chunks\}} (admin auth required).

\reqid{Requirement FR-ADM-002:} \texttt{GET /admin/status} shall return service identity (\texttt{component:\allowbreak "C2-\allowbreak DynamicPlanner"}), admin authentication state, Redis health, KB chunk count, and metrics summary (admin auth required).

\reqid{Requirement FR-ADM-003:} Admin endpoints shall require \texttt{X-Admin-Key} matching \texttt{admin\_key} (fallback: \texttt{X-API-Key} carrying the admin value).

\subsection{Feature: Health and Tool Availability}

\reqid{Requirement FR-HLT-001:} \texttt{GET /health} shall report \texttt{status} (\texttt{healthy}|\texttt{degraded}/\texttt{ok}), \texttt{version}, \texttt{uptime}, Redis reachability, and KB chunk count.

\reqid{Requirement FR-HLT-002:} \texttt{GET /tools/available} (public) shall return the supported tool list and the current number of KB chunks.

\reqid{Requirement FR-HLT-003:} Degradation shall be reported when Redis is unreachable (in-memory fallback active) in the wrapper app health payload.

\section{External Interface Requirements}

\subsection{Software Interfaces}

\subsubsection{REST API (HTTP/JSON)}

\begin{table}[h]
\caption{Table 4.}
\label{tab:t4}
\begin{tabularx}{\columnwidth}{l >{\raggedright\arraybackslash}X >{\raggedright\arraybackslash}X >{\raggedright\arraybackslash}X}
\toprule
\textbf{Method} & \textbf{Path (wrapper)} & \textbf{Auth} & \textbf{Description}\\
\midrule
POST & \texttt{/api/v1/plan} & API key & Plan synthesis pipeline (core)\\
GET & \texttt{/api/v1/tools/available} & none & Tool list + KB chunk count\\
GET & \texttt{/api/v1/metrics} & API key & Metrics summary\\
POST & \texttt{/api/v1/kb/update} & Admin & Knowledge base ingestion\\
POST & \texttt{/\allowbreak api/\allowbreak v1/\allowbreak admin/\allowbreak session/\allowbreak clear} & Admin & Clear a session (query \texttt{session\_id})\\
POST & \texttt{/\allowbreak api/\allowbreak v1/\allowbreak admin/\allowbreak metrics/\allowbreak reset} & Admin & Reset metrics\\
GET & \texttt{/api/v1/admin/status} & Admin & Administration status\\
GET & \texttt{/health} & none & Health probe\\
\bottomrule
\end{tabularx}
\end{table}

\begin{quote}\small Standalone engine entry point (\texttt{src/api/main.py}) exposes the same endpoints with the same semantics \textbf{without} the \texttt{/api/v1} prefix (dev port 8002).\end{quote}

\subsubsection{POST /plan --- Request (Format A: enriched flat)}

\begin{lstlisting}
{
  "session_id": "s-001",
  "intent": "NETWORK_SCAN",
  "sub_intent": "full_port_scan",
  "confidence": 0.94,
  "human_readable_intent": "scan target",
  "target_type": "IP",
  "target_value": "192.168.1.1",
  "target_in_scope": true,
  "primary_tool": "nmap",
  "tool_parameters": {},
  "safety_notes": [],
  "target": "192.168.1.1",
  "ports": [22, 80, 443],
  "modifiers": ["fast"],
  "cve_ids": [],
  "tool_hint": "nmap",
  "schedule": null,
  "rejection_reason": null,
  "scope_warnings": []
}
\end{lstlisting}

\subsubsection{POST /plan --- Request (Format B: nested)}

\begin{lstlisting}
{
  "session_id": "s-001",
  "intent_contract": {
    "intent": "DIRECTORY_BRUTEFORCE",
    "sub_intent": null,
    "target_type": "IP",
    "target_value": "10.0.0.5",
    "primary_tool": "gobuster",
    "tool_parameters": {},
    "target": "10.0.0.5",
    "ports": [80],
    "modifiers": ["wordlist:common"],
    "cve_ids": [],
    "tool_hint": "gobuster",
    "schedule": null,
    "confidence": 0.91,
    "rejection_reason": null,
    "scope_warnings": []
  }
}
\end{lstlisting}

\subsubsection{POST /plan --- Response (200)}

\begin{lstlisting}
{
  "status": "success",
  "command": "nmap -sS -p 22,80,443 192.168.1.1",
  "command_sequence": [],
  "tool": "nmap",
  "session_id": "s-001",
  "intent_ref": "NETWORK_SCAN",
  "estimated_duration": "medium",
  "retrieval_sources": ["nmap_documentation.txt"],
  "validation_passed": true,
  "safety_flags": [],
  "latency_ms": 1850
}
\end{lstlisting}

\subsubsection{POST /plan --- Error responses}

\begin{table}[h]
\caption{Table 5.}
\label{tab:t5}
\begin{tabularx}{\columnwidth}{l >{\raggedright\arraybackslash}X >{\raggedright\arraybackslash}X}
\toprule
\textbf{HTTP} & \textbf{Body (example)} & \textbf{When}\\
\midrule
401 & \texttt{\{"detail":\allowbreak  "Invalid API key"\}} & key mismatch\\
400 & \texttt{\{"stage":\allowbreak "intent\_\allowbreak filter","error":\allowbreak "INTENT\_\allowbreak REJECTED"\}} & REJECTED contract\\
400 & \texttt{\{"detail":\allowbreak "Safety check failed:\allowbreak  <flags>"\}} & blocked pattern / public target\\
422 & \texttt{\{"status":\allowbreak "error","stage":\allowbreak "command\_\allowbreak validator","error":\allowbreak "VALIDATION\_\allowbreak FAILED","detail":\allowbreak ".\allowbreak .\allowbreak .\allowbreak ","latency\_\allowbreak ms":\allowbreak N\}} & structural validation failure\\
500 & \texttt{\{"detail":\allowbreak "<exception>"\}} & unhandled error\\
\bottomrule
\end{tabularx}
\end{table}

\subsection{Hardware Interfaces}

\begin{itemize}
\item No direct hardware interfaces. C2 is a stateless service; all heavy computation (embeddings, LLM inference) runs on the host CPU/GPU via local processes (SentenceTransformer in-process; Ollama out-of-process).
\end{itemize}

\subsection{Software Interfaces (outbound)}

\begin{table}[h]
\caption{Table 6.}
\label{tab:t6}
\begin{tabularx}{\columnwidth}{l >{\raggedright\arraybackslash}X >{\raggedright\arraybackslash}X}
\toprule
\textbf{Interface} & \textbf{Protocol} & \textbf{Purpose}\\
\midrule
Ollama & local HTTP chat (\texttt{ollama.chat()}) & LLM command synthesis\\
ChromaDB & embedded \texttt{PersistentClient} (local disk) & vector knowledge base\\
Redis & TCP (\texttt{redis.from\_url}) & session persistence (optional)\\
SentenceTransformer & in-process & \texttt{all-MiniLM-L6-v2} embeddings (384-dim)\\
\bottomrule
\end{tabularx}
\end{table}

\subsection{Communication Interfaces}

\begin{itemize}
\item \textbf{Transport}: HTTP/1.1 (optionally behind the Kong gateway with TLS termination at the gateway)
\item \textbf{Serialization}: JSON (Pydantic v2 models)
\item \textbf{Auth in transport}: \texttt{X-API-Key} / \texttt{X-Admin-Key} request headers for protected endpoints
\item \textbf{CORS}: wrapper app uses settings-derived origins; dev default \texttt{["*"]} with credentials enabled; standalone dev engine uses \texttt{allow\_\allowbreak credentials=False, allow\_\allowbreak origins=["*"]}
\end{itemize}

\section{Non-Functional Requirements}

\subsection{Performance}

\reqid{Requirement NFR-PERF-001:} End-to-end plan latency shall be measured via \texttt{MetricsCollector} and reported in \texttt{PlannerOutput.\allowbreak latency\_\allowbreak ms} and \texttt{/metrics}.

\reqid{Requirement NFR-PERF-002:} Retrieval (embedding + ChromaDB query, top\_k=3) shall complete in under 300 ms on typical hardware.

\reqid{Requirement NFR-PERF-003:} The dominant latency driver is LLM inference; the system shall tolerate first-call model warm-up (30--90 s) without failing the request but shall record the latency.

\reqid{Requirement NFR-PERF-004:} Target average end-to-end latency shall be under 30 s for single-command plans under the default model (observed 21.7 s avg Gemma4 / 8.4 s avg tuned Gemma2 in evaluation).

\subsection{Security}

\reqid{Requirement NFR-SEC-001:} All protected endpoints shall require a valid \texttt{X-API-Key}; admin endpoints shall additionally require \texttt{X-Admin-Key} (or API key with admin value).

\reqid{Requirement NFR-SEC-002:} The default development keys (\texttt{neuroshell-secret-key} / \texttt{neuroshell-admin-key}) shall never be used in production; production keys shall be supplied via environment/secret injection.

\reqid{Requirement NFR-SEC-003:} The service shall never execute synthesized commands; \texttt{BLOCKED\_PATTERNS} and \texttt{ALLOWED\_TOOLS} enforcement shall be guaranteed before a command leaves the API.

\reqid{Requirement NFR-SEC-004:} Target scope enforcement shall restrict planning to RFC-1918 private ranges and localhost; public targets shall be rejected.

\reqid{Requirement NFR-SEC-005:} Key comparisons shall be constant-time-safe in any hardened deployment (note: current implementation uses direct string comparison --- flagged for hardening).

\reqid{Requirement NFR-SEC-006:} CORS in production shall be restricted to configured origins (current dev default \texttt{["*"]} shall not be shipped to production unchanged).

\reqid{Requirement NFR-SEC-007:} No secrets shall be committed to version control; \texttt{.env} files and model artifacts (\texttt{models/}) are gitignored.

\reqid{Requirement NFR-SEC-008:} LLM prompt grounding shall prevent hallucinated flags by instructing the model to use only flag names present in the retrieved documentation (RULES block).

\subsection{Reliability}

\reqid{Requirement NFR-REL-001:} Redis unavailability shall degrade gracefully to in-memory session storage with a logged warning (no request failure).

\reqid{Requirement NFR-REL-002:} A rejected or failed plan shall still be recorded in metrics for observability (fail-open for metrics, fail-safe for command release).

\reqid{Requirement NFR-REL-003:} Delivery shall reuse a single ingestion pipeline so the knowledge base can be refreshed without service restart.

\reqid{Requirement NFR-REL-004:} The health endpoint shall expose dependency state (Redis, KB count) for orchestrator health checks (container healthcheck: \texttt{curl /health}, 30 s interval, 3 retries, 40 s start period).

\subsection{Maintainability}

\reqid{Requirement NFR-MNT-001:} All tunable parameters (model name, retrieval threshold, top\_k, step cap, keys, Redis URL) shall be environment-configurable without code changes.

\reqid{Requirement NFR-MNT-002:} The pipeline shall be implemented as independently testable modules (intent, rag, synthesis, validation, session, utils) with unit coverage.

\reqid{Requirement NFR-MNT-003:} The knowledge base content and tool map shall be data-driven (file-based: \texttt{knowledge\_\allowbreak base/\allowbreak raw\_\allowbreak docs/\allowbreak *.\allowbreak txt}, \texttt{data/tool\_map.json}) to enable content updates independent of code.

\reqid{Requirement NFR-MNT-004:} Documentation shall be maintained under \texttt{docs/} covering API, data models, pipeline flow, RAG subsystem, safety, sessions, metrics, configuration, and deployment.

\subsection{Portability}

\reqid{Requirement NFR-PORT-001:} The service shall run on Linux, macOS, and Windows (development) with Python 3.9+.

\reqid{Requirement NFR-PORT-002:} Docker deployment shall use the \texttt{python:3.11-slim} base image.

\reqid{Requirement NFR-PORT-003:} All configuration shall be environment-variable-driven via \texttt{.env} or Docker environment.

\reqid{Requirement NFR-PORT-004:} The ChromaDB knowledge base directory and session store shall be persistent volumes in containerized deployments.

\section{Data Requirements}

\subsection{Data Models (src/schemas/models.py)}

\begin{table}[h]
\caption{Table 7.}
\label{tab:t7}
\begin{tabularx}{\columnwidth}{l >{\raggedright\arraybackslash}X >{\raggedright\arraybackslash}X}
\toprule
\textbf{Model} & \textbf{Key fields} & \textbf{Notes}\\
\midrule
\texttt{ToolParameters} & \texttt{target}, \texttt{flags[]}, \texttt{ports[]}, \texttt{wordlist}, \texttt{suggested\_command} & tool parameter bag\\
\texttt{IREIntentContract} & \texttt{intent} (9-literal enum), \texttt{sub\_intent}, \texttt{target\_type}, \texttt{target\_value}, \texttt{primary\_tool}, \texttt{tool\_parameters}, \texttt{target}, \texttt{ports}, \texttt{modifiers}, \texttt{cve\_ids}, \texttt{tool\_hint}, \texttt{schedule}, \texttt{confidence} (float), \texttt{rejection\_reason}, \texttt{scope\_warnings} & C1$\rightarrow$C2 nested contract\\
\texttt{PlanRequest} & \texttt{intent\_contract?}, \texttt{session\_id} + 20 property shims; \texttt{extra="allow"} & inbound request\\
\texttt{PlannerOutput} & \texttt{status}, \texttt{command}, \texttt{command\_sequence[]}, \texttt{tool}, \texttt{session\_id}, \texttt{intent\_ref}, \texttt{estimated\_\allowbreak duration="medium"}, \texttt{retrieval\_sources[]}, \texttt{validation\_passed}, \texttt{safety\_flags[]}, \texttt{latency\_ms} & outbound plan\\
\texttt{ErrorResponse} & \texttt{status}, \texttt{error}, \texttt{stage}, \texttt{detail}, \texttt{latency\_ms} & error envelope\\
\texttt{KBUpdateRequest} & \texttt{content} ($\leq$51,200 chars), \texttt{metadata} (dict) & KB ingestion\\
\bottomrule
\end{tabularx}
\end{table}

\subsection{Intent Enumeration}

The 9 intent literals shared with C1: \texttt{NETWORK\_SCAN}, \texttt{SERVICE\_ENUMERATION}, \texttt{VULNERABILITY\_AUDIT}, \texttt{DIRECTORY\_BRUTEFORCE}, \texttt{EXPLOITATION}, \texttt{PASSWORD\_ATTACK}, \texttt{PASSIVE\_RECON}, plus \texttt{AMBIGUOUS} / \texttt{REJECTED} handling (REJECTED short-circuits).

\subsection{Knowledge Base (ChromaDB)}

\begin{itemize}
\item Collection: \texttt{tool\_documentation}; distance space: cosine (\texttt{hnsw:space})
\item Source documents: \texttt{knowledge\_\allowbreak base/\allowbreak raw\_\allowbreak docs/\allowbreak \{nmap,nikto,gobuster\}.\allowbreak txt} (101/98/91 lines)
\item Embedding: \texttt{all-MiniLM-L6-v2}, 384 dimensions
\item Runtime: \textasciitilde\{\}30 chunks live; chunk scheme \textasciitilde\{\}80/20 words; doc ids \texttt{\{tool\}\_chunk\_\{n\}} / \texttt{\{tool\}\_\allowbreak \{uuid\}\_\allowbreak chunk\_\allowbreak \{idx\}}
\item Metadata per chunk: \texttt{\{tool, chunk\_index\}}
\end{itemize}

\subsection{Session Store (Redis)}

\begin{itemize}
\item Key: \texttt{session:\{session\_id\}}; TTL: 3600 s sliding
\item Fields: \texttt{session\_id}, \texttt{targets\_seen[]}, \texttt{commands\_run[]}, \texttt{tools\_used[]}, \texttt{created\_at}, \texttt{last\_updated}
\item Fallback: in-memory dict on Redis errors
\end{itemize}

\subsection{Metrics Store}

\begin{itemize}
\item In-memory list of \texttt{\{timestamp, session\_\allowbreak id, intent, tool, latency\_\allowbreak ms, success, safety\_\allowbreak flags\}}
\item Ephemeral: reset on restart or via admin reset
\end{itemize}

\subsection{Tool Map (data/tool\_map.json)}

Maps intent $\rightarrow$ fallback tool for \texttt{NETWORK\_SCAN}$\rightarrow$nmap, \texttt{SERVICE\_ENUMERATION}$\rightarrow$nmap, \texttt{VULNERABILITY\_AUDIT}$\rightarrow$nikto, \texttt{DIRECTORY\_BRUTEFORCE}$\rightarrow$gobuster.

\section{Verification and Validation}

\subsection{Test Suite}

\begin{table}[h]
\caption{Table 8.}
\label{tab:t8}
\begin{tabularx}{\columnwidth}{l >{\raggedright\arraybackslash}X >{\raggedright\arraybackslash}X}
\toprule
\textbf{Area} & \textbf{File} & \textbf{Coverage}\\
\midrule
Intent interpretation & \texttt{tests/\allowbreak unit/\allowbreak test\_\allowbreak intent\_\allowbreak interpreter.\allowbreak py} (4) & enriched/legacy format, default ports, multi-step\\
Query construction & \texttt{tests/\allowbreak unit/\allowbreak test\_\allowbreak query\_\allowbreak constructor.\allowbreak py} (3) & templates, modifiers, whitespace\\
Vector retrieval & \texttt{tests/\allowbreak unit/\allowbreak test\_\allowbreak vector\_\allowbreak retriever.\allowbreak py} (3) & init, tool-filtered retrieve, threshold\\
Session context & \texttt{tests/\allowbreak unit/\allowbreak test\_\allowbreak session\_\allowbreak context\_\allowbreak manager.\allowbreak py} (1) & target recovery / UNKNOWN / clear\\
Safety filter & \texttt{tests/\allowbreak unit/\allowbreak test\_\allowbreak safety\_\allowbreak filter.\allowbreak py} (3) & private pass, dangerous block, public block\\
Command validator & \texttt{tests/\allowbreak unit/\allowbreak test\_\allowbreak command\_\allowbreak validator.\allowbreak py} (3) & valid, mismatched tool, empty/multiline\\
Admin endpoints & \texttt{tests/\allowbreak unit/\allowbreak test\_\allowbreak admin.\allowbreak py} (7) & admin auth, \texttt{/kb/update}, admin ops\\
Multi-command plans & \texttt{tests/\allowbreak unit/\allowbreak test\_\allowbreak multi\_\allowbreak command.\allowbreak py} (6) & suggested-command prompt, multi-step, output parser, sequence safety/validator\\
API integration & \texttt{tests/\allowbreak integration/\allowbreak test\_\allowbreak plan\_\allowbreak api.\allowbreak py} (6) & auth fail, REJECTED 400, public-IP 400, validator 422, session recovery, health/tools/metrics\\
\bottomrule
\end{tabularx}
\end{table}

\reqid{Requirement NFR-TST-001:} The verified suite shall contain 36 test functions (30 unit, 6 integration) as currently present in the repository.

\reqid{Requirement NFR-TST-002:} Note: README/docs currently describe "51/51 tests passing"; the SRS shall flag this discrepancy and treat the 36 verified tests as the source of truth until the docs are reconciled.

\subsection{Empirical Evaluation}

\begin{itemize}
\item \textbf{Benchmark}: 25-case evaluation campaign (docs/15\_evaluation\_report.md); average 21.7 s (Gemma4) vs 8.4 s (tuned Gemma2)
\item \textbf{Success criteria}: RAG retrieval returns relevant tool docs (score threshold 0.30); synthesized commands pass safety + structural validation; multi-step sequences capped at 5 commands
\end{itemize}

\subsection{Validation Checks on Delivery}

\reqid{Requirement FR-VER-001:} Delivery shall pass \texttt{pytest tests/\allowbreak unit tests/\allowbreak integration}.

\reqid{Requirement FR-VER-002:} The health endpoint shall report healthy and the \texttt{/tools/available} endpoint shall report 3 tools and current KB chunk count.

\reqid{Requirement FR-VER-003:} A smoke plan against \texttt{192.168.x.x} target with a mocked synthesizer shall produce a \texttt{validation\_passed=true} plan; a public-IP target shall be rejected with HTTP 400.

\section{Deployment}

\subsection{Prerequisites}

\begin{itemize}
\item Docker Engine + docker-compose (deployment) or Python 3.9+ venv (development)
\item Ollama running locally with the configured model (\texttt{gemma4:e4b}) --- tuned deployment additionally requires \texttt{.\allowbreak /\allowbreak models/\allowbreak gemma2-\allowbreak c2-\allowbreak lora.\allowbreak gguf} (external, not committed) to create \texttt{neuroshell-\allowbreak c2-\allowbreak gemma2:\allowbreak latest}
\item Redis reachable at \texttt{REDIS\_URL} (optional; in-memory fallback applies otherwise)
\item Seeded ChromaDB KB (run \texttt{scripts/ingest\_docs.py} at build/first boot)
\end{itemize}

\subsection{Env Configuration}

\begin{table}[h]
\caption{Table 9.}
\label{tab:t9}
\begin{tabularx}{\columnwidth}{l >{\raggedright\arraybackslash}X >{\raggedright\arraybackslash}X}
\toprule
\textbf{Variable} & \textbf{Default} & \textbf{Purpose}\\
\midrule
\texttt{API\_KEY} & \texttt{neuroshell-secret-key} & plan/metrics auth (replace in production)\\
\texttt{ADMIN\_KEY} & \texttt{neuroshell-admin-key} & admin auth (replace in production)\\
\texttt{REDIS\_URL} & \texttt{redis://localhost:6379} & session store\\
\texttt{OLLAMA\_MODEL} & \texttt{gemma4:e4b} & synthesis model\\
\texttt{RAG\_SCORE\_THRESHOLD} & \texttt{0.30} & retrieval similarity floor\\
\texttt{PORT} & \texttt{8000} & uvicorn port (root main.py)\\
\texttt{RAG\_TOP\_K} / \texttt{MAX\_PLAN\_STEPS} / \texttt{RAG\_RETRIEVAL\_LIMIT} & \texttt{3} / \texttt{10} / \texttt{5} & pipeline limits (settings)\\
\bottomrule
\end{tabularx}
\end{table}

\subsection{Build and Start (Docker)}

\begin{lstlisting}
# seed KB (optional refresh)
python scripts/ingest_docs.py

# build and start
docker build -t neuroshell-dynamic-planner-rag .
docker run -p 8003:8000 -e OLLAMA_MODEL=neuroshell-c2-gemma2:latest neuroshell-dynamic-planner-rag
\end{lstlisting}

\subsection{Verify}

\begin{lstlisting}
curl http://localhost:8003/health
curl "http://localhost:8003/api/v1/tools/available"
curl -X POST http://localhost:8003/api/v1/plan -H "Content-Type: application/json" -H "X-API-Key: <API_KEY>" -d @plan.json
\end{lstlisting}

\subsection{Known Deployment Gaps}

\begin{itemize}
\item Production/staging \texttt{.env} files at monorepo \texttt{configs/\allowbreak \{dev,prod,staging\}/\allowbreak } contain placeholder secrets and must be replaced with real injected secrets
\item C2 tests are not yet wired into CI (\texttt{.github/workflows/ci.yml} builds the image but does not run C2 tests) --- flagged for hardening
\item No TLS termination in process; TLS is expected at the reverse proxy (Kong)
\item \texttt{knowledge\_\allowbreak base/\allowbreak chroma\_\allowbreak db} should be a persistent volume (dev compose mounts the whole directory)
\end{itemize}

\section{Appendices}

\subsection{Appendix A --- Intent Enumeration and Tool Map}

\begin{table}[h]
\caption{Table 10.}
\label{tab:t10}
\begin{tabularx}{\columnwidth}{l >{\raggedright\arraybackslash}X >{\raggedright\arraybackslash}X >{\raggedright\arraybackslash}X}
\toprule
\textbf{Intent literal} & \textbf{Primary tool (tool\_map)} & \textbf{KB source doc} & \textbf{Notes}\\
\midrule
\texttt{NETWORK\_SCAN} & nmap & nmap.txt & full port / NSE\\
\texttt{SERVICE\_ENUMERATION} & nmap & nmap.txt & distinct query template\\
\texttt{VULNERABILITY\_AUDIT} & nikto & nikto.txt & CVE-ID query part\\
\texttt{DIRECTORY\_BRUTEFORCE} & gobuster & gobuster.txt & wordlist modifier expansion\\
\texttt{EXPLOITATION} & (none) & --- & multi-step; RAG/hint resolution\\
\texttt{PASSWORD\_ATTACK} & (none) & --- & not in KB raw docs\\
\texttt{PASSIVE\_RECON} & (none) & --- & enumeration-style query\\
\bottomrule
\end{tabularx}
\end{table}

\subsection{Appendix B --- Blocked Patterns (SafetyFilter)}

\texttt{rm -rf}; \texttt{mkfs}; \texttt{dd if=...of=/dev/}; fork bomb \texttt{:()\{ :|:\& \};:}; \texttt{chmod 777 /}; \texttt{> /dev/sd*}; \texttt{wget|curl ...|bash}; \texttt{nc -e /bin*}; \texttt{bash -i >\&/dev/tcp}; and additional dangerous shell constructs (10 \texttt{BLOCKED\_PATTERNS} regexes).

\subsection{Appendix C --- Modifier Text Map (16 entries)}

\texttt{stealth}, \texttt{aggressive}, \texttt{all-ports}, \texttt{quick}, \texttt{os-detection}, \texttt{ssl}, \texttt{silent}, \texttt{full}, \texttt{evasion}, \texttt{extension:php}, \texttt{extension:html}, \texttt{dns}, \texttt{vhost}, \texttt{wordlist:common}, \texttt{wordlist:medium}, plus remaining map entries.

\subsection{Appendix D --- Query Templates (7 intents + default)}

Intent-specific templates for \texttt{NETWORK\_SCAN}, \texttt{SERVICE\_ENUMERATION}, \texttt{VULNERABILITY\_AUDIT}, \texttt{DIRECTORY\_BRUTEFORCE}, \texttt{EXPLOITATION}, \texttt{PASSWORD\_ATTACK}, \texttt{PASSIVE\_RECON}, plus a generic default template.

\subsection{Appendix E --- PlannerOutput Schema (stable hand-off contract for C3)}

\begin{table}[h]
\caption{Table 11.}
\label{tab:t11}
\begin{tabularx}{\columnwidth}{l >{\raggedright\arraybackslash}X >{\raggedright\arraybackslash}X}
\toprule
\textbf{Field} & \textbf{Type} & \textbf{Notes}\\
\midrule
\texttt{status} & str & \texttt{success}\\
\texttt{command} & str & single primary command\\
\texttt{command\_sequence} & list[str] & multi-step ($\leq$5)\\
\texttt{tool} & str & resolved tool\\
\texttt{session\_id} & str & passthrough\\
\texttt{intent\_ref} & str & intent literal\\
\texttt{estimated\_duration} & str & \texttt{medium}\\
\texttt{retrieval\_sources} & list[str] & grounding provenance\\
\texttt{validation\_passed} & bool & true\\
\texttt{safety\_flags} & list[str] & warnings\\
\texttt{latency\_ms} & int & recorded E2E latency\\
\bottomrule
\end{tabularx}
\end{table}

C3 consumes the stable fields (\texttt{command}, \texttt{tool}, \texttt{session\_id}, \texttt{intent\_ref}, \texttt{estimated\_duration}) defined in the C3 \texttt{PlannerOutput} input schema.


\srsdivider{Component 03 --- Adaptive Execution \& Error Recovery Engine (AEERE)}
\section{Introduction}

\subsection{Purpose}

This Software Requirements Specification (SRS) defines the complete functional and non-functional requirements for the \textbf{NeuroShell Adaptive Execution \& Error Recovery Engine (AEERE)}, designated Component 03 of the NeuroShell Framework. AEERE is an autonomous self-healing execution engine that receives tool-grounded command plans from the Component 02 Dynamic Planner, executes them inside an isolated Kali Linux Docker container, detects and classifies failures (including silent failures), and iteratively corrects and retries the command within strict safety and budget constraints.

This document is intended for:

\begin{itemize}
\item \textbf{Research evaluators} assessing the AEERE design and empirical recovery performance
\item \textbf{Developers} maintaining or extending the execution, classification, mutation, and reporting modules
\item \textbf{Upstream / downstream consumers} (Component 02 planner feeding plans; Component 04 AVAE consuming execution results) integrating with the AEERE API
\item \textbf{Examiners} reviewing the MSc Cybersecurity research project (ref R26-CS-018)
\end{itemize}

\subsection{Scope}

\subsubsection{Product Scope}

The AEERE is a \textbf{self-contained FastAPI microservice} that:

\begin{enumerate}
\item Accepts a \texttt{PlannerOutput} plan (command, tool, session\_id, intent\_ref) from Component 02 over REST
\item Executes the command inside a fresh, disposable Docker Kali Linux container with strict resource limits
\item Detects execution success, hard failure, and \textbf{silent failure} (exit 0 with no useful output)
\item Classifies errors through a 3-tier cascade: regex patterns $\rightarrow$ embedding similarity $\rightarrow$ LLM causal reasoning
\item Produces a corrected command via a priority mutation chain: dataset lookup $\rightarrow$ LLM generation $\rightarrow$ rule-based transformation
\item Validates every corrected command through a 4-layer safety gate (emptiness, danger patterns, confidence, RFC-1918 scope)
\item Retries execution with budgeted attempts and backoff, always honouring a hard wall-clock deadline
\item Reports a structured \texttt{ExecutionResult} with a full recovery log, correction history, and failure report
\end{enumerate}

\subsubsection{In Scope}

\begin{itemize}
\item All source code under \texttt{app/src/} (API, execution, capture, classification, mutation, safety, retry, reporting, session)
\item Error taxonomy and tool substitution data under \texttt{app/taxonomy/}
\item Kali error dataset and loader under \texttt{app/kali\_data/}
\item Audit persistence (\texttt{app/data/audit.db} SQLite)
\item Test suite under \texttt{app/tests/} (unit + engine/integration)
\item Benchmark suites (\texttt{app/benchmark\_*.py}) and safety/recovery evaluation scripts
\item Deployment artifacts (Dockerfile, Dockerfile.kali, docker-compose.yml, startup.ps1)
\end{itemize}

\subsubsection{Out of Scope}

\begin{itemize}
\item Component 01 (IRE), Component 02 (Dynamic Planner), Component 04 (AVAE) --- referenced only for interface design
\item Command \textbf{synthesis} (planning) --- AEERE executes and repairs commands; it does not originate them
\item Vulnerability analysis / exploitation results interpretation --- forwarded to C4
\item Network infrastructure, host-level isolation, and physical security
\item Long-running listener/interactive sessions (explicitly rejected by the mutation engine)
\end{itemize}

\subsection{Definitions, Acronyms, Abbreviations}

\begin{table}[h]
\caption{Table 1.}
\label{tab:t1}
\begin{tabularx}{\columnwidth}{l >{\raggedright\arraybackslash}X}
\toprule
\textbf{Term} & \textbf{Definition}\\
\midrule
C1/C2/C3/C4 & Component 01 (IRE) / 02 (Planner) / 03 (AEERE) / 04 (AVAE)\\
AEERE & Adaptive Execution \& Error Recovery Engine\\
PlannerOutput & C2$\rightarrow$C3 inbound plan (command, tool, session\_id, intent\_ref, estimated\_duration)\\
ExecutionResult & C3 outbound structured result\\
RawExecutionResult & Container execution output (command, stdout, stderr, exit\_code, timed\_out)\\
ErrorContext & Sanitized context passed to classifiers\\
RecoveryPlan & Classified error, root cause, strategy, corrected command, confidence\\
Silent failure & Exit code 0 but no useful output / stderr error signal\\
RFC-1918 & Address Allocation for Private Internets\\
KALI\_PACKAGE\_WHITELIST & Tools allowed for automatic \texttt{apt-get install}\\
Tool substitution & Replacing a tool with a mapped equivalent (gobuster$\rightarrow$ffuf, nmap$\rightarrow$masscan, etc.)\\
Retry budget & Per-error-class max attempts and backoff from \texttt{error\_taxonomy.json}\\
Correction source & Origin of a corrected command: \texttt{rule-based}, \texttt{dataset}, or \texttt{llm}\\
Safety gate & 4-layer validation: emptiness, danger patterns, confidence, scope\\
\bottomrule
\end{tabularx}
\end{table}

\subsection{References}

\begin{table}[h]
\caption{Table 2.}
\label{tab:t2}
\begin{tabularx}{\columnwidth}{l >{\raggedright\arraybackslash}X >{\raggedright\arraybackslash}X}
\toprule
\textbf{Ref ID} & \textbf{Document} & \textbf{Location}\\
\midrule
R1 & Component 03 README / architecture & \texttt{README.md}, module docstrings\\
R2 & C2 Dynamic Planner SRS (upstream plans) & \texttt{dynamic-\allowbreak planner-\allowbreak rag/\allowbreak SRS-\allowbreak Doc/\allowbreak }\\
R3 & C4 AVAE SRS (downstream consumer of results) & \texttt{ai-\allowbreak vulnerability-\allowbreak analysis/\allowbreak SRS-\allowbreak Doc/\allowbreak }\\
R4 & Error taxonomy & \texttt{app/\allowbreak taxonomy/\allowbreak error\_\allowbreak taxonomy.\allowbreak json}\\
R5 & Tool substitution map & \texttt{app/\allowbreak taxonomy/\allowbreak tool\_\allowbreak substitution\_\allowbreak map.\allowbreak json}\\
R6 & Kali error dataset (labeled corrections) & \texttt{app/\allowbreak kali\_\allowbreak data/\allowbreak kali\_\allowbreak linux\_\allowbreak error.\allowbreak csv}\\
R7 & Benchmark results & \texttt{app/benchmark\_*.py} + \texttt{.results.json}\\
R8 & IEEE 830-1998 SRS Standard & IEEE Computer Society\\
R9 & RFC 1918 --- Private Address Space & IETF\\
R10 & OWASP Top 10 (2021) & owasp.org\\
\bottomrule
\end{tabularx}
\end{table}

\subsection{Overview}

The remainder of this document is organized as follows:

\begin{itemize}
\item \textbf{Section 2} describes the overall product perspective, user classes, and operating environment
\item \textbf{Section 3} specifies all system features and functional requirements in detail
\item \textbf{Section 4} defines external interface requirements (API, hardware, software, communication)
\item \textbf{Section 5} specifies non-functional requirements (performance, security, reliability)
\item \textbf{Section 6} covers data requirements (models, taxonomy, dataset, audit store)
\item \textbf{Section 7} describes verification, validation, and empirical results
\item \textbf{Section 8} covers deployment procedures
\item \textbf{Section 9} provides appendices with enumerations, strategies, and schemas
\end{itemize}

\section{Overall Description}

\subsection{Product Perspective}

The AEERE is \textbf{Component 03} of the NeuroShell Framework. It closes the loop between LLM-generated plans and real offensive-security tooling by executing commands in a disposable sandbox and recovering from failures autonomously.

\begin{lstlisting}
C1 IRE ──intent contract──▶ C2 Planner ──PlannerOutput──▶ C3 AEERE ──ExecutionResult──▶ C4 AVAE
                         (synthesis)         │ (cash)         (this component)
                                              │
                                              ▼
                                 macOS: Docker Kali container
                                 (isolated, resource-bounded, disposable)
\end{lstlisting}

\subsection{Product Functions Summary}

\begin{itemize}
\item \textbf{Sandboxed command execution} --- Docker Kali container per attempt, bridge network, privileged tooling, 512 m memory, 50,000 CPU-quota, 60 s exec timeout
\item \textbf{Silent failure detection} --- flags exit-0 runs that produce no output for output-requiring tools, with stderr error-pattern override and version/help exemption
\item \textbf{3-tier error classification} --- deterministic regex $\rightarrow$ embedding cosine similarity $\rightarrow$ LLM causal reasoning
\item \textbf{Prompt injection defense} --- redacts known injection patterns from stdout/stderr before LLM exposure
\item \textbf{Mutation engine} --- strategy-routed correction: \texttt{INSTALL\_TOOL}, \texttt{ADD\_SUDO}, \texttt{ADJUST\_PARAMETERS}, \texttt{TOOL\_SUBSTITUTION}, \texttt{FIX\_SYNTAX}, \texttt{ADJUST\_TIMEOUT}, \texttt{FIX\_AUTH}, \texttt{VERSION\_DOWNGRADE}, \texttt{ESCALATE}
\item \textbf{Target preservation} --- dataset/LLM corrections are re-targeted to the user's original host so a fix never silently redirects a scan
\item \textbf{4-layer safety gate} --- emptiness, danger patterns (16 regexes), confidence threshold (0.60), RFC-1918/domain-scope enforcement
\item \textbf{Budgeted retry with backoff} --- per-error-class max retries × fixed/exponential backoff, bounded by a hard wall-clock deadline (default 240 s)
\item \textbf{Audit \& reporting} --- SQLite append-only audit trail, structured failure reports with recommendations, per-session query
\item \textbf{Session state} --- Redis (optional) tracking of prior recovery strategies per session
\end{itemize}

\subsection{User Classes and Characteristics}

\begin{table}[h]
\caption{Table 3.}
\label{tab:t3}
\begin{tabularx}{\columnwidth}{l >{\raggedright\arraybackslash}X >{\raggedright\arraybackslash}X}
\toprule
\textbf{User class} & \textbf{Description} & \textbf{Interactions}\\
\midrule
\textbf{End users / operator} & Human operator via frontend or direct REST & \texttt{POST /execute}, \texttt{GET /sessions/\{id\}}\\
\textbf{Upstream C2 (Planner)} & Submits \texttt{PlannerOutput} plans directly & \texttt{POST /execute}\\
\textbf{Downstream C4 (AVAE)} & Consumes \texttt{ExecutionResult} for analysis & (outbound hand-off, C4 side)\\
\textbf{Administrator / analyst} & Reviews audit trails and failure reports & \texttt{GET /sessions/\{id\}}, \texttt{GET /health}\\
\textbf{Maintainer / researcher} & Extends taxonomy, dataset, classifiers & source, \texttt{app/kali\_data/}, benchmark scripts\\
\bottomrule
\end{tabularx}
\end{table}

\subsection{Operating Environment}

\begin{itemize}
\item \textbf{Runtime}: Python 3.11+, FastAPI + Uvicorn ASGI
\item \textbf{Execution sandbox}: Docker Engine with \texttt{kalilinux/kali-rolling} image (or \texttt{neuroshell-kali} build)
\item \textbf{State stores}: SQLite (\texttt{app/data/audit.db}); Redis (optional session state)
\item \textbf{LLM}: local Ollama server --- default model \texttt{qwen2.5-coder:latest} (health reports \texttt{OLLAMA\_MODEL})
\item \textbf{Embedding}: sentence-transformers model for embedding classifier (graceful skip if unavailable)
\item \textbf{Ports}: default 8000 (uvicorn, from \texttt{PORT} env); compose \texttt{6501:8000} in standalone compose
\item \textbf{OS}: Linux (production); Windows development via \texttt{startup.ps1}; macOS support documented via Docker
\end{itemize}

\subsection{Design and Implementation Constraints}

\begin{itemize}
\item \textbf{Deadline before upstream timeouts}: recovery loop must return a proper \texttt{ExecutionResult} before the upstream chain client (\textasciitilde\{\}300 s) times out; default \texttt{AEERE\_\allowbreak DEADLINE\_\allowbreak SECONDS=240}
\item \textbf{Deterministic classification first}: regex path is preferred; LLM is the last-resort classifier to control cost and latency
\item \textbf{Safety over automation}: any candidate correction failing the safety gate stops the loop
\item \textbf{Root inside container}: Docker containers run as root; \texttt{sudo} stripping is rule-based (\texttt{ADD\_SUDO} removes \texttt{sudo})
\item \textbf{Target fidelity}: corrections are rejected if the user's original target cannot be preserved
\item \textbf{Package whitelist}: only \texttt{KALI\_PACKAGE\_WHITELIST} tools may be auto-installed (rustscan excluded --- not in Kali apt repos)
\end{itemize}

\subsection{Assumptions and Dependencies}

\begin{itemize}
\item Docker Engine is reachable by the service (\texttt{CommandRunner.\allowbreak check\_\allowbreak docker\_\allowbreak available})
\item \texttt{kalilinux/kali-rolling} image (or \texttt{neuroshell-kali}) is present; wordlist paths in \texttt{VALID\_WORDLIST\_PATHS} exist in the image
\item Ollama is reachable and hosts \texttt{OLLAMA\_MODEL} for LLM classification/correction (degradation if absent)
\item \texttt{app/\allowbreak kali\_\allowbreak data/\allowbreak kali\_\allowbreak linux\_\allowbreak error.\allowbreak csv} is present for dataset corrections (dataset load is optional/graceful)
\item Redis is optional for session state (in-memory degradation on error)
\end{itemize}

\section{System Features and Functional Requirements}

\subsection{Feature: Sandboxed Command Execution}

\subsubsection{Description}

\texttt{CommandRunner} executes each command in a disposable Docker Kali container with enforced resource limits and captures stdout, stderr, exit code, and timeout state.

\reqid{Requirement FR-EXEC-001:} \texttt{CommandRunner.\allowbreak execute(command)} shall run the command as \texttt{bash -c <command>} inside a fresh container from \texttt{KALI\_IMAGE} (default \texttt{kalilinux/kali-rolling}), detached, on the bridge network, with \texttt{DEBIAN\_\allowbreak FRONTEND=noninteractive}.

\reqid{Requirement FR-EXEC-002:} Resource limits shall be enforced per container: memory \texttt{CONTAINER\_MEM\_LIMIT} (default \texttt{512m}), CPU quota \texttt{CONTAINER\_CPU\_QUOTA} (default \texttt{50000}).

\reqid{Requirement FR-EXEC-003:} The exec timeout shall be \texttt{CONTAINER\_EXEC\_TIMEOUT} (default 60 s); on expiry the container shall be stopped and \texttt{RawExecutionResult.\allowbreak timed\_\allowbreak out=True} with exit code 124.

\reqid{Requirement FR-EXEC-004:} The runner shall capture and return \texttt{RawExecutionResult\{command, stdout, stderr, exit\_\allowbreak code, timed\_\allowbreak out\}}.

\reqid{Requirement FR-EXEC-005:} Containers shall be force-removed in a \texttt{finally} block regardless of outcome.

\reqid{Requirement FR-EXEC-006:} \texttt{check\_\allowbreak docker\_\allowbreak available()} shall ping the Docker daemon and report connectivity for \texttt{/health}.

\subsection{Feature: Silent Failure Detection}

\subsubsection{Description}

\texttt{is\_silent\_failure} identifies executions that return exit 0 but were nevertheless unsuccessful, so the recovery loop treats them as failures.

\reqid{Requirement FR-SIL-001:} For tools in \texttt{TOOLS\_\allowbreak REQUIRING\_\allowbreak OUTPUT = \{nmap, nikto, ffuf, dirb, hydra, sqlmap, whatweb, masscan\}}, an exit-0 run with stdout shorter than 10 characters shall be flagged as a silent failure.

\reqid{Requirement FR-SIL-002:} If stderr contains any error pattern (\texttt{unknown option}, \texttt{invalid option}, \texttt{unrecognized option}, \texttt{requires an argument}, \texttt{requires a value}, \texttt{unknown argument}), the run shall be flagged as a silent failure regardless of exit code 0.

\reqid{Requirement FR-SIL-003:} Commands containing version/help keywords (\texttt{--version}, \texttt{-version}, \texttt{version}, \texttt{--help}, \texttt{-help}, \texttt{--h}) shall be exempted from the short-stdout check (but the stderr override shall still apply).

\subsection{Feature: Error Capture and Context Building}

\subsubsection{Description}

\texttt{build\_error\_context} constructs a sanitized \texttt{ErrorContext} from the raw execution result for downstream classifiers and mutators.

\reqid{Requirement FR-CAP-001:} The context shall contain \texttt{session\_\allowbreak id, command, stdout, stderr, exit\_\allowbreak code, tool, intent\_\allowbreak ref, attempt\_\allowbreak number, prior\_\allowbreak strategies, silent\_\allowbreak failure}.

\reqid{Requirement FR-CAP-002:} Prompt injection defense (\texttt{\_strip\_injections}) shall redact known injection patterns from stdout/stderr before they are passed to the LLM --- including "ignore previous instructions", "disregard the above", "forget previous instructions", and chat-template tokens such as \texttt{<|im\_start|>}, \texttt{<|im\_end|>}.

\reqid{Requirement FR-CAP-003:} The tool name shall be extracted from the command string (\texttt{\_extract\_tool\_name}) for classifier enrichment.

\subsection{Feature: 3-Tier Error Classification}

\subsubsection{Description}

The recovery loop classifies errors using a cascade that prioritizes cheap deterministic methods and reserves LLM reasoning as the last resort. Confidence thresholds and \texttt{error\_taxonomy.json} pattern sets drive the cascade.

\reqid{Requirement FR-CLS-001:} A container timeout shall be deterministically classified as \texttt{TIMEOUT} with confidence 1.0, without classifier invocation.

\reqid{Requirement FR-CLS-002:} \texttt{RegexClassifier.\allowbreak classify(stderr, exit\_\allowbreak code)} shall match stderr/exit-code signals against \texttt{error\_taxonomy.json} patterns and return \texttt{(error\_\allowbreak class, confidence)} or \texttt{(None, 0)}.

\reqid{Requirement FR-CLS-003:} When regex classification returns \texttt{None}, \texttt{EmbeddingClassifier.\allowbreak classify(stderr)} shall compute cosine similarity between the stderr text and labeled scenario embeddings; an unavailable embedding model shall be treated as a graceful skip.

\reqid{Requirement FR-CLS-004:} When both deterministic classifiers return \texttt{None}, \texttt{call\_llm\_reasoner(ctx)} shall invoke the LLM to produce a full \texttt{RecoveryPlan} (error class, root cause, strategy, corrected command, confidence, reasoning).

\reqid{Requirement FR-CLS-005:} All eight error classes from the taxonomy shall be supported: \texttt{TOOL\_NOT\_INSTALLED}, \texttt{PERMISSION\_DENIED}, \texttt{NETWORK\_UNREACHABLE}, \texttt{WRONG\_SYNTAX}, \texttt{RESOURCE\_EXHAUSTION}, \texttt{AUTH\_FAILURE}, \texttt{TIMEOUT}, \texttt{VERSION\_MISMATCH}; anything else resolves to \texttt{UNKNOWN}.

\reqid{Requirement FR-CLS-006:} The classification method actually used (\texttt{regex}, \texttt{embedding}, \texttt{llm}) and classification latency shall be recorded in the audit log for each attempt.

\subsection{Feature: Recovery Planning}

\reqid{Requirement FR-PLAN-001:} \texttt{\_build\_default\_plan} shall map each error class to a default strategy: \texttt{TOOL\_\allowbreak NOT\_\allowbreak INSTALLED$\rightarrow$INSTALL\_\allowbreak TOOL}, \texttt{PERMISSION\_\allowbreak DENIED$\rightarrow$ADD\_\allowbreak SUDO}, \texttt{NETWORK\_\allowbreak UNREACHABLE$\rightarrow$ADJUST\_\allowbreak PARAMETERS}, \texttt{WRONG\_\allowbreak SYNTAX$\rightarrow$FIX\_\allowbreak SYNTAX}, \texttt{RESOURCE\_\allowbreak EXHAUSTION$\rightarrow$ADJUST\_\allowbreak PARAMETERS}, \texttt{AUTH\_\allowbreak FAILURE$\rightarrow$FIX\_\allowbreak AUTH}, \texttt{TIMEOUT$\rightarrow$ADJUST\_\allowbreak TIMEOUT}, \texttt{VERSION\_\allowbreak MISMATCH$\rightarrow$VERSION\_\allowbreak DOWNGRADE}, \texttt{UNKNOWN$\rightarrow$ESCALATE}.

\reqid{Requirement FR-PLAN-002:} \texttt{RecoveryPlan} shall carry \texttt{error\_class}, \texttt{root\_cause} (10--500 chars), \texttt{strategy}, \texttt{corrected\_command} (3--1000 chars), \texttt{confidence} (0--1), and \texttt{reasoning}.

\reqid{Requirement FR-PLAN-003:} When the LLM reasoner returns a plan, its corrected command shall be used; otherwise the default plan's command (current command) is used as the mutation input.

\subsection{Feature: Command Mutation}

\subsubsection{Description}

\texttt{MutationEngine.\allowbreak apply\_\allowbreak mutation(plan, ctx)} routes on the plan strategy and produces the next command. Sources are strictly ordered per strategy: rule-based transformation, dataset lookup, then LLM generation.

\reqid{Requirement FR-MUT-001:} \texttt{INSTALL\_TOOL} shall produce \texttt{apt-\allowbreak get update -\allowbreak qq -\allowbreak -\allowbreak fix-\allowbreak missing \&\& apt-\allowbreak get install -\allowbreak y -\allowbreak -\allowbreak fix-\allowbreak missing <tool> \&\& <cmd>} only for tools in \texttt{KALI\_PACKAGE\_WHITELIST} (\{nmap, gobuster, nikto, ffuf, dirb, hydra, sqlmap, whatweb, masscan, medusa, ncrack, wfuzz, crackmapexec, enum4linux, netcat, curl, wget\}); otherwise no mutation is returned.

\reqid{Requirement FR-MUT-002:} \texttt{ADD\_SUDO} shall remove a leading \texttt{sudo } (or an embedded \texttt{ sudo }) from the command, since containers run as root.

\reqid{Requirement FR-MUT-003:} \texttt{ADJUST\_PARAMETERS} shall apply rule-based timing fixes (e.g., \texttt{-t N}$\rightarrow$\texttt{-t 10}, \texttt{-T[45]}$\rightarrow$\texttt{-T2}, \texttt{--timeout N}$\rightarrow$\texttt{--timeout 30}) before falling back to dataset, then LLM.

\reqid{Requirement FR-MUT-004:} \texttt{TOOL\_SUBSTITUTION} shall replace the tool via \texttt{ToolSubstitutionMapper.\allowbreak get\_\allowbreak substitute(tool, intent\_\allowbreak ref, prior\_\allowbreak strategies)} using \texttt{tool\_\allowbreak substitution\_\allowbreak map.\allowbreak json}.

\reqid{Requirement FR-MUT-005:} \texttt{FIX\_SYNTAX}, \texttt{ADJUST\_TIMEOUT}, \texttt{FIX\_AUTH}, and \texttt{VERSION\_DOWNGRADE} shall: fix missing-wordlist/file paths from stderr signals, use any LLM-provided corrected command (sanitized), then dataset, then LLM generation.

\reqid{Requirement FR-MUT-006:} Dataset corrections shall be rejected (returning \texttt{None}) when the user's original target cannot be preserved or when the correction is a tool name / shorter than 5 characters.

\reqid{Requirement FR-MUT-007:} Corrections containing listener/interactive patterns (\texttt{nc|netcat|ncat -l}, \texttt{-lvnp|-lvp|-lp}) shall be hard-rejected.

\reqid{Requirement FR-MUT-008:} For gobuster/dirb/ffuf/wfuzz corrections, \texttt{\_fix\_wordlist\_path} shall replace wordlist paths not in \texttt{VALID\_WORDLIST\_PATHS} with \texttt{DEFAULT\_\allowbreak WORDLIST = /\allowbreak usr/\allowbreak share/\allowbreak dirb/\allowbreak wordlists/\allowbreak common.\allowbreak txt}.

\reqid{Requirement FR-MUT-009:} \texttt{\_\allowbreak preserve\_\allowbreak original\_\allowbreak target} shall re-inject the user's original target into dataset/LLM corrections, or return \texttt{None} if the target cannot be identified (fail safe).

\reqid{Requirement FR-MUT-010:} \texttt{apply\_mutation} shall return \texttt{(corrected\_\allowbreak command, source)} where source ∈ \texttt{\{rule-\allowbreak based, dataset, llm\}}; \texttt{(None, '')} indicates no mutation possible.

\subsection{Feature: Safety Gate (4-Layer Validation)}

\subsubsection{Description}

\texttt{safety\_\allowbreak gate.\allowbreak validate(corrected\_\allowbreak command, plan)} must approve every command (original and corrected) before execution.

\reqid{Requirement FR-SAF-001:} Layer 1 --- an empty or whitespace-only command shall be rejected with \texttt{approved=False}.

\reqid{Requirement FR-SAF-002:} Layer 2 --- any match against the 16 \texttt{DANGER\_PATTERNS} (\texttt{rm -rf}, \texttt{rm -fr}, \texttt{mkfs.}, \texttt{dd if=}, \texttt{wget|curl ... |sh|bash}, \texttt{nc -e /bin}, \texttt{bash -i >\&}, \texttt{> /etc/passwd}, \texttt{> /etc/shadow}, \texttt{cat /etc/shadow}, \texttt{chmod (777|000|666) /}, recursive chmod on \texttt{/}, fork bomb \texttt{:()\{...\}}, \texttt{shred}, \texttt{wipefs}) shall reject the command and record the matched danger words.

\reqid{Requirement FR-SAF-003:} Layer 3 --- a recovery confidence below 0.60 shall reject the command (\texttt{confidence\_\allowbreak too\_\allowbreak low=True}).

\reqid{Requirement FR-SAF-004:} Layer 4 --- all extracted targets must be RFC-1918 private IPs (\texttt{10.0.0.0/8}, \texttt{172.16.0.0/12}, \texttt{192.168.0.0/16}), loopback (\texttt{127.0.0.0/8}), or the literal hostname \texttt{localhost}; any public IP or unapproved hostname shall be rejected as a scope violation.

\reqid{Requirement FR-SAF-005:} \texttt{SafetyCheckResult} shall report \texttt{approved}, \texttt{danger\_words\_found}, \texttt{scope\_violation}, \texttt{confidence\_too\_low}, and \texttt{reason}.

\reqid{Requirement FR-SAF-006:} The safety gate shall be applied to the initial command before the first execution (stage 0); a blocked initial command shall return an immediate \texttt{SAFETY\_BLOCKED} failure without executing anything.

\subsection{Feature: Retry Orchestration and Deadlines}

\subsubsection{Description}

\texttt{RetryOrchestrator} enforces per-error-class budgets from \texttt{error\_taxonomy.json}; the outer loop enforces a hard wall-clock deadline.

\reqid{Requirement FR-RET-001:} \texttt{should\_\allowbreak retry(error\_\allowbreak class, attempt)} shall permit retry while \texttt{attempt <= max\_retries} for the error class; the \texttt{UNKNOWN} default budget is \texttt{\{max\_\allowbreak retries:\allowbreak  1, backoff:\allowbreak  fixed, backoff\_\allowbreak seconds:\allowbreak  0\}}.

\reqid{Requirement FR-RET-002:} Backoff shall be fixed or exponential (\texttt{backoff\_\allowbreak seconds * 2\textasciicircum\{\}(attempt-\allowbreak 1)}) per the taxonomy entry.

\reqid{Requirement FR-RET-003:} Taxonomied budgets shall be: \texttt{TOOL\_NOT\_INSTALLED} 2, \texttt{PERMISSION\_DENIED} 2, \texttt{NETWORK\_UNREACHABLE} 3 (exponential, 2 s), \texttt{WRONG\_SYNTAX} 2, \texttt{RESOURCE\_EXHAUSTION} 2, \texttt{AUTH\_FAILURE} 2, \texttt{TIMEOUT} 3 (exponential, 3 s), \texttt{VERSION\_MISMATCH} 1.

\reqid{Requirement FR-RET-004:} When the budget is exhausted, the loop shall record \texttt{BUDGET\_EXHAUSTED} and stop.

\reqid{Requirement FR-RET-005:} The recovery loop shall abort as \texttt{failed} when elapsed time exceeds \texttt{AEERE\_DEADLINE\_SECONDS} (default 240 s), classifying an unresolved error as \texttt{TIMEOUT}.

\subsection{Feature: Audit Logging and Failure Reporting}

\subsubsection{Description}

\texttt{audit\_logger} persists every attempt to SQLite; \texttt{FailureReporter} builds structured failure reports when recovery fails.

\reqid{Requirement FR-AUD-001:} \texttt{init\_db(DB\_PATH)} shall create the SQLite \texttt{audit\_log} table; the audit log shall be append-only.

\reqid{Requirement FR-AUD-002:} Each attempt shall log: \texttt{session\_id}, \texttt{attempt\_number}, \texttt{original\_command}, \texttt{stderr}, \texttt{exit\_code}, \texttt{error\_class}, \texttt{classification\_method}, \texttt{classification\_\allowbreak latency\_\allowbreak ms}, \texttt{strategy\_applied}, \texttt{corrected\_command}, \texttt{safety\_approved}, \texttt{danger\_words}, \texttt{confidence}, \texttt{llm\_reasoning}, \texttt{outcome}, \texttt{total\_latency\_ms}.

\reqid{Requirement FR-AUD-003:} Audit outcomes shall include \texttt{approved}, \texttt{blocked}, \texttt{budget\_exhausted}, \texttt{no\_mutation}, \texttt{success}, \texttt{recovered}, and \texttt{failed}.

\reqid{Requirement FR-AUD-004:} \texttt{GET /\allowbreak sessions/\allowbreak \{session\_\allowbreak id\}} shall return the full audit trail for a session with event count.

\reqid{Requirement FR-RPT-001:} On final failure, \texttt{FailureReporter.\allowbreak build\_\allowbreak report(session\_\allowbreak id, error\_\allowbreak class, attempt, DB\_\allowbreak PATH)} shall produce a report with \texttt{attempts}, \texttt{final\_error\_class}, and a \texttt{recommendation}.

\reqid{Requirement FR-RPT-002:} A global exception handler shall convert unhandled errors into a structured \texttt{ExecutionResult} with \texttt{failure\_\allowbreak report.\allowbreak final\_\allowbreak error\_\allowbreak class='INTERNAL\_\allowbreak ERROR'} and HTTP 500.

\subsection{Feature: Session Context and Health}

\reqid{Requirement FR-SES-001:} \texttt{SessionContext} shall persist per-session prior recovery strategies in Redis (optional) and degrade to in-memory state on \texttt{redis.RedisError}.

\reqid{Requirement FR-SES-002:} Prior strategies shall feed tool-substitution decisions to avoid repeating failed strategies.

\reqid{Requirement FR-HLT-001:} \texttt{GET /health} shall report \texttt{status: healthy}, Docker availability, Redis availability, dataset correction availability (\texttt{is\_loaded}), dataset entry count, configured \texttt{OLLAMA\_MODEL}, and version.

\section{External Interface Requirements}

\subsection{Software Interfaces}

\subsubsection{REST API (HTTP/JSON)}

\begin{table}[h]
\caption{Table 4.}
\label{tab:t4}
\begin{tabularx}{\columnwidth}{l >{\raggedright\arraybackslash}X >{\raggedright\arraybackslash}X >{\raggedright\arraybackslash}X}
\toprule
\textbf{Method} & \textbf{Path} & \textbf{Auth} & \textbf{Description}\\
\midrule
POST & \texttt{/execute} & none (gateway-level) & Execute a plan with recovery loop\\
GET & \texttt{/\allowbreak sessions/\allowbreak \{session\_\allowbreak id\}} & none & Audit trail query\\
GET & \texttt{/health} & none & Health probe\\
\bottomrule
\end{tabularx}
\end{table}

\subsubsection{POST /execute --- Request (PlannerOutput)}

\begin{lstlisting}
{
  "command": "nmap -sS -p 22,80,443 192.168.1.1",
  "tool": "nmap",
  "session_id": "s-001",
  "intent_ref": "NETWORK_SCAN",
  "estimated_duration": "medium"
}
\end{lstlisting}

\subsubsection{POST /execute --- Response (200, ExecutionResult)}

\begin{lstlisting}
{
  "session_id": "s-001",
  "status": "success",
  "command_executed": "nmap -sS -p 22,80,443 192.168.1.1",
  "stdout": "...",
  "stderr": "",
  "exit_code": 0,
  "recovery_log": [],
  "correction_history": [],
  "failure_report": null,
  "tool_substituted": false,
  "latency_ms": 4210
}
\end{lstlisting}

Recovered example (status \texttt{recovered}) shows \texttt{correction\_history} entries with \texttt{attempt}, \texttt{original\_command}, \texttt{corrected\_command}, \texttt{error\_class}, \texttt{strategy}, \texttt{approved}, \texttt{correction\_source}.

\subsubsection{POST /execute --- Error behavior}

\begin{table}[h]
\caption{Table 5.}
\label{tab:t5}
\begin{tabularx}{\columnwidth}{l >{\raggedright\arraybackslash}X}
\toprule
\textbf{Condition} & \textbf{Result}\\
\midrule
Initial command blocked by safety gate & 200 \texttt{ExecutionResult} status \texttt{failed}, \texttt{failure\_\allowbreak report.\allowbreak final\_\allowbreak error\_\allowbreak class='SAFETY\_\allowbreak BLOCKED'}\\
Recovery loop error & 200 \texttt{ExecutionResult} status \texttt{failed} with failure report\\
Unhandled exception & HTTP 500 with \texttt{ExecutionResult} \texttt{INTERNAL\_ERROR}\\
\bottomrule
\end{tabularx}
\end{table}

\subsection{Hardware Interfaces}

\begin{itemize}
\item Docker Engine (local or remote via daemon socket) --- the sole execution substrate.
\item No other direct hardware interfaces; container CPU/memory limits are specified resource parameters.
\end{itemize}

\subsection{Software Interfaces (outbound)}

\begin{table}[h]
\caption{Table 6.}
\label{tab:t6}
\begin{tabularx}{\columnwidth}{l >{\raggedright\arraybackslash}X >{\raggedright\arraybackslash}X}
\toprule
\textbf{Interface} & \textbf{Protocol} & \textbf{Purpose}\\
\midrule
Docker Engine & Docker SDK (\texttt{docker.from\_env}) & container lifecycle, logs, wait, remove\\
Ollama & local HTTP (\texttt{ollama} client) & LLM classification + command correction\\
Redis & TCP (\texttt{redis.from\_url}) & optional session state\\
SQLite & file-backed (\texttt{sqlite3}) & audit log persistence\\
SentenceTransformers & in-process & embedding classifier\\
\bottomrule
\end{tabularx}
\end{table}

\subsection{Communication Interfaces}

\begin{itemize}
\item \textbf{Transport}: HTTP/1.1; JSON over Pydantic v2 models
\item \textbf{CORS}: \texttt{allow\_origins=['*']}, all methods/headers (gateway-terminated auth is expected in production)
\item \textbf{Timing}: the loop guarantees a returned result before the upstream chain timeout (\textasciitilde\{\}300 s)
\end{itemize}

\section{Non-Functional Requirements}

\subsection{Performance}

\reqid{Requirement NFR-PERF-001:} The recovery loop shall respect \texttt{AEERE\_DEADLINE\_SECONDS} (default 240 s) so it always returns before upstream client timeouts.

\reqid{Requirement NFR-PERF-002:} Per-container exec timeout shall be \texttt{CONTAINER\_EXEC\_TIMEOUT} (default 60 s); timeout runs shall be classified \texttt{TIMEOUT} deterministically.

\reqid{Requirement NFR-PERF-003:} Classification latency per attempt shall be measured (\texttt{classification\_\allowbreak latency\_\allowbreak ms}) and audited.

\reqid{Requirement NFR-PERF-004:} Deterministic (regex) classification shall be preferred to minimize LLM invocations; the embedding classifier runs only when regex fails; the LLM only when both fail.

\reqid{Requirement NFR-PERF-005:} End-to-end \texttt{latency\_ms} shall be reported in every \texttt{ExecutionResult}.

\subsection{Security}

\reqid{Requirement NFR-SEC-001:} Corrected commands shall never include danger patterns (\texttt{DANGER\_PATTERNS}), destructively recursive chmod on \texttt{/}, \texttt{dd if=}, or fork bombs.

\reqid{Requirement NFR-SEC-002:} All execution targets shall be restricted to RFC-1918 private ranges, loopback, or \texttt{localhost}; public-IP/domain targets shall be blocked by the safety gate.

\reqid{Requirement NFR-SEC-003:} The container's wordlist autofix shall only substitute against the \texttt{VALID\_WORDLIST\_PATHS} allowlist.

\reqid{Requirement NFR-SEC-004:} Prompt injection content (\texttt{\_strip\_injections}) shall be redacted from stdout/stderr before LLM exposure.

\reqid{Requirement NFR-SEC-005:} Listener/interactive commands (\texttt{nc -l}, \texttt{-lvnp}/\texttt{-lvp}/\texttt{-lp}) shall never be auto-corrected into execution.

\reqid{Requirement NFR-SEC-006:} Containers shall be force-removed after every execution (no container reuse/leakage).

\reqid{Requirement NFR-SEC-007:} Secrets and credentials in commands (e.g., from \texttt{FIX\_AUTH} mutations) shall not be committed; audit data shall not expose plaintext secrets beyond the command text required for traceability.

\reqid{Requirement NFR-SEC-008:} Auto-install of tools shall be limited to \texttt{KALI\_PACKAGE\_WHITELIST}; tools like \texttt{rustscan} (GitHub-release only) are excluded by design.

\subsection{Reliability}

\reqid{Requirement NFR-REL-001:} Dataset load failure shall not prevent startup (graceful skip with log).

\reqid{Requirement NFR-REL-002:} Embedding-model unavailability shall not prevent recovery (only the embedding tier is skipped).

\reqid{Requirement NFR-REL-003:} Audit logging failures on the success path shall not alter the returned result (audit failure is logged, result still returned).

\reqid{Requirement NFR-REL-004:} Redis unavailability shall degrade to in-memory session state.

\reqid{Requirement NFR-REL-005:} Every loop exit path shall return a valid, structured \texttt{ExecutionResult} (never an unstructured exception).

\reqid{Requirement NFR-REL-006:} The audit trail shall be append-only with no delete operations exposed via API.

\subsection{Maintainability}

\reqid{Requirement NFR-MNT-001:} Error taxonomy shall be data-driven (\texttt{app/\allowbreak taxonomy/\allowbreak error\_\allowbreak taxonomy.\allowbreak json}) and reloadable without code changes.

\reqid{Requirement NFR-MNT-002:} Tool substitution mapping shall be data-driven (\texttt{app/\allowbreak taxonomy/\allowbreak tool\_\allowbreak substitution\_\allowbreak map.\allowbreak json}).

\reqid{Requirement NFR-MNT-003:} Kali error dataset shall be a labeled CSV (\texttt{app/\allowbreak kali\_\allowbreak data/\allowbreak kali\_\allowbreak linux\_\allowbreak error.\allowbreak csv}) driving both classifier enrichment and dataset corrections.

\reqid{Requirement NFR-MNT-004:} Module boundaries (execution, capture, classification, mutation, safety, retry, reporting, session) shall be independently unit-testable.

\reqid{Requirement NFR-MNT-005:} Configuration shall be environment-driven (\texttt{.env}, Docker env): \texttt{PORT}, \texttt{AUDIT\_DB\_PATH}, \texttt{KALI\_IMAGE}, \texttt{CONTAINER\_*}, \texttt{OLLAMA\_MODEL}, \texttt{AEERE\_DEADLINE\_SECONDS}.

\subsection{Portability}

\reqid{Requirement NFR-PORT-001:} The service shall run on Linux (production), Windows (development via \texttt{startup.ps1}), and macOS, with Docker Engine as the execution dependency.

\reqid{Requirement NFR-PORT-002:} The Dockerfile shall be multi-stage with a \texttt{python:3.11-slim} base; \texttt{Dockerfile.kali} provides the Kali image with test target.

\reqid{Requirement NFR-PORT-003:} The standalone \texttt{docker-compose.yml} shall provision Redis + AEERE together.

\section{Data Requirements}

\subsection{Data Models (app/src/schemas/models.py)}

\begin{table}[h]
\caption{Table 7.}
\label{tab:t7}
\begin{tabularx}{\columnwidth}{l >{\raggedright\arraybackslash}X}
\toprule
\textbf{Model} & \textbf{Purpose / key fields}\\
\midrule
\texttt{PlannerOutput} & inbound plan: \texttt{command}, \texttt{tool}, \texttt{session\_id}, \texttt{intent\_ref}, \texttt{estimated\_duration}\\
\texttt{RawExecutionResult} & \texttt{command}, \texttt{stdout}, \texttt{stderr}, \texttt{exit\_code}, \texttt{timed\_out}\\
\texttt{ErrorContext} & \texttt{session\_id}, \texttt{command}, \texttt{stdout}, \texttt{stderr}, \texttt{exit\_code}, \texttt{tool}, \texttt{intent\_ref}, \texttt{attempt\_number}, \texttt{prior\_strategies}, \texttt{silent\_failure}\\
\texttt{RecoveryPlan} & \texttt{error\_class} (9-literal), \texttt{root\_cause} (10--500), \texttt{strategy} (9-literal), \texttt{corrected\_command} (3--1000), \texttt{confidence} (0--1), \texttt{reasoning}\\
\texttt{SafetyCheckResult} & \texttt{approved}, \texttt{danger\_words\_found}, \texttt{scope\_violation}, \texttt{confidence\_too\_low}, \texttt{reason}\\
\texttt{RecoveryLogEntry} & \texttt{attempt}, \texttt{error\_class}, \texttt{strategy}, \texttt{clf\_method}, \texttt{confidence}, \texttt{original\_command}, \texttt{corrected\_command}\\
\texttt{CorrectionEntry} & \texttt{attempt}, \texttt{original\_command}, \texttt{corrected\_command}, \texttt{error\_class}, \texttt{strategy}, \texttt{approved}, \texttt{correction\_source}\\
\texttt{ExecutionResult} & \texttt{session\_id}, \texttt{status} (\texttt{success|recovered|failed}), \texttt{command\_executed}, \texttt{stdout}, \texttt{stderr}, \texttt{exit\_code}, \texttt{recovery\_log}, \texttt{correction\_history}, \texttt{failure\_report}, \texttt{tool\_substituted}, \texttt{latency\_ms}\\
\bottomrule
\end{tabularx}
\end{table}

\subsection{Error Taxonomy (app/taxonomy/error\_taxonomy.json)}

Eight error classes, each with detection \texttt{patterns}, recovery \texttt{strategy}, \texttt{max\_retries}, and \texttt{backoff}:

\begin{table}[h]
\caption{Table 8.}
\label{tab:t8}
\begin{tabularx}{\columnwidth}{l >{\raggedright\arraybackslash}X >{\raggedright\arraybackslash}X >{\raggedright\arraybackslash}X}
\toprule
\textbf{Error class} & \textbf{Strategy} & \textbf{Max retries} & \textbf{Backoff}\\
\midrule
\texttt{TOOL\_NOT\_INSTALLED} & \texttt{INSTALL\_TOOL} & 2 & fixed 0 s\\
\texttt{PERMISSION\_DENIED} & \texttt{ADD\_SUDO} & 2 & fixed 1 s\\
\texttt{NETWORK\_UNREACHABLE} & \texttt{ADJUST\_PARAMETERS} & 3 & exponential 2 s\\
\texttt{WRONG\_SYNTAX} & \texttt{FIX\_SYNTAX} & 2 & fixed 1 s\\
\texttt{RESOURCE\_EXHAUSTION} & \texttt{ADJUST\_PARAMETERS} & 2 & fixed 2 s\\
\texttt{AUTH\_FAILURE} & \texttt{FIX\_AUTH} & 2 & fixed 1 s\\
\texttt{TIMEOUT} & \texttt{ADJUST\_TIMEOUT} & 3 & exponential 3 s\\
\texttt{VERSION\_MISMATCH} & \texttt{VERSION\_DOWNGRADE} & 1 & fixed 0 s\\
\bottomrule
\end{tabularx}
\end{table}

\subsection{Tool Substitution Map (app/taxonomy/tool\_substitution\_map.json)}

\begin{table}[h]
\caption{Table 9.}
\label{tab:t9}
\begin{tabularx}{\columnwidth}{l >{\raggedright\arraybackslash}X}
\toprule
\textbf{Intent} & \textbf{Tool $\rightarrow$ substitutes}\\
\midrule
\texttt{DIRECTORY\_BRUTEFORCE} & gobuster$\rightarrow$[ffuf, dirb, wfuzz]; ffuf$\rightarrow$[gobuster, dirb, wfuzz]; dirb$\rightarrow$[gobuster, ffuf, wfuzz]\\
\texttt{VULNERABILITY\_AUDIT} & nikto$\rightarrow$[whatweb, nmap --script vuln]; whatweb$\rightarrow$[nikto]\\
\texttt{NETWORK\_SCAN} & nmap$\rightarrow$[masscan, rustscan]; masscan$\rightarrow$[nmap, rustscan]\\
\texttt{PASSWORD\_ATTACK} & hydra$\rightarrow$[medusa, ncrack, crackmapexec]; medusa$\rightarrow$[hydra, ncrack]\\
\bottomrule
\end{tabularx}
\end{table}

\subsection{Kali Error Dataset (app/kali\_data/kali\_linux\_error.csv)}

\begin{itemize}
\item Labeled error/correction scenario pairs used for classifier enrichment (\texttt{regex\_\allowbreak clf.\allowbreak enrich\_\allowbreak from\_\allowbreak dataset}) and dataset corrections (\texttt{ds\_\allowbreak corrector.\allowbreak load\_\allowbreak from\_\allowbreak scenarios})
\item Loaded at startup; failure to load is non-fatal and reported in \texttt{/health}
\end{itemize}

\subsection{Audit Store (app/data/audit.db)}

\begin{itemize}
\item SQLite database; table \texttt{audit\_log}; append-only; per-attempt event records as defined in FR-AUD-002
\item Queryable per session via \texttt{GET /\allowbreak sessions/\allowbreak \{session\_\allowbreak id\}}
\end{itemize}

\subsection{Session State}

\begin{itemize}
\item Redis optional (keyed per \texttt{session\_id}, tracks \texttt{prior\_strategies}); in-memory fallback on Redis errors
\end{itemize}

\section{Verification and Validation}

\subsection{Test Suite}

\begin{table}[h]
\caption{Table 10.}
\label{tab:t10}
\begin{tabularx}{\columnwidth}{l >{\raggedright\arraybackslash}X >{\raggedright\arraybackslash}X}
\toprule
\textbf{Area} & \textbf{File} & \textbf{Coverage}\\
\midrule
Recovery pipeline & \texttt{app/\allowbreak tests/\allowbreak test\_\allowbreak execution\_\allowbreak engine.\allowbreak py} & full 8-step pipeline (single execution + recovery path)\\
Tool-specific recovery & \texttt{app/\allowbreak tests/\allowbreak test\_\allowbreak tool\_\allowbreak specific\_\allowbreak recovery.\allowbreak py} & live API recovery scenarios per tool\\
Error capture & \texttt{app/\allowbreak tests/\allowbreak unit/\allowbreak test\_\allowbreak error\_\allowbreak capture.\allowbreak py} & 20 methods: context building, injection stripping, tool extraction\\
Regex classifier & \texttt{app/\allowbreak tests/\allowbreak unit/\allowbreak test\_\allowbreak regex\_\allowbreak classifier.\allowbreak py} & 14 methods: pattern matching, dataset enrichment\\
Safety gate & \texttt{app/\allowbreak tests/\allowbreak unit/\allowbreak test\_\allowbreak safety\_\allowbreak gate.\allowbreak py} & 34 methods: all 4 layers, danger patterns, scope, confidence\\
Silent failure detector & \texttt{app/\allowbreak tests/\allowbreak unit/\allowbreak test\_\allowbreak silent\_\allowbreak failure\_\allowbreak detector.\allowbreak py} & 14 methods: output-requiring tools, stderr override, version/help exemption\\
\bottomrule
\end{tabularx}
\end{table}

\reqid{Requirement NFR-TST-001:} The safety-gate unit suite shall cover all 16 danger patterns and all four validation layers.

\reqid{Requirement NFR-TST-002:} The silent-failure suite shall cover all eight output-requiring tools plus stderr override and exemption cases.

\subsection{Benchmarks}

\begin{itemize}
\item \texttt{benchmark\_safety.py} --- 10-scenario safety block-rate benchmark
\item \texttt{benchmark\_recovery.py} --- 15-scenario recovery-rate benchmark
\item \texttt{benchmark\_f1.py}, \texttt{benchmark\_hard.py}, \texttt{benchmark\_baselines.py}, \texttt{benchmark\_pilot.py} --- recovery/classification measurement
\item \texttt{benchmark\_baselines.py} etc. produce \texttt{.results.json} artifacts in the component root
\end{itemize}

\reqid{Requirement NFR-VER-001:} Delivery validation shall run the full unit suite (\texttt{app/tests/unit}), the engine test, and the tool-specific recovery test.

\reqid{Requirement NFR-VER-002:} \texttt{/health} shall report \texttt{status: healthy} with Docker available in the target environment.

\reqid{Requirement NFR-VER-003:} A smoke \texttt{POST /execute} for a well-formed plan shall return status \texttt{success} for a clean command and \texttt{recovered} for a command that fails once then is corrected (e.g., a timing-parameter fix).

\reqid{Requirement NFR-VER-004:} A plan targeting a public IP or containing a danger pattern shall be blocked (status \texttt{failed}, class \texttt{SAFETY\_BLOCKED}).

\section{Deployment}

\subsection{Prerequisites}

\begin{itemize}
\item Docker Engine running locally (sandbox requirement)
\item \texttt{kalilinux/kali-rolling} image (default) or a built \texttt{neuroshell-kali} image with wordlists at \texttt{VALID\_WORDLIST\_PATHS}
\item Ollama reachable with \texttt{OLLAMA\_MODEL} (default \texttt{qwen2.5-coder:latest})
\item Optional Redis at \texttt{REDIS\_URL}
\end{itemize}

\subsection{Env Configuration}

\begin{table}[h]
\caption{Table 11.}
\label{tab:t11}
\begin{tabularx}{\columnwidth}{l >{\raggedright\arraybackslash}X >{\raggedright\arraybackslash}X}
\toprule
\textbf{Variable} & \textbf{Default} & \textbf{Purpose}\\
\midrule
\texttt{PORT} & \texttt{8000} & uvicorn port\\
\texttt{AUDIT\_DB\_PATH} & \texttt{./data/audit.db} & SQLite audit store\\
\texttt{KALI\_IMAGE} & \texttt{kalilinux/kali-rolling} & execution image\\
\texttt{CONTAINER\_MEM\_LIMIT} & \texttt{512m} & container memory\\
\texttt{CONTAINER\_CPU\_QUOTA} & \texttt{50000} & CPU quota\\
\texttt{CONTAINER\_EXEC\_TIMEOUT} & \texttt{60} & per-exec timeout (s)\\
\texttt{AEERE\_DEADLINE\_SECONDS} & \texttt{240} & loop wall-clock deadline\\
\texttt{OLLAMA\_MODEL} & \texttt{qwen2.5-coder:latest} & LLM model\\
\texttt{REDIS\_URL} & optional & session state\\
\bottomrule
\end{tabularx}
\end{table}

\subsection{Build and Start}

\begin{lstlisting}
# option A: standalone compose (Redis + AEERE)
docker compose -f app/docker-compose.yml up

# option B: build service image
docker build -t neuroshell-aeere .
docker run -p 8000:8000 -e OLLAMA_MODEL=qwen2.5-coder:latest neuroshell-aeere

# Windows development
app/startup.ps1
\end{lstlisting}

\subsection{Verify}

\begin{lstlisting}
curl http://localhost:8000/health
curl -X POST http://localhost:8000/execute \
  -H "Content-Type: application/json" \
  -d '{"command":"nmap -sV 192.168.1.1","tool":"nmap","session_id":"s-001","intent_ref":"NETWORK_SCAN","estimated_duration":"medium"}'
curl http://localhost:8000/sessions/s-001
\end{lstlisting}

\section{Appendices}

\subsection{Appendix A --- Error Class / Strategy Enumeration}

\textbf{Error classes (9):} \texttt{TOOL\_NOT\_INSTALLED}, \texttt{PERMISSION\_DENIED}, \texttt{NETWORK\_UNREACHABLE}, \texttt{WRONG\_SYNTAX}, \texttt{RESOURCE\_EXHAUSTION}, \texttt{AUTH\_FAILURE}, \texttt{TIMEOUT}, \texttt{VERSION\_MISMATCH}, \texttt{UNKNOWN}.

\textbf{Strategies (9):} \texttt{INSTALL\_TOOL}, \texttt{ADD\_SUDO}, \texttt{ADJUST\_PARAMETERS}, \texttt{TOOL\_SUBSTITUTION}, \texttt{FIX\_SYNTAX}, \texttt{ADJUST\_TIMEOUT}, \texttt{FIX\_AUTH}, \texttt{VERSION\_DOWNGRADE}, \texttt{ESCALATE}.

\textbf{Correction sources (3):} \texttt{rule-based}, \texttt{dataset}, \texttt{llm}.

\subsection{Appendix B --- Recovery Loop Pseudocode}

\begin{lstlisting}
START (deadline = now + AEERE_DEADLINE_SECONDS)
attempt = 0
while now < deadline:
    attempt += 1
    if attempt == 1:  safety_validate(original)      # stage 0
                       if not approved → SAFETY_BLOCKED
    raw = run_in_docker(command)                      # stage 1
    if raw.exit_code == 0 and not silent_failure: success/recovered → DONE
    error_class, conf = classify(raw)                 # stage 2 (regex→embed→llm)
    if not retry_budget(error_class, attempt):        # stage 3
        BUDGET_EXHAUSTED → DONE
    plan = build_plan(error_class, conf, command)     # stage 4
    corrected, src = mutate(plan, ctx)                # stage 5
    if corrected is None: NO_MUTATION → DONE
    safety = safety_validate(corrected, plan)         # stage 6
    audit_attempt(...)                                # stage 7
    if not safety.approved: SAFETY_BLOCKED → DONE
    commit correction, update session, backoff        # stages 8–9
END
return ExecutionResult(status, logs, history, report, latency)
\end{lstlisting}

\subsection{Appendix C --- Safety Gate Danger Patterns (16)}

\texttt{rm -rf}, \texttt{rm -fr}, \texttt{mkfs.}, \texttt{dd if=}, \texttt{wget|...|sh}, \texttt{curl|...|bash}, \texttt{nc -e /bin}, \texttt{bash -i >\&}, \texttt{> /etc/passwd}, \texttt{> /etc/shadow}, \texttt{cat /etc/shadow}, \texttt{chmod (777|000|666) /}, recursive \texttt{chmod -R ... /}, fork bomb \texttt{:()\{...\}}, \texttt{shred}, \texttt{wipefs}.

\subsection{Appendix D --- Wordlist Allowlist}

\texttt{/\allowbreak usr/\allowbreak share/\allowbreak dirb/\allowbreak wordlists/\allowbreak \{common,big,small\}.\allowbreak txt}, \texttt{/\allowbreak usr/\allowbreak share/\allowbreak dirbuster/\allowbreak wordlists/\allowbreak directory-\allowbreak list-\allowbreak 2.\allowbreak 3-\allowbreak medium.\allowbreak txt}, \texttt{/\allowbreak usr/\allowbreak share/\allowbreak seclists/\allowbreak Discovery/\allowbreak Web-\allowbreak Content/\allowbreak common.\allowbreak txt}.

\subsection{Appendix E --- C3 $\rightarrow$ C4 Hand-off Contract}

\texttt{ExecutionResult} fields consumed by C4 AVAE (\texttt{C3ExecutionResult} in the AVAE SRS): \texttt{session\_id}, \texttt{status} (\texttt{success|recovered|failed}), \texttt{command\_executed}, \texttt{stdout}, \texttt{stderr}, \texttt{exit\_code}, \texttt{tool}, \texttt{intent\_ref}, \texttt{latency\_ms}.


\srsdivider{Component 04 --- AI Vulnerability Analysis Engine (AVAE)}
\section{Introduction}

\subsection{Purpose}

This Software Requirements Specification (SRS) defines the complete functional and non-functional requirements for the \textbf{NeuroShell AI Vulnerability Analysis Engine (AVAE)}, designated Component 04 of the NeuroShell Framework. AVAE is the vulnerability intelligence and prioritization layer that ingests scanner execution results from Component 03 (AEERE), parses tool output, extracts network entities, enriches CVEs from authoritative threat-intelligence sources (NVD, CISA KEV, FIRST EPSS), predicts exploitability using a trained machine-learning model with SHAP explanations, and returns triaged, ranked vulnerability findings.

This document is intended for:

\begin{itemize}
\item \textbf{Research evaluators} assessing the AVAE design and predictive performance
\item \textbf{Developers} maintaining or extending parsers, enrichment, ML scoring, triage, and graph persistence
\item \textbf{Upstream / downstream consumers} (Component 03 feeding results; Component 01 orchestration and frontend displaying findings) integrating with the AVAE API
\item \textbf{Examiners} reviewing the MSc Cybersecurity research project (ref R26-CS-018)
\end{itemize}

\subsection{Scope}

\subsubsection{Product Scope}

The AVAE is a \textbf{self-contained FastAPI microservice} that:

\begin{enumerate}
\item Accepts a \texttt{C3ExecutionResult} JSON (session\_id, status, command\_executed, stdout, stderr, exit\_code, tool, intent\_ref, latency\_ms) from Component 03
\item Parses scanner outputs (Nmap XML/plain-text, Nikto JSON/plain-text, Gobuster, generic fallback) into structured entities
\item Enriches detected CVEs with NVD metadata (CVSS v2/v3, CWEs, age, Exploit-DB refs), CISA KEV known-exploited status, and FIRST EPSS score/percentile
\item Predicts CVE exploitability with a trained XGBoost classifier and explains predictions with SHAP feature attributions
\item Falls back to a heuristic service-risk scorer when no CVE is identifiable (\texttt{analysis\_\allowbreak mode="service\_\allowbreak based\_\allowbreak fallback"})
\item Triages findings --- filtering, deduplicating, ranking, and computing a false-positive reduction rate
\item Persists the analysis as a knowledge graph in Neo4j (Session $\rightarrow$ Host $\rightarrow$ Service $\rightarrow$ Vulnerability)
\item Returns a structured JSON analysis to the caller for display and follow-up
\end{enumerate}

\subsubsection{In Scope}

\begin{itemize}
\item All application source under \texttt{app/src/} (API, parsers, ingestion, intelligence clients, ML, triage, graph, feedback, schemas)
\item ML artifacts under \texttt{app/models/} (xgboost model JSON, SHAP images)
\item Data processing scripts under \texttt{app/scripts/} (dataset build, merge, training, prediction test)
\item Raw/processed datasets and offline cache under \texttt{app/data/}
\item Test suite under \texttt{app/tests/} (unit parser + ML scoring)
\item Component documentation (\texttt{app/README.md}) and deployment artifacts (Dockerfile, compose, CI)
\end{itemize}

\subsubsection{Out of Scope}

\begin{itemize}
\item Component 01 (IRE), Component 02 (Planner), Component 03 (AEERE) --- referenced only for interface design
\item Actual vulnerability exploitation --- AVAE only analyzes and prioritizes
\item Frontend UI and end-user-facing presentation (separate project)
\item The \texttt{reporting} package (currently an empty placeholder) --- out of the implemented feature set
\item Training-dataset expansion and model re-training pipeline execution (scripts exist; runtime model is a shipped artifact)
\end{itemize}

\subsection{Definitions, Acronyms, Abbreviations}

\begin{table}[h]
\caption{Table 1.}
\label{tab:t1}
\begin{tabularx}{\columnwidth}{l >{\raggedright\arraybackslash}X}
\toprule
\textbf{Term} & \textbf{Definition}\\
\midrule
C1/C2/C3/C4 & Component 01 (IRE) / 02 (Planner) / 03 (AEERE) / 04 (AVAE)\\
AVAE & AI Vulnerability Analysis Engine\\
NVD & National Vulnerability Database (NIST CVE API 2.0)\\
KEV & CISA Known Exploited Vulnerabilities catalog\\
EPSS & Exploit Prediction Scoring System (FIRST.org)\\
CVSS & Common Vulnerability Scoring System (v2/v3)\\
CVE & Common Vulnerabilities and Exposures\\
CWE & Common Weakness Enumeration\\
CPE & Common Platform Enumeration\\
SHAP & SHapley Additive exPlanations (feature attributions)\\
XGBoost & Gradient-boosted decision tree classifier (runtime predictor)\\
FP-reduction & False-positive reduction rate: \texttt{1 - ranked/total\_input}\\
EntityCollection & Parsed entities: hosts, services, raw CVEs, paths\\
EnrichedVulnerability & CVE + NVD/KEV/EPSS + operational context\\
VulnerabilityPrediction & exploit probability, composite risk score, risk tier, SHAP top features\\
Risk tier & CRITICAL / HIGH / MEDIUM / LOW / FILTERED\\
CCEE & Corporate/community CVE extraction dataset (training source)\\
C3ExecutionResult & C3$\rightarrow$C4 inbound execution result schema\\
\bottomrule
\end{tabularx}
\end{table}

\subsection{References}

\begin{table}[h]
\caption{Table 2.}
\label{tab:t2}
\begin{tabularx}{\columnwidth}{l >{\raggedright\arraybackslash}X >{\raggedright\arraybackslash}X}
\toprule
\textbf{Ref ID} & \textbf{Document} & \textbf{Location}\\
\midrule
R1 & Component 04 README / project documentation & \texttt{app/README.md}\\
R2 & C3 AEERE SRS (upstream execution results) & \texttt{adaptive-\allowbreak execution-\allowbreak error-\allowbreak recovery/\allowbreak SRS-\allowbreak Doc/\allowbreak }\\
R3 & C2 Dynamic Planner SRS (planner context) & \texttt{dynamic-\allowbreak planner-\allowbreak rag/\allowbreak SRS-\allowbreak Doc/\allowbreak }\\
R4 & NVD CVE API 2.0 & services.nvd.nist.gov\\
R5 & CISA KEV feed & cisa.gov\\
R6 & FIRST EPSS API & api.first.org/data/v1/epss\\
R7 & IEEE 830-1998 SRS Standard & IEEE Computer Society\\
R8 & OWASP Top 10 (2021) & owasp.org\\
R9 & CVE Format Standard & mitre.org\\
R10 & CCEE training dataset & \texttt{app/data/raw/CCEE.csv}\\
\bottomrule
\end{tabularx}
\end{table}

\subsection{Overview}

The remainder of this document is organized as follows:

\begin{itemize}
\item \textbf{Section 2} describes the overall product perspective, user classes, and operating environment
\item \textbf{Section 3} specifies all system features and functional requirements in detail
\item \textbf{Section 4} defines external interface requirements (API, hardware, software, communication)
\item \textbf{Section 5} specifies non-functional requirements (performance, security, reliability)
\item \textbf{Section 6} covers data requirements (entities, enrichment, ML, persistence)
\item \textbf{Section 7} describes verification, validation, and empirical results
\item \textbf{Section 8} covers deployment procedures
\item \textbf{Section 9} provides appendices with enumerations, thresholds, and schemas
\end{itemize}

\section{Overall Description}

\subsection{Product Perspective}

The AVAE is \textbf{Component 04} of the NeuroShell Framework --- the final intelligence stage of the operation pipeline. It converts raw scanner output into prioritized, evidence-backed vulnerability findings.

\begin{lstlisting}
C2 Planner ──▶ C3 AEERE ──ExecutionResult──▶ C4 AVAE (this component)
                                            │  1. parse scanner output
                                            │  2. extract entities (hosts, services, CVEs, paths)
                                            │  3. enrich (NVD, KEV, EPSS)
                                            │  4. predict exploitability (XGBoost + SHAP)
                                            │  5. fallback heuristic service risk
                                            │  6. triage / rank / FP-reduction
                                            │  7. persist to Neo4j graph
                                            ▼
                                        JSON analysis → C1 orchestration / frontend
\end{lstlisting}

\subsection{Product Functions Summary}

\begin{itemize}
\item \textbf{Scanner parsing}: \texttt{NmapParser} (XML + plain-text), \texttt{NiktoParser} (JSON + plain-text), \texttt{GobusterParser} (classic + fallback), \texttt{GenericParser} (last-resort CVE/IP extraction), dispatched by \texttt{ParserRegistry}
\item \textbf{Entity extraction}: hosts (IPv4), open ports, protocols, service names/versions, banners, CPEs, CVE hints (from NSE/CVSS output), discovered paths
\item \textbf{Threat intelligence enrichment}: \texttt{NVDClient} (CVE metadata + keyword search), \texttt{KEVClient} (CISA known-exploited), \texttt{EPSSClient} (score + percentile, batch)
\item \textbf{Intelligence caching}: \texttt{CacheManager} (SQLite \texttt{cve\_cache} + in-memory dict)
\item \textbf{Exploitability prediction}: \texttt{ExploitabilityPredictor} --- XGBClassifier (300 trees, 5 features) + \texttt{shap.TreeExplainer}; composite risk score and \texttt{\_risk\_tier} thresholds
\item \textbf{Heuristic fallback}: \texttt{ServiceRiskEngine} --- port/service/version additive scoring when no CVE is available
\item \textbf{Triage}: filter below threshold (default 0.20), dedupe by CVE, rank by composite risk score, compute FP-reduction rate
\item \textbf{Graph persistence}: \texttt{Neo4jPersistence} --- Session/Host/Service/Vulnerability with \texttt{HAS\_SERVICE}, \texttt{HAS\_VULNERABILITY}, \texttt{FOUND\_VULNERABILITY} relationships
\item \textbf{Feedback generation}: \texttt{FeedbackIntentGenerator} maps ranked findings to follow-up intents for C2 (implemented; not yet wired into the API response)
\end{itemize}

\subsection{User Classes and Characteristics}

\begin{table}[h]
\caption{Table 3.}
\label{tab:t3}
\begin{tabularx}{\columnwidth}{l >{\raggedright\arraybackslash}X >{\raggedright\arraybackslash}X}
\toprule
\textbf{User class} & \textbf{Description} & \textbf{Interactions}\\
\midrule
\textbf{End users / operator} & Human operator via frontend & reads ranked findings, initiates follow-ups\\
\textbf{Upstream C3 (AEERE)} & Submits \texttt{C3ExecutionResult} for analysis & \texttt{POST /analyze}\\
\textbf{C1 orchestration / frontend} & Displays AVAE analysis within \texttt{/execute} chain & reads \texttt{/analyze} response via \texttt{VulnerabilityClient}\\
\textbf{Analyst} & Reviews Neo4j graph and cached enrichment & graph traversal, cache inspection\\
\textbf{Maintainer / researcher} & Retrains models, extends parsers, ingests datasets & \texttt{app/scripts/*.py}, tests\\
\bottomrule
\end{tabularx}
\end{table}

\subsection{Operating Environment}

\begin{itemize}
\item \textbf{Runtime}: Python 3.9+ (deployed image \texttt{python:3.11-slim}), FastAPI + Uvicorn
\item \textbf{Deployment}: Docker container (\texttt{neuroshell-\allowbreak ai-\allowbreak vulnerability-\allowbreak analysis}), behind Kong gateway (upstream \texttt{ai\_\allowbreak vulnerability\_\allowbreak analysis} at 8002)
\item \textbf{External services}: NVD API, CISA KEV feed, FIRST EPSS API, Neo4j (\texttt{bolt://localhost:7687}, optional fail-open)
\item \textbf{Model}: \texttt{app/\allowbreak models/\allowbreak xgboost\_\allowbreak model.\allowbreak json} (XGBClassifier JSON, 300 trees, 5 features)
\item \textbf{Ports}: container 8000; documented service identities and compose map 8002/\texttt{/} endpoint reports 8004 (see deployment notes)
\item \textbf{OS}: Linux (production); Windows/macOS supported for development through \texttt{.venv}
\end{itemize}

\subsection{Design and Implementation Constraints}

\begin{itemize}
\item \textbf{Fail-open intelligence}: NVD/KEV/EPSS/Neo4j/SHAP lookups fail open with defaults so pipeline analysis continues under degraded conditions
\item \textbf{Deterministic structure}: server startup loads the KEV feed into memory (\texttt{KEV loaded: N CVEs}); parser dispatch is first-match-wins
\item \textbf{Runtime/model parity}: the shipped runtime model uses exactly 5 \texttt{FEATURE\_NAMES} (\texttt{cvss\_\allowbreak v3\_\allowbreak score, epss\_\allowbreak score, epss\_\allowbreak percentile, age\_\allowbreak days, cvss\_\allowbreak epss\_\allowbreak product}); the 12-feature set belongs to the offline dataset pipeline
\item \textbf{No LLM in the analysis path}: AVAE consumes structured output only, minimizing prompt-injection surface
\item \textbf{Persistent cache discipline}: enrichment results are keyed by \texttt{cve\_id} and cached to SQLite for offline repeatability
\end{itemize}

\subsection{Assumptions and Dependencies}

\begin{itemize}
\item NVD/KEV/EPSS endpoints are reachable (with graceful degradation when not)
\item \texttt{app/\allowbreak models/\allowbreak xgboost\_\allowbreak model.\allowbreak json} and SHAP explainer are loadable (SHAP load is non-fatal)
\item Neo4j is reachable at \texttt{NEO4J\_URI} for graph persistence (unreachable $\rightarrow$ \texttt{available=False}, still returns analysis)
\item Scanner outputs conform to their tool's standard formats (XML for nmap when available; plain-text fallback otherwise)
\end{itemize}

\section{System Features and Functional Requirements}

\subsection{Feature: Analysis Orchestration (POST /analyze)}

\subsubsection{Description}

\texttt{analyze(result:\allowbreak  C3ExecutionResult)} runs the 7-step pipeline and returns the full analysis JSON.

\reqid{Requirement FR-ANA-001:} The service shall expose \texttt{POST /analyze} accepting a \texttt{C3ExecutionResult} validated by Pydantic v2 (\texttt{status} literal \texttt{success|recovered|failed}; required \texttt{session\_id}, \texttt{command\_executed}).

\reqid{Requirement FR-ANA-002:} The pipeline shall execute in order: parse $\rightarrow$ enrich $\rightarrow$ predict $\rightarrow$ (fallback if no CVEs) $\rightarrow$ triage $\rightarrow$ persist graph $\rightarrow$ build response.

\reqid{Requirement FR-ANA-003:} The response shall report \texttt{session\_id}, \texttt{parse\_result} (hosts, services, raw\_cve\_ids, paths\_found), \texttt{enriched\_count}, \texttt{enriched\_\allowbreak vulnerabilities}, \texttt{predictions}, \texttt{ranked\_findings}, \texttt{false\_\allowbreak positive\_\allowbreak reduction\_\allowbreak rate}, and \texttt{latency\_ms}.

\reqid{Requirement FR-ANA-004:} The service shall expose \texttt{GET /} (\texttt{component: "C4 AVAE"}, status, version, port) and \texttt{GET /health} (component + status \texttt{healthy}).

\reqid{Requirement FR-ANA-005:} CORS shall allow origins \texttt{http://localhost:5173} and \texttt{http://127.0.0.1:5173} with credentials for the frontend.

\subsection{Feature: Scanner Output Parsing}

\subsubsection{Description}

\texttt{ParserRegistry} dispatches to the first parser whose \texttt{can\_parse} returns true; \texttt{GenericParser} is the terminal fallback.

\reqid{Requirement FR-PAR-001:} \texttt{NmapParser} shall parse nmap XML (\texttt{xml.etree.ElementTree}, CPE + NSE-script CVE hints via \texttt{CVE\_REGEX}) with plain-text fallback on XML errors, extracting IPv4 hosts, open ports, protocols, service name/product/version, and CPEs.

\reqid{Requirement FR-PAR-002:} \texttt{NiktoParser} shall parse nikto JSON and plain-text output, extracting target IP/hostname/port, Server banner, \texttt{+ /path} entries, and CVEs, deriving \texttt{service\_name} http/https.

\reqid{Requirement FR-PAR-003:} \texttt{GobusterParser} shall parse classic \texttt{(Status: NNN)} and fallback \texttt{path NNN} formats, keep statuses \{200, 201, 204, 301, 302, 303, 307, 308, 401, 403\}, and infer host/port/scheme from the \texttt{Url:} line (with dedup).

\reqid{Requirement FR-PAR-004:} \texttt{GenericParser} shall be the always-true terminal fallback, extracting only CVEs and IPs via regex.

\reqid{Requirement FR-PAR-005:} Parsing shall short-circuit on the first matching parser (\texttt{[NmapParser, NiktoParser, GobusterParser]} then \texttt{GenericParser}).

\reqid{Requirement FR-PAR-006:} All parsers shall return an \texttt{EntityCollection\{session\_\allowbreak id, tool, hosts, services, raw\_\allowbreak cve\_\allowbreak ids, paths\_\allowbreak found, parse\_\allowbreak success\}}.

\subsection{Feature: Threat Intelligence Enrichment}

\subsubsection{Description}

\texttt{LogIngestor.\allowbreak enrich(entities)} collects CVE hints, runs service/version-based CVE discovery when none exist, batches EPSS lookups, and per-CVE fetches NVD metadata + feature extraction with caching.

\reqid{Requirement FR-ENR-001:} The ingestor shall collect CVE hints from \texttt{services[].cve\_hints} and \texttt{entities.raw\_cve\_ids}, normalizing (\texttt{strip().upper()}) and building a \texttt{cve\_source\_map} (auto-association only when exactly one source service matches).

\reqid{Requirement FR-ENR-002:} When no direct CVEs exist, the ingestor shall perform service/version-based discovery: normalize banners (e.g., \texttt{Apache/2.4.49} $\rightarrow$ \texttt{Apache HTTP Server 2.\allowbreak 4.\allowbreak 49}), generate keywords (Apache/OpenSSH/nginx/others), and query \texttt{NVDClient.\allowbreak search\_\allowbreak by\_\allowbreak keyword(keyword, limit=1)}.

\reqid{Requirement FR-ENR-003:} EPSS shall be resolved in a batch (\texttt{EPSSClient.get\_batch}, capped at 50 CVEs per request).

\reqid{Requirement FR-ENR-004:} Each CVE shall resolve through \texttt{CacheManager.get} $\rightarrow$ \texttt{NVDClient.get\_cve} $\rightarrow$ feature extraction $\rightarrow$ \texttt{CacheManager.set} (cache write \texttt{INSERT OR REPLACE}).

\reqid{Requirement FR-ENR-005:} \texttt{NVDClient.\allowbreak extract\_\allowbreak features} shall yield \texttt{cvss\_v3}, \texttt{cvss\_v2}, \texttt{cwe\_ids} (prefix \texttt{CWE-}), \texttt{exploit\_db\_refs} (references containing "exploit-db"), \texttt{age\_days}, and \texttt{v3\_available}.

\reqid{Requirement FR-ENR-006:} Operational context shall be attached per vulnerability: \texttt{service\_confirmed}, \texttt{port\_open}, \texttt{service\_type} (from \texttt{\_detect\_service\_type}: web/smb/ssh/ftp/other), and \texttt{version\_\allowbreak match\_\allowbreak confidence} (0.8 with version, 0.3 without).

\reqid{Requirement FR-ENR-007:} NVD requests shall honour rate limits (sleep 0.6 s without API key, 0.1 s with key; on HTTP 429 sleep 60 s and retry) with a 15 s client timeout.

\subsection{Feature: CISA KEV and EPSS Intelligence}

\subsubsection{Description}

Startup loads the CISA KEV feed into an in-memory set; the EPSS client resolves scores for CVEs.

\reqid{Requirement FR-KEV-001:} \texttt{KEVClient.load()} shall fetch the CISA KEV JSON feed on startup (30 s timeout) and populate a set of CVE IDs (case-insensitive membership); load failure shall be logged with \texttt{is\_kev} returning \texttt{False}.

\reqid{Requirement FR-KEV-002:} \texttt{is\_kev(cve\_id)} shall return \texttt{True} if the CVE is listed as known-exploited.

\reqid{Requirement FR-EPSS-001:} \texttt{EPSSClient.\allowbreak get\_\allowbreak score(cve\_\allowbreak id)} shall return \texttt{(epss, percentile)} from the FIRST API (10 s timeout).

\reqid{Requirement FR-EPSS-002:} \texttt{EPSSClient.\allowbreak get\_\allowbreak batch(cve\_\allowbreak ids)} shall query up to 50 CVEs per request (20 s timeout) and return a score map; failures shall return zeros (fail-open).

\subsection{Feature: ML Exploitability Prediction and SHAP}

\subsubsection{Description}

\texttt{ExploitabilityPredictor} scores each enriched vulnerability with the trained XGBoost model and optional SHAP explanations.

\reqid{Requirement FR-ML-001:} The predictor shall load \texttt{XGBOOST\_MODEL\_PATH} (default \texttt{.\allowbreak /\allowbreak models/\allowbreak xgboost\_\allowbreak model.\allowbreak json}) as a binary XGBClassifier and build a \texttt{shap.TreeExplainer} (SHAP load failure non-fatal).

\reqid{Requirement FR-ML-002:} The feature vector shall use exactly \texttt{FEATURE\_\allowbreak NAMES = ["cvss\_\allowbreak v3\_\allowbreak score", "epss\_\allowbreak score", "epss\_\allowbreak percentile", "age\_\allowbreak days", "cvss\_\allowbreak epss\_\allowbreak product"]} with median imputation from \texttt{TRAINING\_MEDIANS} for missing values.

\reqid{Requirement FR-ML-003:} \texttt{predict\_batch(vulns)} shall return \texttt{list[VulnerabilityPrediction]} with \texttt{exploit\_probability} (predict\_proba) and SHAP top features per row (\texttt{\_calculate\_shap} sorted by absolute impact).

\reqid{Requirement FR-ML-004:} The composite risk score shall be computed as the exploit probability adjusted by: \texttt{kev\_mult=1.5} if KEV-listed, \texttt{epss\_boost=0.15} if \texttt{epss\_score > 0.5}, \texttt{age\_penalty=-0.05} if \texttt{age\_days < 30}, clamped to [0,1].

\reqid{Requirement FR-ML-005:} \texttt{\_risk\_tier(score)} shall map: CRITICAL $\geq$ 0.75, HIGH $\geq$ 0.55, MEDIUM $\geq$ 0.40, LOW $\geq$ 0.20, else FILTERED.

\subsection{Feature: Heuristic Service-Risk Fallback}

\subsubsection{Description}

When no CVE is identifiable, \texttt{ServiceRiskEngine.\allowbreak score(services)} scores discovered services heuristically.

\reqid{Requirement FR-SRV-001:} Base score 0.30; +0.35 for ports \{80, 443\}; +0.20 for \{21, 22\}; +0.20 for http/https, +0.10 for ssh, +0.15 for ftp service types; +0.10 if version is known; clamped [0,1].

\reqid{Requirement FR-SRV-002:} The fallback shall tier as CRITICAL $\geq$ 0.85, HIGH $\geq$ 0.70, MEDIUM $\geq$ 0.50, else LOW, and return the \textbf{top-1} service only.

\reqid{Requirement FR-SRV-003:} Fallback entries shall use \texttt{cve\_id="SERVICE\_BASED"} and \texttt{analysis\_\allowbreak mode="service\_\allowbreak based\_\allowbreak fallback"}.

\subsection{Feature: Triage, Ranking, and FP Reduction}

\subsubsection{Description}

\texttt{TriageEngine.\allowbreak triage(predictions, total\_\allowbreak input)} filters, dedupes, and ranks findings.

\reqid{Requirement FR-TRG-001:} Predictions with \texttt{composite\_\allowbreak risk\_\allowbreak score < TRIAGE\_\allowbreak THRESHOLD} (default 0.20) shall be filtered out.

\reqid{Requirement FR-TRG-002:} Findings shall be deduplicated per \texttt{cve\_id} (highest score retained) and sorted descending by composite risk score.

\reqid{Requirement FR-TRG-003:} The false-positive reduction rate shall be computed as \texttt{1.\allowbreak 0 -\allowbreak  len(ranked)/\allowbreak total\_\allowbreak input} (0 if total\_input is 0).

\subsection{Feature: Graph Persistence (Neo4j)}

\subsubsection{Description}

\texttt{Neo4jPersistence.\allowbreak save\_\allowbreak findings(session\_\allowbreak id, services, enriched, predictions)} models the analysis as a property graph.

\reqid{Requirement FR-GRP-001:} Unique constraints shall be created for \texttt{Host.ip}, \texttt{Vulnerability.cve\_id}, and \texttt{Session.session\_id}.

\reqid{Requirement FR-GRP-002:} The persistence layer shall MERGE \texttt{Session}, \texttt{Host} (with \texttt{HAS\_SERVICE}), \texttt{Service \{ip, port\}} (with name/version/protocol), and \texttt{Vulnerability \{cve\_\allowbreak id\}} (with \texttt{exploit\_probability}, \texttt{risk\_score}, \texttt{risk\_tier}).

\reqid{Requirement FR-GRP-003:} When source mapping exists (\texttt{source\_ip}, \texttt{source\_port > 0}), a \texttt{HAS\_VULNERABILITY} relationship shall link the service with \texttt{service\_name}, \texttt{service\_version}, \texttt{version\_\allowbreak match\_\allowbreak confidence}, \texttt{service\_confirmed}, and \texttt{session\_id}.

\reqid{Requirement FR-GRP-004:} A \texttt{FOUND\_VULNERABILITY} relationship shall link the session to each vulnerability.

\reqid{Requirement FR-GRP-005:} On connection failure, the layer shall set \texttt{available=False} and continue without graph persistence (fail-open), logging the condition.

\subsection{Feature: Offline Dataset and Training Scripts}

\subsubsection{Description}

\texttt{app/scripts/} provides CLI entry points for building the training dataset and training the model.

\reqid{Requirement FR-TRN-001:} \texttt{build\_\allowbreak training\_\allowbreak dataset.\allowbreak py} shall read \texttt{data/raw/CCEE.csv}, drop duplicates, label positives from the \texttt{cisa\_kev} ground truth (\texttt{true/1/yes}), and produce a 12-feature \texttt{training\_\allowbreak dataset.\allowbreak parquet}.

\reqid{Requirement FR-TRN-002:} Three scripts shall produce the 12 features: \texttt{cvss\_v3\_score}, \texttt{exploitability\_score}, \texttt{impact\_score}, \texttt{epss\_score}, \texttt{epss\_percentile}, \texttt{age\_days}, \texttt{cvss\_epss\_product}, \texttt{attack\_vector\_network}, \texttt{attack\_complexity\_low}, \texttt{privileges\_\allowbreak required\_\allowbreak none}, \texttt{user\_interaction\_none}, \texttt{scope\_changed}.

\reqid{Requirement FR-TRN-003:} \texttt{train\_model.py} shall split 70/10/20 (stratified, \texttt{random\_state=42}), grid-search XGBoost configurations (n\_estimators 200--300, max\_depth 4--6, learning\_rate 0.03--0.05, min\_child\_weight 1--5, scale\_pos\_weight 0.50--1.00), optimize the F1-sweep threshold (0.10--0.95), fit on the 5 \texttt{FEATURE\_COLS}, and persist \texttt{xgboost\_model.json} with SHAP plots.

\reqid{Requirement FR-TRN-004:} \texttt{merge\_\allowbreak global\_\allowbreak dataset.\allowbreak py} shall normalize and deduplicate \texttt{Global\_\allowbreak Vulnerability\_\allowbreak Dataset.\allowbreak xlsx} into \texttt{global\_cleaned.csv}.

\reqid{Requirement FR-TRN-005:} \texttt{test\_prediction.py} shall smoke-test the runtime predictor + triage on two sample vulnerabilities (Log4Shell high-risk and a low-risk mock).

\section{External Interface Requirements}

\subsection{Software Interfaces}

\subsubsection{REST API (HTTP/JSON)}

\begin{table}[h]
\caption{Table 4.}
\label{tab:t4}
\begin{tabularx}{\columnwidth}{l >{\raggedright\arraybackslash}X >{\raggedright\arraybackslash}X >{\raggedright\arraybackslash}X}
\toprule
\textbf{Method} & \textbf{Path} & \textbf{Auth} & \textbf{Description}\\
\midrule
POST & \texttt{/analyze} & none (gateway-level) & Full analysis pipeline\\
GET & \texttt{/} & none & Component identity/status\\
GET & \texttt{/health} & none & Health probe\\
GET & \texttt{/docs} & none & Swagger/OpenAPI UI\\
\bottomrule
\end{tabularx}
\end{table}

\subsubsection{POST /analyze --- Request (C3ExecutionResult)}

\begin{lstlisting}
{
  "session_id": "s-001",
  "status": "success",
  "command_executed": "nmap -sV -p 22,80,443 192.168.1.1",
  "stdout": "Nmap scan report for 192.168.1.1 ... 80/tcp open http Apache httpd 2.4.49 ...",
  "stderr": "",
  "exit_code": 0,
  "tool": "nmap",
  "intent_ref": "SERVICE_ENUMERATION",
  "latency_ms": 4210
}
\end{lstlisting}

\subsubsection{POST /analyze --- Response (200)}

\begin{lstlisting}
{
  "session_id": "s-001",
  "parse_result": {
    "hosts": ["192.168.1.1"],
    "services": [
      {
        "ip": "192.168.1.1",
        "port": 80,
        "protocol": "tcp",
        "service_name": "http",
        "service_version": "Apache/2.4.49",
        "banner": "Apache/2.4.49",
        "cpe": "cpe:/a:apache:http_server:2.4.49",
        "cve_hints": ["CVE-2021-41773"]
      }
    ],
    "raw_cve_ids": ["CVE-2021-41773"],
    "paths_found": []
  },
  "enriched_count": 1,
  "enriched_vulnerabilities": [
    {
      "cve_id": "CVE-2021-41773",
      "cvss_v3": 9.8,
      "cvss_v2": 7.5,
      "epss_score": 0.97,
      "epss_percentile": 0.99,
      "kev_listed": true,
      "cwe_ids": ["CWE-22"],
      "age_days": 120,
      "exploit_db_refs": 3,
      "service_confirmed": true,
      "version_match_confidence": 0.8,
      "port_open": true,
      "service_type": "web",
      "source_ip": "192.168.1.1",
      "source_port": 80,
      "source_service": "http",
      "source_version": "Apache/2.4.49"
    }
  ],
  "predictions": [
    {
      "cve_id": "CVE-2021-41773",
      "exploit_probability": 0.92,
      "composite_risk_score": 0.96,
      "risk_tier": "CRITICAL",
      "top_features": {"epss_score": 0.8, "cvss_v3_score": 0.15}
    }
  ],
  "ranked_findings": [],
  "false_positive_reduction_rate": 0.0,
  "latency_ms": 1240
}
\end{lstlisting}

\subsubsection{Error behavior}

\begin{table}[h]
\caption{Table 5.}
\label{tab:t5}
\begin{tabularx}{\columnwidth}{l >{\raggedright\arraybackslash}X}
\toprule
\textbf{Condition} & \textbf{Result}\\
\midrule
Validation failure on \texttt{C3ExecutionResult} & FastAPI 422\\
Intelligence/Neo4j/SHAP failures & fail-open; analysis proceeds with defaults\\
Unhandled exception & HTTP 500\\
\bottomrule
\end{tabularx}
\end{table}

\subsection{Hardware Interfaces}

\begin{itemize}
\item No direct hardware interfaces. Heavy ML inference (XGBoost scoring, SHAP) runs in-process on the host CPU.
\end{itemize}

\subsection{Software Interfaces (outbound HTTP)}

\begin{table}[h]
\caption{Table 6.}
\label{tab:t6}
\begin{tabularx}{\columnwidth}{l >{\raggedright\arraybackslash}X >{\raggedright\arraybackslash}X}
\toprule
\textbf{Interface} & \textbf{Endpoint} & \textbf{Purpose}\\
\midrule
NVD API 2.0 & \texttt{https:\allowbreak /\allowbreak /\allowbreak services.\allowbreak nvd.\allowbreak nist.\allowbreak gov/\allowbreak rest/\allowbreak json/\allowbreak cves/\allowbreak 2.\allowbreak 0} & CVE metadata, CVSS, CWE, keyword search\\
CISA KEV & \texttt{https:\allowbreak /\allowbreak /\allowbreak www.\allowbreak cisa.\allowbreak gov/\allowbreak sites/\allowbreak default/\allowbreak files/\allowbreak feeds/\allowbreak known\_\allowbreak exploited\_\allowbreak vulnerabilities.\allowbreak json} & known-exploited membership\\
FIRST EPSS & \texttt{https:\allowbreak /\allowbreak /\allowbreak api.\allowbreak first.\allowbreak org/\allowbreak data/\allowbreak v1/\allowbreak epss} & EPSS score + percentile (batch $\leq$ 50)\\
Neo4j & \texttt{bolt://localhost:7687} & graph persistence (optional, fail-open)\\
\bottomrule
\end{tabularx}
\end{table}

\subsection{Communication Interfaces}

\begin{itemize}
\item \textbf{Transport}: HTTP/1.1; JSON (Pydantic v2)
\item \textbf{Auth}: none in-process; gateway/network boundary expected (framework-level JWT/API-key at Kong)
\item \textbf{Client-to-AVAE timeout}: IRE \texttt{VulnerabilityClient} uses 300 s (\texttt{component4\_\allowbreak timeout\_\allowbreak seconds}); NVD 15 s; KEV 30 s; EPSS 10 s/20 s
\item \textbf{CORS}: frontend origins \texttt{localhost:5173} / \texttt{127.0.0.1:5173}
\end{itemize}

\section{Non-Functional Requirements}

\subsection{Performance}

\reqid{Requirement NFR-PERF-001:} NVD requests shall time out at 15 s and rate-limit at 0.6 s/request (0.1 s with API key); HTTP 429 shall trigger a 60 s sleep and retry.

\reqid{Requirement NFR-PERF-002:} KEV startup load shall time out at 30 s.

\reqid{Requirement NFR-PERF-003:} EPSS single lookups shall time out at 10 s; batch at 20 s, capped at 50 CVEs per call.

\reqid{Requirement NFR-PERF-004:} End-to-end analysis latency shall be measured and returned as \texttt{latency\_ms}.

\reqid{Requirement NFR-PERF-005:} The containerized deployment shall be limited to 1024 MB (reserve 512 MB) and 1.0 CPU (reserve 0.5).

\subsection{Security}

\reqid{Requirement NFR-SEC-001:} The service shall never execute scanner output as code; all parsing shall be regex/XPath/data-only (no shell evaluation).

\reqid{Requirement NFR-SEC-002:} Scanner stdout/stderr shall be treated strictly as data; prompt-injection defense is not required in-process because AVAE invokes no LLM, but upstream C1 defenses apply to raw language input.

\reqid{Requirement NFR-SEC-003:} CVE IDs shall be normalized (\texttt{strip().upper()}) and validated by regex before enrichment to prevent injection into intelligence queries.

\reqid{Requirement NFR-SEC-004:} NVD API key shall be read from environment (\texttt{NVD\_API\_KEY}); no key shall be committed to VCS.

\reqid{Requirement NFR-SEC-005:} The Neo4j password shall not use hardcoded defaults in production (\texttt{NEO4J\_PASSWORD} via environment); a hardcoded dev default is a known hardening gap.

\reqid{Requirement NFR-SEC-006:} CORS shall be restricted to the frontend origins (\texttt{localhost:5173}, \texttt{127.0.0.1:5173}).

\reqid{Requirement NFR-SEC-007:} The cache database, raw datasets, and logs (\texttt{data/raw}, \texttt{offline\_cache.db}, logs) shall be excluded from VCS.

\subsection{Reliability}

\reqid{Requirement NFR-REL-001:} All intelligence lookups (NVD, KEV, EPSS) shall fail open with defaults so analysis never aborts on external service failure.

\reqid{Requirement NFR-REL-002:} SHAP explainer unavailability shall be non-fatal (empty explanations, pipeline continues).

\reqid{Requirement NFR-REL-003:} Neo4j unavailability shall be non-fatal (\texttt{available=False}), with the analysis returned regardless.

\reqid{Requirement NFR-REL-004:} The offline cache (\texttt{OFFLINE\_CACHE\_PATH}, SQLite) shall enable repeatable enrichment and reduce external dependency; cache write failures shall be logged and ignored.

\reqid{Requirement NFR-REL-005:} Splitting batch EPSS lookups ($\leq$50) shall prevent request-size failure.

\subsection{Maintainability}

\reqid{Requirement NFR-MNT-001:} Parser addition shall follow the \texttt{BaseToolParser} ABC contract (\texttt{can\_parse}, \texttt{parse}) and register with \texttt{ParserRegistry}.

\reqid{Requirement NFR-MNT-002:} Model artifact (\texttt{xgboost\_model.json}) and SHAP plots shall be replaceable without code changes (path via \texttt{XGBOOST\_MODEL\_PATH}).

\reqid{Requirement NFR-MNT-003:} Dataset pipelining (build $\rightarrow$ merge $\rightarrow$ train $\rightarrow$ test) shall be standalone CLI scripts under \texttt{app/scripts/}.

\reqid{Requirement NFR-MNT-004:} All tunables (\texttt{PORT}, \texttt{XGBOOST\_MODEL\_PATH}, \texttt{TRIAGE\_THRESHOLD}, \texttt{NVD\_API\_KEY}, \texttt{OFFLINE\_CACHE\_PATH}, \texttt{NEO4J\_*}) shall be environment-driven via dotenv.

\reqid{Requirement NFR-MNT-005:} Database schema (cve cache) and graph constraints shall be created automatically at first use.

\subsection{Portability}

\reqid{Requirement NFR-PORT-001:} The service shall run on Linux (production) and Windows/macOS (development) with Python 3.9+.

\reqid{Requirement NFR-PORT-002:} Docker deployment shall use a 2-stage \texttt{python:3.11-slim} image.

\reqid{Requirement NFR-PORT-003:} The runtime model shall be packaged as a single XGBoost JSON artifact (no build-time compilation required).

\reqid{Requirement NFR-PORT-004:} Neo4j, NVD, KEV, and EPSS shall be out-of-process dependencies; the service shall start without them (fail-open).

\section{Data Requirements}

\subsection{Data Models (app/src/schemas/models.py)}

\begin{table}[h]
\caption{Table 7.}
\label{tab:t7}
\begin{tabularx}{\columnwidth}{l >{\raggedright\arraybackslash}X}
\toprule
\textbf{Model} & \textbf{Key fields}\\
\midrule
\texttt{C3ExecutionResult} & \texttt{session\_id}, \texttt{status} (Literal \texttt{success|recovered|failed}), \texttt{command\_executed}, \texttt{stdout}, \texttt{stderr}, \texttt{exit\_code}, \texttt{tool}, \texttt{intent\_ref}, \texttt{latency\_ms}\\
\texttt{DiscoveredService} & \texttt{ip}, \texttt{port}, \texttt{protocol}, \texttt{service\_name}, \texttt{service\_version}, \texttt{banner}, \texttt{cpe}, \texttt{cve\_hints[]}\\
\texttt{EntityCollection} & \texttt{session\_id}, \texttt{tool}, \texttt{hosts[]}, \texttt{services[]}, \texttt{raw\_cve\_ids[]}, \texttt{paths\_found[]}, \texttt{parse\_success}\\
\texttt{EnrichedVulnerability} & \texttt{cve\_id}, \texttt{cvss\_v3}, \texttt{cvss\_v2}, \texttt{epss\_score}, \texttt{epss\_percentile}, \texttt{kev\_listed}, \texttt{cwe\_ids[]}, \texttt{age\_days}, \texttt{exploit\_db\_refs}, \texttt{service\_confirmed}, \texttt{version\_\allowbreak match\_\allowbreak confidence}, \texttt{port\_open}, \texttt{service\_type}, \texttt{source\_ip}, \texttt{source\_port}, \texttt{source\_service}, \texttt{source\_version}\\
\texttt{VulnerabilityPrediction} & \texttt{cve\_id}, \texttt{exploit\_probability}, \texttt{composite\_risk\_score}, \texttt{risk\_tier}, \texttt{top\_features} (SHAP)\\
\texttt{FeedbackIntent} & \texttt{intent} (Literal \texttt{EXPLOIT\_\allowbreak TARGET|DEEP\_\allowbreak SCAN|SERVICE\_\allowbreak FINGERPRINT|CREDENTIAL\_\allowbreak TEST|PATH\_\allowbreak TRAVERSAL\_\allowbreak TEST}), \texttt{target}, \texttt{cve\_id}, \texttt{service\_port}, \texttt{risk\_score}, \texttt{confidence}, \texttt{rationale}, \texttt{session\_id}\\
\texttt{AnalysisReport} / \texttt{AnalysisResult} & defined but currently unused in the API (report + feedback intents)\\
\bottomrule
\end{tabularx}
\end{table}

\subsection{Risk Tier Taxonomy}

\begin{table}[h]
\caption{Table 8.}
\label{tab:t8}
\begin{tabularx}{\columnwidth}{l >{\raggedright\arraybackslash}X}
\toprule
\textbf{Path} & \textbf{Thresholds}\\
\midrule
ML composite score & CRITICAL $\geq$ 0.75; HIGH $\geq$ 0.55; MEDIUM $\geq$ 0.40; LOW $\geq$ 0.20; else FILTERED\\
Heuristic fallback & CRITICAL $\geq$ 0.85; HIGH $\geq$ 0.70; MEDIUM $\geq$ 0.50; else LOW\\
\bottomrule
\end{tabularx}
\end{table}

\subsection{Service Type Detection (\texttt{\_detect\_service\_type})}

\texttt{web} (http/https/www/http-alt), \texttt{smb} (smb/microsoft-ds/netbios-ssn), \texttt{ssh}, \texttt{ftp}, else \texttt{other}.

\subsection{ML Feature Contract}

\begin{itemize}
\item \textbf{Runtime} (5 features): \texttt{cvss\_v3\_score}, \texttt{epss\_score}, \texttt{epss\_percentile}, \texttt{age\_days}, \texttt{cvss\_epss\_product}; training medians used for imputation.
\item \textbf{Dataset} (12 features, offline only): adds \texttt{exploitability\_score}, \texttt{impact\_score}, \texttt{attack\_vector\_network}, \texttt{attack\_complexity\_low}, \texttt{privileges\_\allowbreak required\_\allowbreak none}, \texttt{user\_interaction\_none}, \texttt{scope\_changed}.
\item \textbf{Deployed model}: XGBClassifier, 300 trees, binary classifier. Dev metrics per README: accuracy \textasciitilde\{\}89\%, ROC-AUC \textasciitilde\{\}0.8979 (flagged for PP2 leakage/validation work).
\end{itemize}

\subsection{Persistence Stores}

\begin{table}[h]
\caption{Table 9.}
\label{tab:t9}
\begin{tabularx}{\columnwidth}{l >{\raggedright\arraybackslash}X >{\raggedright\arraybackslash}X}
\toprule
\textbf{Store} & \textbf{Type} & \textbf{Purpose}\\
\midrule
\texttt{cve\_cache} & SQLite (\texttt{OFFLINE\_CACHE\_PATH}, default \texttt{./data/offline\_cache.db}) & per-CVE enrichment cache (\texttt{INSERT OR REPLACE}); key \texttt{cve\_id}\\
Neo4j graph & bolt & Session/Host/Service/Vulnerability with \texttt{HAS\_SERVICE}, \texttt{HAS\_VULNERABILITY}, \texttt{FOUND\_VULNERABILITY}\\
Raw datasets & \texttt{data/raw/} & CCEE.csv (17.8 MB), Global VDD xlsx (18.8 MB), KEV.csv (450 KB), NVD.csv (35.8 MB)\\
Processed datasets & \texttt{data/processed/} & \texttt{global\_cleaned.csv} (53.6 MB), \texttt{training\_\allowbreak dataset.\allowbreak parquet} (3.1 MB)\\
SHAP artifacts & \texttt{models/} & \texttt{shap\_summary.png}, \texttt{shap\_\allowbreak feature\_\allowbreak importance.\allowbreak png}\\
\bottomrule
\end{tabularx}
\end{table}

\subsection{KEV Runtime State}

\begin{itemize}
\item In-memory set of known-exploited CVEs loaded at startup (log shows \textasciitilde\{\}1685 CVEs in the runtime snapshot).
\end{itemize}

\section{Verification and Validation}

\subsection{Test Suite}

\begin{table}[h]
\caption{Table 10.}
\label{tab:t10}
\begin{tabularx}{\columnwidth}{l >{\raggedright\arraybackslash}X >{\raggedright\arraybackslash}X >{\raggedright\arraybackslash}X}
\toprule
\textbf{File} & \textbf{Test classes} & \textbf{Count} & \textbf{Covers}\\
\midrule
\texttt{tests/\allowbreak unit/\allowbreak test\_\allowbreak nmap\_\allowbreak parser.\allowbreak py} & \texttt{TestCanParse}, \texttt{TestParseXml}, \texttt{TestParsePlainText}, \texttt{TestParseEdgeCases} & 20 & tool sniffing, XML/plain-text parsing, open-port filtering, versions, CVE hints, CPEs, session\_id, fallback\\
\texttt{tests/\allowbreak unit/\allowbreak test\_\allowbreak gobuster\_\allowbreak parser.\allowbreak py} & \texttt{TestCanParse}, \texttt{TestParseClassic}, \texttt{TestParseFallback}, \texttt{TestParseEdgeCases} & 18 & detection, classic/fallback parsing, status filtering, host/port/scheme inference, dedup\\
\texttt{tests/\allowbreak unit/\allowbreak test\_\allowbreak ml\_\allowbreak scoring.\allowbreak py} & \texttt{TestFeatureEngineer}, \texttt{TestServiceRiskEngine} & 14 & vector/median imputation, shapes, service-risk scoring, sorting, clamping, tiers\\
\bottomrule
\end{tabularx}
\end{table}

\reqid{Requirement NFR-TST-001:} The verified suite shall contain 52 test functions across three files.

\reqid{Requirement NFR-TST-002:} CI shall include AVAE lint/typecheck/docker-build/bandit in addition to any component-specific tests (currently only the IRE unit suite runs in \texttt{.github/workflows/ci.yml} --- flagged for hardening).

\subsection{Empirical Evaluation}

\begin{itemize}
\item Model metrics: accuracy \textasciitilde\{\}89\%, ROC-AUC \textasciitilde\{\}0.8979 (README; PP2 leakage/validation caveat documented)
\item KEV runtime load verified in server logs (\texttt{KEV loaded: ... CVEs})
\end{itemize}

\subsection{Validation Checks on Delivery}

\reqid{Requirement FR-VER-001:} Delivery shall pass the unit suite \texttt{app/tests/unit}.

\reqid{Requirement FR-VER-002:} \texttt{/health} shall report \texttt{status: healthy}.

\reqid{Requirement FR-VER-003:} A smoke \texttt{POST /analyze} with an nmap XML sample shall produce parsed hosts/services, \texttt{enriched\_count}, predictions, and \texttt{latency\_ms}.

\reqid{Requirement FR-VER-004:} A \texttt{POST /analyze} with an nmap run that produced no CVEs shall return a \texttt{service\_based\_fallback} entry (\texttt{cve\_id="SERVICE\_BASED"}).

\section{Deployment}

\subsection{Prerequisites}

\begin{itemize}
\item Reachable NVD / KEV / EPSS endpoints (optional, fail-open)
\item Optional Neo4j at \texttt{NEO4J\_URI} (default \texttt{bolt://localhost:7687})
\item Model artifact \texttt{app/\allowbreak models/\allowbreak xgboost\_\allowbreak model.\allowbreak json} present (shipped with the repo)
\end{itemize}

\subsection{Env Configuration}

\begin{table}[h]
\caption{Table 11.}
\label{tab:t11}
\begin{tabularx}{\columnwidth}{l >{\raggedright\arraybackslash}X >{\raggedright\arraybackslash}X}
\toprule
\textbf{Variable} & \textbf{Default} & \textbf{Purpose}\\
\midrule
\texttt{PORT} & \texttt{8000} & uvicorn port\\
\texttt{XGBOOST\_MODEL\_PATH} & \texttt{.\allowbreak /\allowbreak models/\allowbreak xgboost\_\allowbreak model.\allowbreak json} & model artifact\\
\texttt{TRIAGE\_THRESHOLD} & \texttt{0.20} & triage filter floor\\
\texttt{NVD\_API\_KEY} & \texttt{''} & NVD rate-limit key (optional)\\
\texttt{OFFLINE\_CACHE\_PATH} & \texttt{./data/offline\_cache.db} & SQLite enrichment cache\\
\texttt{NEO4J\_URI} & \texttt{bolt://localhost:7687} & graph endpoint\\
\texttt{NEO4J\_USER} & \texttt{neo4j} & graph auth (env)\\
\texttt{NEO4J\_PASSWORD} & dev default in code (hardening gap) & graph auth (env)\\
\bottomrule
\end{tabularx}
\end{table}

\subsection{Build and Start}

\begin{lstlisting}
docker build -t neuroshell-ai-vulnerability-analysis .
docker run -p 8000:8000 -e NVD_API_KEY=... neuroshell-ai-vulnerability-analysis
# or via dev-compose (host 8002:8000) / prod-compose with resource limits + healthcheck
\end{lstlisting}

\subsection{Verify}

\begin{lstlisting}
curl http://localhost:8000/health
curl -X POST http://localhost:8000/analyze -H "Content-Type: application/json" -d @sample.json
\end{lstlisting}

\subsection{Known Deployment Notes / Gaps}

\begin{itemize}
\item \textbf{Port ambiguity}: \texttt{main.py}/README default 8000; compose host mapping \texttt{8002:8000}; the \texttt{/} endpoint reports port 8004; IRE \texttt{component4\_url} default \texttt{http://127.0.0.1:8004}. The SRS treats 8000 as the container port and requires explicit port alignment at deployment.
\item \textbf{Startup inconsistency}: Docker \texttt{CMD ["python","main.py"]} runs \texttt{uvicorn.run("app:app")}, but the FastAPI app lives in \texttt{app.src.api.main}; documented dev invocation is \texttt{python -\allowbreak m uvicorn src.\allowbreak main:\allowbreak app -\allowbreak -\allowbreak reload}. Flagged as a deployment defect to resolve.
\item \textbf{No in-process authentication}: relies on the gateway/network boundary; framework-level JWT/API-key enforcement is expected at Kong.
\item \textbf{CI coverage}: AVAE tests are not yet executed in CI; lint/typecheck/docker-build/bandit jobs exist.
\item \textbf{Hardcoded Neo4j default password} in source --- replace via env in production.
\end{itemize}

\section{Appendices}

\subsection{Appendix A --- Pipeline Sequence (POST /analyze)}

\begin{lstlisting}
C3ExecutionResult
  1. ParserRegistry.parse(stdout, stderr, tool)          → EntityCollection
  2. LogIngestor.enrich(entities)  [NVD + KEV + EPSS]    → list[EnrichedVulnerability]
  3. ExploitabilityPredictor.predict_batch(enriched)     → list[VulnerabilityPrediction]
  4. if no CVEs: ServiceRiskEngine.score(services)       → top-1 heuristic
  5. TriageEngine.triage(predictions, total_input)       → ranked, FP-reduction
  6. Neo4jPersistence.save_findings(...)                 → graph (fail-open)
  7. response dict (session_id, parse_result, enriched, predictors, ranked, latency_ms)
\end{lstlisting}

\subsection{Appendix B --- Parser Dispatch Rules}

\begin{table}[h]
\caption{Table 12.}
\label{tab:t12}
\begin{tabularx}{\columnwidth}{l >{\raggedright\arraybackslash}X >{\raggedright\arraybackslash}X}
\toprule
\textbf{Parser} & \textbf{Detection (\texttt{can\_parse})} & \textbf{Output}\\
\midrule
\texttt{NmapParser} & tool name / XML sniff & hosts, ports, services, CPE, CVEs (NSE)\\
\texttt{NiktoParser} & tool name + \texttt{+ "..."} lines & targets, banners, paths, CVEs\\
\texttt{GobusterParser} & classic \texttt{(Status:)} / fallback \texttt{path NNN} & discovered paths + inferred host/port/scheme\\
\texttt{GenericParser} & always True (terminal) & CVEs + IPs only\\
\bottomrule
\end{tabularx}
\end{table}

\subsection{Appendix C --- Risk Threshold Table}

\begin{table}[h]
\caption{Table 13.}
\label{tab:t13}
\begin{tabularx}{\columnwidth}{l >{\raggedright\arraybackslash}X >{\raggedright\arraybackslash}X}
\toprule
\textbf{Tier} & \textbf{ML composite score} & \textbf{Heuristic score}\\
\midrule
CRITICAL & $\geq$ 0.75 & $\geq$ 0.85\\
HIGH & $\geq$ 0.55 & $\geq$ 0.70\\
MEDIUM & $\geq$ 0.40 & $\geq$ 0.50\\
LOW & $\geq$ 0.20 & < 0.50\\
FILTERED & < 0.20 & ---\\
\bottomrule
\end{tabularx}
\end{table}

\subsection{Appendix D --- Neo4j Graph Schema}

\begin{itemize}
\item Nodes: \texttt{Session\{session\_id\}}, \texttt{Host\{ip\}}, \texttt{Service\{ip, port, name, version, protocol\}}, \texttt{Vulnerability\{cve\_\allowbreak id, exploit\_\allowbreak probability, risk\_\allowbreak score, risk\_\allowbreak tier\}}
\item Relationships: \texttt{Session -\allowbreak [:\allowbreak FOUND\_\allowbreak VULNERABILITY]-\allowbreak > Vulnerability}, \texttt{Host -\allowbreak [:\allowbreak HAS\_\allowbreak SERVICE]-\allowbreak > Service}, \texttt{Service -\allowbreak [:\allowbreak HAS\_\allowbreak VULNERABILITY \{version\_\allowbreak match\_\allowbreak confidence, service\_\allowbreak confirmed, source\_\allowbreak *\}]-\allowbreak > Vulnerability}
\item Constraints: unique on \texttt{Host.ip}, \texttt{Vulnerability.cve\_id}, \texttt{Session.session\_id}
\end{itemize}

\subsection{Appendix E --- Feedback Intent Mapping (FeedbackIntentGenerator)}

\begin{table}[h]
\caption{Table 14.}
\label{tab:t14}
\begin{tabularx}{\columnwidth}{l >{\raggedright\arraybackslash}X}
\toprule
\textbf{Risk tier} & \textbf{Follow-up intent}\\
\midrule
CRITICAL & \texttt{EXPLOIT\_TARGET}\\
HIGH & \texttt{SERVICE\_FINGERPRINT}\\
MEDIUM & \texttt{DEEP\_SCAN}\\
LOW & \texttt{CREDENTIAL\_TEST}\\
FILTERED/other & \texttt{DEEP\_SCAN}\\
\bottomrule
\end{tabularx}
\end{table}

\begin{quote}\small Note: the generator is implemented but not yet wired into the \texttt{/analyze} response (target placeholder \texttt{"auto-target"}, \texttt{service\_port=0}).\end{quote}


\end{document}
